% La réaction chimique — PCT 3ème — Cours signé OnBuch+
% Programme officiel MINESEC. Compiler avec : tectonic N04-reaction-chimique.tex
\newcommand{\DOCMATIERE}{PCT}
\newcommand{\DOCNIVEAU}{3ème}
\newcommand{\DOCMODULE}{Module 1 : Chimie}
\newcommand{\DOCLECON}{N04}
\newcommand{\DOCTITRE}{La réaction chimique}
% OnBuch+ — préambule « cours signés » ENRICHI (compilé avec tectonic / XeLaTeX)
% Système visuel « Pop » d'OnBuch+ : fond crème (jamais blanc), bordures encre
% épaisses, ombres « dures » décalées, titres Archivo Black en capitales.
% Version MATHÉMATIQUES : graphes (pgfplots/tikz), boîtes pédagogiques
% (définition, propriété, méthode, attention, le savais-tu, exercice résolu,
% exercice, TP, bilan), page de garde et signature OnBuch+ ancrée en bas.
%
% Macros attendues dans main.tex AVANT \input{preamble} :
%   \DOCMATIERE  \DOCNIVEAU  \DOCMODULE  \DOCLECON (numéro)  \DOCTITRE
\documentclass[11pt]{article}

\usepackage{fontspec}
\usepackage[french]{babel}
\usepackage[a4paper,margin=2cm,top=3.4cm,bottom=2.8cm,headheight=1.4cm,footskip=1.3cm]{geometry}
\usepackage{amsmath,amssymb,amsfonts}
\usepackage{mathtools} % \xmapsto, \coloneqq... souvent utilisés par les modèles
\usepackage{cancel} % simplifications barrées (bilans de forces)
% \abs{z} (module, valeur absolue) et \norm{u} (norme), utilisés par les modèles
\providecommand{\abs}{}\let\abs\relax\DeclarePairedDelimiter{\abs}{\lvert}{\rvert}
\providecommand{\norm}{}\let\norm\relax\DeclarePairedDelimiter{\norm}{\lVert}{\rVert}
\usepackage[table]{xcolor} % [table] : fournit aussi \rowcolors (alternance de lignes)
\usepackage{pagecolor}
\usepackage{tikz}
\usetikzlibrary{arrows.meta,positioning,calc,shapes.geometric,shapes.misc,shapes.arrows,decorations.pathmorphing,decorations.pathreplacing,decorations.markings,patterns,shadows,mindmap,trees,fit,backgrounds,matrix,intersections,angles,quotes}
\usepackage{pgfplots}
\usepackage[version=4]{mhchem} % équations nucléaires \ce{^{235}_{92}U}
\usepackage[european,siunitx]{circuitikz} % schémas électriques
\pgfplotsset{compat=1.18}
\usepgfplotslibrary{fillbetween,groupplots} % aires entre courbes (calcul intégral)
\usepackage{enumitem}
\usepackage{fancyhdr}
% [most] charge toutes les bibliothèques tcolorbox (listings, minted, raster,
% documentation...) et gonfle inutilement la mémoire TeX sur un document long
% (~300 boîtes) : on ne charge que ce qui est réellement utilisé.
\usepackage[skins,breakable]{tcolorbox}
\usepackage{graphicx}
\usepackage{colortbl}
\usepackage{booktabs}
\usepackage{tabularx}
\usepackage{multirow}
\usepackage{diagbox} % case d'en-tête diagonale des tableaux à double entrée
% Colonnes extensibles centrée / gauche / droite (C, L, R), souvent utilisées par les modèles
\newcolumntype{C}{>{\centering\arraybackslash}X}
\newcolumntype{L}{>{\raggedright\arraybackslash}X}
\newcolumntype{R}{>{\raggedleft\arraybackslash}X}
\usepackage{array}
\usepackage{multicol}
\usepackage{siunitx}
% parse-numbers=false : accepte aussi \SI{2,5 \times 10^{-3}}{...}
\sisetup{locale=FR,output-decimal-marker={,},group-minimum-digits=4,parse-numbers=false}
\DeclareSIUnit\ohm{\ensuremath{\symup{\Omega}}} % U+2126 absent de la police
\DeclareSIUnit\division{div} % réglages d'oscilloscope (ms/div, V/div)
\DeclareSIUnit\tr{tr} % tours (vitesse de rotation, ex. \SI{1500}{\tr\per\minute})
\DeclareSIUnit\molar{M} % concentration molaire (ex. \SI{3}{\molar}, \SI{1}{\micro\molar})
\DeclareSIUnit\an{an} % année (le modèle confond parfois avec \year, un compteur TeX) — ex. \SI{5}{\kilo\meter\per\an}
\DeclareSIUnit\FCFA{FCFA} % franc CFA (ex. \SI{500}{\FCFA\per\kilo\gram})
\DeclareSIUnit\degreeBrix{\textdegree Brix} % teneur en sucre d'un sirop/jus (réfractométrie)
\DeclareSIUnit\month{mois} % le modèle utilise parfois \month, un compteur TeX, comme unité de durée
\DeclareSIUnit\microliter{\micro\liter} % le modèle écrit parfois \microliter en un seul token
\DeclareSIUnit\million{millions} % dénombrement (ex. \SI{15}{\million\per\mL} spermatozoïdes)
\usepackage{qrcode}
% \degree : commande fréquemment utilisée par les modèles pour un angle ou une
% température en dehors de \SI{}{} (non fournie par les paquets chargés).
\providecommand{\degree}{^{\circ}}
% \qed (fin de démonstration), utilisé par les modèles sans amsthm
\providecommand{\qed}{\hfill$\square$}
% \barre : double barre des valeurs interdites dans un tableau de variations (tkz-tab)
\providecommand{\barre}{\ensuremath{\|}}
% environnement proof (amsthm non chargé) utilisé par les modèles
\ifdefined\proof\else\newenvironment{proof}[1][Démonstration]{\par\noindent\textit{#1.}\ }{\qed\par}\fi
\usepackage{lastpage}
\usepackage{hyperref}
\hypersetup{hidelinks}

% ============================================================
% Palette « Pop » OnBuch+ — mêmes hex que la classe P (plus_ui.dart)
% ============================================================
\definecolor{poporange}{HTML}{F59321}
\definecolor{popdark}{HTML}{E07A0C}
\definecolor{popcream}{HTML}{FFF6E8}
\definecolor{popink}{HTML}{1C1714}
\definecolor{poppurple}{HTML}{7A5AE0}
\definecolor{popgreen}{HTML}{2E9E6A}
\definecolor{poppink}{HTML}{E0457B}
\definecolor{popblue}{HTML}{2F7DE1}
\definecolor{popgold}{HTML}{C98A12}
\definecolor{popmuted}{HTML}{8A7F76}
\definecolor{popline}{HTML}{EADFD2}
\definecolor{popgray}{HTML}{8A7F76}
\colorlet{popgrey}{popgray}
% Alias tolérants (le modèle nomme parfois les couleurs différemment)
\colorlet{popwhite}{white}
\colorlet{popblack}{popink}
\colorlet{popred}{poppink}
\colorlet{popyellow}{popgold}
% Teintes claires (fonds de tableaux/encadrés) — le modèle nomme parfois
% les couleurs avec un suffixe L (ex. popblueL) sans les avoir définies.
\colorlet{poporangeL}{poporange!20}
\colorlet{popdarkL}{popdark!20}
\colorlet{poppurpleL}{poppurple!20}
\colorlet{popgreenL}{popgreen!20}
\colorlet{poppinkL}{poppink!20}
\colorlet{popblueL}{popblue!20}
\colorlet{popgoldL}{popgold!20}
\colorlet{popredL}{popred!20}
\colorlet{popgrayL}{popgray!20}
% Teintes claires pour fonds de tableaux / zones
\colorlet{popblueL}{popblue!10!white}
\colorlet{poporangeL}{poporange!14!white}
\colorlet{popgreenL}{popgreen!12!white}
\colorlet{poppurpleL}{poppurple!12!white}
\colorlet{poppinkL}{poppink!12!white}
\colorlet{popgoldL}{popgold!14!white}

\pagecolor{popcream}

% ============================================================
% Typographie
% ============================================================
\defaultfontfeatures{Ligatures=TeX}
\setmainfont{PlusJakartaSans-Medium.ttf}[Path=../../../fonts/,BoldFont=PlusJakartaSans-Bold.ttf]
% Maths en sans-serif (Fira Math) avec chiffres et lettres droites de Plus
% Jakarta, pour que les formules se fondent dans le texte.
\usepackage[math-style=ISO,bold-style=ISO]{unicode-math}
\setmathfont{FiraMath-Regular.otf}[Path=../../../fonts/]
\setmathfont{PlusJakartaSans-Medium.ttf}[Path=../../../fonts/,range={up/num,up/{Latin,latin},\mathup}]
\usepackage{icomma} % virgule décimale sans espace en mode math
% Caractères mathématiques Unicode absents des polices de texte (Plus Jakarta,
% Archivo Black) : rendus par leur équivalent mathématique, en texte comme en
% formule — sinon ils s'affichent en carré vide (ex. « Arithmétique dans ℤ »).
\usepackage{newunicodechar}
\newunicodechar{ℝ}{\ensuremath{\mathbb{R}}}
\newunicodechar{ℤ}{\ensuremath{\mathbb{Z}}}
\newunicodechar{ℂ}{\ensuremath{\mathbb{C}}}
\newunicodechar{ℕ}{\ensuremath{\mathbb{N}}}
\newunicodechar{ℚ}{\ensuremath{\mathbb{Q}}}
\newunicodechar{𝔻}{\ensuremath{\mathbb{D}}}
\newunicodechar{α}{\ensuremath{\alpha}}
\newunicodechar{β}{\ensuremath{\beta}}
\newunicodechar{γ}{\ensuremath{\gamma}}
\newunicodechar{δ}{\ensuremath{\delta}}
\newunicodechar{ε}{\ensuremath{\varepsilon}}
\newunicodechar{θ}{\ensuremath{\theta}}
\newunicodechar{λ}{\ensuremath{\lambda}}
\newunicodechar{μ}{\ensuremath{\mu}}
\newunicodechar{σ}{\ensuremath{\sigma}}
\newunicodechar{τ}{\ensuremath{\tau}}
\newunicodechar{φ}{\ensuremath{\varphi}}
\newunicodechar{ψ}{\ensuremath{\psi}}
\newunicodechar{ω}{\ensuremath{\omega}}
\newunicodechar{Σ}{\ensuremath{\Sigma}}
% majuscules produites par \MakeUppercase dans les titres (α-aminés -> Α)
\newunicodechar{Α}{\ensuremath{\alpha}}
\newunicodechar{Β}{\ensuremath{\beta}}
\newunicodechar{Γ}{\ensuremath{\Gamma}}
\newunicodechar{Δ}{\ensuremath{\Delta}}
\newunicodechar{Ε}{\ensuremath{\varepsilon}}
\newunicodechar{Θ}{\ensuremath{\Theta}}
\newunicodechar{Λ}{\ensuremath{\Lambda}}
\newunicodechar{Μ}{\ensuremath{\mu}}
\newunicodechar{Τ}{\ensuremath{\tau}}
\newunicodechar{Φ}{\ensuremath{\Phi}}
\newunicodechar{Ψ}{\ensuremath{\Psi}}
\newunicodechar{Ω}{\ensuremath{\Omega}}
\newunicodechar{↦}{\ensuremath{\mapsto}}
\newunicodechar{∈}{\ensuremath{\in}}
\newunicodechar{∉}{\ensuremath{\notin}}
\newunicodechar{∧}{\ensuremath{\wedge}}
\newunicodechar{⇔}{\ensuremath{\Leftrightarrow}}
\newunicodechar{⟺}{\ensuremath{\Longleftrightarrow}}
\newunicodechar{⇒}{\ensuremath{\Rightarrow}}
\newunicodechar{⟹}{\ensuremath{\Longrightarrow}}
\newunicodechar{∘}{\ensuremath{\circ}}
\newunicodechar{∪}{\ensuremath{\cup}}
\newunicodechar{∩}{\ensuremath{\cap}}
\newunicodechar{≡}{\ensuremath{\equiv}}
\newunicodechar{⊂}{\ensuremath{\subset}}
\newunicodechar{⊥}{\ensuremath{\perp}}
\newunicodechar{∃}{\ensuremath{\exists}}
\newunicodechar{∀}{\ensuremath{\forall}}
\newunicodechar{∛}{\ensuremath{\sqrt[3]{\,}}}
\newunicodechar{∜}{\ensuremath{\sqrt[4]{\,}}}
\newunicodechar{⃗}{\ensuremath{{}^{\to}}}
\newunicodechar{ⁿ}{\ifmmode^{n}\else\textsuperscript{\ensuremath{n}}\fi}
\newunicodechar{ˣ}{\ifmmode^{x}\else\textsuperscript{\ensuremath{x}}\fi}
\newunicodechar{ᵇ}{\ifmmode^{b}\else\textsuperscript{\ensuremath{b}}\fi}
\newunicodechar{ʳ}{\ifmmode^{r}\else\textsuperscript{\ensuremath{r}}\fi}
\newunicodechar{ᵘ}{\ifmmode^{u}\else\textsuperscript{\ensuremath{u}}\fi}
\newunicodechar{ʸ}{\ifmmode^{y}\else\textsuperscript{\ensuremath{y}}\fi}
\newunicodechar{ᵃ}{\ifmmode^{a}\else\textsuperscript{\ensuremath{a}}\fi}
\newunicodechar{ᵉ}{\ifmmode^{e}\else\textsuperscript{\ensuremath{e}}\fi}
\newunicodechar{ᵏ}{\ifmmode^{k}\else\textsuperscript{\ensuremath{k}}\fi}
\newunicodechar{ᵖ}{\ifmmode^{p}\else\textsuperscript{\ensuremath{p}}\fi}
\newunicodechar{ᴺ}{\ifmmode^{N}\else\textsuperscript{\ensuremath{N}}\fi}
\newunicodechar{ᶜ}{\ifmmode^{c}\else\textsuperscript{\ensuremath{c}}\fi}
\newunicodechar{ʰ}{\ifmmode^{h}\else\textsuperscript{\ensuremath{h}}\fi}
\newunicodechar{ᵠ}{\ifmmode^{q}\else\textsuperscript{\ensuremath{q}}\fi}
\newunicodechar{ₙ}{\ifmmode_{n}\else\textsubscript{\ensuremath{n}}\fi}
\newunicodechar{ₐ}{\ifmmode_{a}\else\textsubscript{\ensuremath{a}}\fi}
\newunicodechar{ᵢ}{\ifmmode_{i}\else\textsubscript{\ensuremath{i}}\fi}
\newunicodechar{ᵣ}{\ifmmode_{r}\else\textsubscript{\ensuremath{r}}\fi}
\newunicodechar{ₖ}{\ifmmode_{k}\else\textsubscript{\ensuremath{k}}\fi}
\newunicodechar{ᵦ}{\ifmmode_{\beta}\else\textsubscript{\ensuremath{\beta}}\fi}
\newunicodechar{ₓ}{\ifmmode_{x}\else\textsubscript{\ensuremath{x}}\fi}
\newfontfamily\popdisplay{ArchivoBlack-Regular.ttf}[Path=../../../fonts/]
\newfontfamily\poplogo{SpaceGrotesk-Bold.ttf}[Path=../../../fonts/]
\newfontfamily\popsemi{PlusJakartaSans-SemiBold.ttf}[Path=../../../fonts/]
\newfontfamily\popxbold{PlusJakartaSans-ExtraBold.ttf}[Path=../../../fonts/]
\newfontfamily\popxboldls{PlusJakartaSans-ExtraBold.ttf}[Path=../../../fonts/,LetterSpace=3.0]

\setlist[enumerate]{leftmargin=1.7em,itemsep=5pt,topsep=4pt}
\setlist[itemize]{leftmargin=1.7em,itemsep=4pt,topsep=4pt,label={\color{poporange}\textbullet}}
\setlength{\parindent}{0pt}
\setlength{\parskip}{5pt}
\renewcommand{\arraystretch}{1.25}

% pgfplots : style Pop par défaut
\pgfplotsset{
  every axis/.append style={axis line style={popink,line width=1pt},tick style={popink},
    grid style={popline,line width=0.5pt},label style={font=\small},tick label style={font=\footnotesize}},
  popcurve/.style={poporange,line width=1.8pt,no markers,smooth},
}
% Même style disponible dans un \draw TikZ simple (hors pgfplots)
\tikzset{popcurve/.style={poporange,line width=1.8pt}}
\tikzset{popgrid/.style={popline,very thin}}
\colorlet{popcurve}{poporange} % le nom du style est parfois utilisé comme couleur

% ============================================================
% Pill / squiggle
% ============================================================
\newcommand{\poppill}[3]{%
  \tikz[baseline=-0.55ex]\node[draw=popink,line width=1.2pt,rounded corners=2.6pt,%
    fill=#2,inner xsep=6.5pt,inner ysep=3pt]{%
    \popxboldls\fontsize{7}{7}\selectfont\color{#3}\MakeUppercase{#1}};%
}
\newcommand{\popsquiggle}[1]{%
  \tikz[baseline=-0.4ex]{
    \draw[#1,line width=2.6pt,line cap=round]
      plot [smooth] coordinates {(0,0.09) (0.34,0.01) (0.68,0.21) (1.02,0.01) (1.36,0.10)};
  }%
}
% Mot-clé surligné (marqueur orange sous le texte)
\newcommand{\cle}[1]{{\popxbold\color{popdark}#1}}

% ============================================================
% En-tête / pied
% ============================================================
\pagestyle{fancy}
\fancyhf{}
\renewcommand{\headrulewidth}{0pt}
\renewcommand{\footrulewidth}{0pt}
\lhead{\raisebox{-0.15ex}{\poplogo\fontsize{13}{13}\selectfont\color{popink}OnBuch\color{poporange}+}}
\rhead{\poppill{\DOCMATIERE\ \textbullet\ \DOCNIVEAU\ \textbullet\ Leçon \DOCLECON}{white}{popink}}
\lfoot{\popxbold\fontsize{7.6}{7.6}\selectfont\color{popmuted}\MakeUppercase{Cours signé OnBuch+}\ \textbullet\ {\popsemi\DOCTITRE}}
\rfoot{\tikz[baseline=-0.6ex]\node[rounded corners=8.5pt,fill=popink,minimum height=17pt,minimum width=17pt,inner xsep=5pt,inner sep=0]{\popxbold\fontsize{8}{8}\selectfont\color{popcream}\thepage\,/\,\pageref*{LastPage}};}

% ============================================================
% Titres
% ============================================================
\newcommand{\chaptitre}[2]{%
  \begin{tcolorbox}[enhanced,colback=poporange,colframe=popink,boxrule=2.6pt,arc=11pt,
      drop shadow=popink,left=15pt,right=15pt,top=12pt,bottom=11pt,before skip=3pt,after skip=18pt]
    \poppill{#2}{popink}{popcream}\par\vspace{9pt}
    {\raggedright\hyphenpenalty=10000\exhyphenpenalty=10000\popdisplay\fontsize{22}{24}\selectfont\color{popcream}\MakeUppercase{#1}\par}
  \end{tcolorbox}
}
% Section numérotée : pastille orange avec le numéro + titre Archivo
\newcounter{popsec}
\newcommand{\coursec}[1]{%
  \stepcounter{popsec}\setcounter{popsub}{0}%
  \par\vspace{16pt}\noindent
  \begin{minipage}{\linewidth}
  \tikz[baseline=-0.7ex]\node[draw=popink,line width=1.6pt,fill=poporange,rounded corners=4pt,minimum size=22pt,inner sep=0]{\popdisplay\fontsize{12}{12}\selectfont\color{popcream}\thepopsec};%
  \hspace{8pt}{\popdisplay\fontsize{15}{17}\selectfont\color{popink}\MakeUppercase{#1}}\par
  \vspace{4pt}\noindent{\color{poporange}\rule{36pt}{3.6pt}}
  \end{minipage}\par\nopagebreak\vspace{8pt}
}
\newcounter{popsub}
\newcommand{\courssub}[1]{%
  \stepcounter{popsub}%
  \par\vspace{9pt}\noindent{\popxbold\fontsize{11.5}{13}\selectfont\color{popink}{\color{poporange}\thepopsec.\thepopsub}\hspace{6pt}#1}\par\nopagebreak\vspace{4pt}
}

% ============================================================
% Boîtes pédagogiques — même recette : fond blanc, bordure épaisse colorée,
% ombre dure encre, badge plein. Titre optionnel [#1].
% ============================================================
\tcbset{popbase/.style={breakable,enhanced,colback=white,boxrule=2.3pt,arc=9pt,
  shadow={3pt}{-3pt}{0pt}{popink},left=14pt,right=14pt,top=11pt,bottom=11pt,before skip=11pt,after skip=15pt,
  fontupper=\popsemi\fontsize{10.6}{14.5}\selectfont\color{popink},
  fontlower=\popsemi\fontsize{10.6}{14.5}\selectfont\color{popink}}}
% Coche dessinée (le glyphe ✓ n'existe pas dans les polices du document)
\newcommand{\popcheck}[1]{\tikz[baseline=-0.5ex]\draw[#1,line width=1.6pt,line cap=round,line join=round](-0.13,0.0)--(-0.04,-0.09)--(0.13,0.1);}
\providecommand{\coche}{\popcheck{popgreen}}
\providecommand{\resultat}[1]{\par\smallskip\noindent\fcolorbox{popgreen}{popgreenL}{\parbox{\dimexpr\linewidth-2\fboxsep-2\fboxrule\relax}{\textbf{Résultat.} #1}}\par\smallskip}
\newcommand{\popboxhead}[3]{\poppill{#1}{#2}{white}\ifx\relax#3\relax\else\hspace{6pt}{\popxbold\fontsize{10.6}{12}\selectfont\color{popink}#3}\fi\par\vspace{7pt}}

\newtcolorbox{aretenir}[1][]{popbase,colframe=popgreen,before upper={\popboxhead{À retenir}{popgreen}{#1}}}
\newtcolorbox{exemplebox}[1][]{popbase,colframe=poporange,before upper={\popboxhead{Exemple}{poporange}{#1}}}
\newtcolorbox{definition}[1][]{popbase,colframe=popblue,before upper={\popboxhead{Définition}{popblue}{#1}}}
\newtcolorbox{propriete}[1][]{popbase,colframe=poppurple,before upper={\popboxhead{Propriété}{poppurple}{#1}}}
\newtcolorbox{methode}[1][]{popbase,colframe=popgold,before upper={\popboxhead{Méthode}{popgold}{#1}}}
\newtcolorbox{attention}[1][]{popbase,colframe=poppink,before upper={\popboxhead{Attention}{poppink}{#1}}}
\newtcolorbox{experience}[1][]{popbase,colframe=popblue,colback=popblueL,before upper={\popboxhead{Expérience}{popblue}{#1}}}
% « Le savais-tu ? » : carte violette pleine (contexte, culture, Cameroun)
\newtcolorbox{savaistu}[1][]{popbase,colback=poppurple,colframe=popink,
  fontupper=\popsemi\fontsize{10.4}{14.2}\selectfont\color{white},
  before upper={\poppill{Le savais-tu\,?}{white}{poppurple}\ifx\relax#1\relax\else\hspace{6pt}{\popxbold\color{white}#1}\fi\par\vspace{7pt}}}
% Exercice résolu : énoncé en haut, \tcblower puis la solution sur fond crème
\newcounter{popexo}\newcounter{popexr}
\newtcolorbox{exoresolu}[1][]{popbase,colframe=poporange,colbacklower=popcream,
  segmentation style={popink,line width=1.2pt,dashed},
  before upper={\stepcounter{popexr}\popboxhead{Exercice résolu \thepopexr}{poporange}{#1}},
  before lower={\poppill{Solution}{popink}{popcream}\par\vspace{6pt}}}
% Exercice (non résolu en place) — la correction est dans la partie Corrigés
\newtcolorbox{exercice}[1][]{popbase,colframe=popink,
  before upper={\stepcounter{popexo}\popboxhead{Exercice \thepopexo}{popink}{#1}}}
% Corrigé
\newtcolorbox{corrige}[1][]{popbase,colframe=popgreen,colback=popgreenL,
  before upper={\popboxhead{Corrigé}{popgreen}{#1}}}
% Objectifs / prérequis / bilan
\newtcolorbox{objectifs}{popbase,colframe=popink,colback=poporangeL,
  before upper={\popboxhead{Objectifs de la leçon}{popink}{}}}
\newtcolorbox{prerequis}{popbase,colframe=popink,colback=popblueL,
  before upper={\popboxhead{Prérequis}{popblue}{}}}
\newtcolorbox{bilan}{popbase,colframe=popink,colback=white,boxrule=2.8pt,
  before upper={\popboxhead{Fiche bilan}{poporange}{L'essentiel à connaître pour l'examen}}}
% Figure encadrée avec légende
\newtcolorbox{popfigure}[1][]{enhanced,breakable=false,colback=white,colframe=popline,boxrule=1.4pt,arc=8pt,
  left=10pt,right=10pt,top=10pt,bottom=8pt,before skip=10pt,after skip=12pt,halign=center,
  lower separated=false,halign lower=center,
  fontlower=\popsemi\fontsize{9}{11.5}\selectfont\color{popmuted},
  #1}
\newcommand{\legende}[1]{\par\vspace{6pt}{\popsemi\fontsize{9}{11.5}\selectfont\color{popmuted}#1}}

% ============================================================
% Signature OnBuch+
% ============================================================
\newcommand{\poplogoinline}[1]{{\poplogo\fontsize{#1}{#1}\selectfont\color{popink}OnBuch\color{poporange}+}}
\newcommand{\signatureonbuch}{%
  \begin{tcolorbox}[enhanced,colback=popink,colframe=popink,boxrule=2.2pt,arc=10pt,
      drop shadow=poporange,left=14pt,right=14pt,top=10pt,bottom=10pt,before skip=0pt,after skip=0pt]
    \tikz[baseline=-0.7ex]\node[circle,fill=popgreen,draw=popcream,line width=1.3pt,minimum size=22pt,inner sep=0]{\popcheck{white}};%
    \hspace{9pt}\begin{minipage}[c]{0.62\linewidth}
      {\popxbold\fontsize{12}{13}\selectfont\color{popcream}Signé\ \textbullet\ {\poplogo OnBuch\color{poporange}+}}\par\vspace{2pt}
      {\raggedright\popsemi\fontsize{8}{10.5}\selectfont\color{popline}\MakeUppercase{Équipe pédagogique OnBuch+}\\\MakeUppercase{Conforme au programme officiel MINESEC}\par}
    \end{minipage}\hfill
    {\poplogo\fontsize{22}{22}\selectfont\color{popcream}OnBuch\color{poporange}+}
  \end{tcolorbox}%
}

% ============================================================
% Page de garde
% #1 titre · #2 matière · #3 niveau · #4 module · #5 n° leçon · #6 durée ·
% #7 description / objectifs (texte ou itemize)
% ============================================================
\newcommand{\pagegarde}[7]{%
  \thispagestyle{empty}
  \newgeometry{a4paper,margin=1.7cm,top=1.6cm,bottom=1.4cm}%
  \begin{tikzpicture}[remember picture,overlay]
    \fill[poporange!18] ([xshift=-3.2cm,yshift=-3.6cm]current page.north east) circle (4.6cm);
    \fill[poppurple!14] ([xshift=2.4cm,yshift=7.6cm]current page.south west) circle (3.8cm);
  \end{tikzpicture}%
  {\poplogo\fontsize{30}{30}\selectfont\color{popink}OnBuch\color{poporange}+}\hfill
  \raisebox{0.6ex}{\poppill{Cours signé \textbullet\ Premium}{popink}{popcream}}\par
  \vspace{1pt}\popsquiggle{poppurple}\par
  \vspace{8pt}
  \begin{tcolorbox}[enhanced,colback=poporange,colframe=popink,boxrule=2.8pt,arc=13pt,
      drop shadow=popink,left=18pt,right=18pt,top=12pt,bottom=13pt,after skip=10pt]
    \poppill{#2 \textbullet\ #3}{popink}{popcream}\hspace{5pt}\poppill{#4}{white}{popink}\par\vspace{8pt}
    {\popdisplay\fontsize{14}{14}\selectfont\color{popink}LEÇON #5}\par\vspace{4pt}
    {\raggedright\hyphenpenalty=10000\exhyphenpenalty=10000\popdisplay\fontsize{25}{27}\selectfont\color{popcream}\MakeUppercase{#1}\par}
    \vspace{8pt}
    {\popxbold\fontsize{9.5}{9.5}\selectfont\color{popink}\MakeUppercase{Durée conseillée : #6}}
  \end{tcolorbox}
  \begin{tcolorbox}[enhanced,colback=white,colframe=popink,boxrule=2.2pt,arc=10pt,
      drop shadow=popink,left=16pt,right=16pt,top=10pt,bottom=10pt,after skip=8pt]
    \poppill{Dans cette leçon}{poppurple}{white}\par\vspace{5pt}
    \popsemi\fontsize{9.3}{12.6}\selectfont\color{popink}#7
  \end{tcolorbox}
  \noindent\begin{minipage}[t]{0.487\linewidth}
    \begin{tcolorbox}[enhanced,colback=popgreen,colframe=popink,boxrule=2.2pt,arc=10pt,
        drop shadow=popink,left=11pt,right=11pt,top=10pt,bottom=10pt,height=3.1cm,valign=center]
      \tcbox[colback=white,colframe=popink,boxrule=1.3pt,arc=5pt,boxsep=0pt,nobeforeafter,
        left=3pt,right=3pt,top=3pt,bottom=3pt,valign=center]{\qrcode[height=1.9cm]{https://play.google.com/store/apps/details?id=com.luvvix.onbuch&hl=fr}}%
      \hspace{8pt}\begin{minipage}[c]{0.5\linewidth}
        {\popdisplay\fontsize{11}{12.5}\selectfont\color{white}\MakeUppercase{Télécharge OnBuch}}\par\vspace{4pt}
        {\raggedright\popsemi\fontsize{8.4}{11}\selectfont\color{white}Gratuit \textbullet\ Google Play\\Tuteur IA, annales, concours, résultats\par}
      \end{minipage}
    \end{tcolorbox}
  \end{minipage}\hfill
  \begin{minipage}[t]{0.487\linewidth}
    \begin{tcolorbox}[enhanced,colback=poppurple,colframe=popink,boxrule=2.2pt,arc=10pt,
        drop shadow=popink,left=11pt,right=11pt,top=10pt,bottom=10pt,height=3.1cm,valign=center]
      \tikz[baseline=-0.7ex]\node[circle,fill=white,draw=popink,line width=1.3pt,minimum size=1.6cm,inner sep=0]{\popdisplay\fontsize{24}{24}\selectfont\color{poppurple}+};%
      \hspace{8pt}\begin{minipage}[c]{0.6\linewidth}
        {\popdisplay\fontsize{11}{12.5}\selectfont\color{white}\MakeUppercase{Passe à OnBuch+}}\par\vspace{4pt}
        {\raggedright\popsemi\fontsize{8.4}{11}\selectfont\color{white}Tous les cours signés, Tuteur IA illimité, fiches PDF — onglet OnBuch+ de l'app\par}
      \end{minipage}
    \end{tcolorbox}
  \end{minipage}
  \vspace{8pt}
  \popcontenu
  \vfill
  \signatureonbuch
  \restoregeometry
  \newpage
}

% Rangée « Ce cours contient » (page de garde)
\newcommand{\poptile}[3]{% #1 couleur #2 gros texte #3 légende
  \begin{tcolorbox}[enhanced,colback=#1,colframe=popink,boxrule=2pt,arc=9pt,shadow={2.5pt}{-2.5pt}{0pt}{popink},
      left=6pt,right=6pt,top=9pt,bottom=9pt,halign=center,valign=center,height=1.75cm,width=\linewidth,nobeforeafter]
    {\popdisplay\fontsize{12.5}{13}\selectfont\color{white}\MakeUppercase{#2}}\par\vspace{3pt}
    {\centering\popsemi\fontsize{7.8}{9.5}\selectfont\color{white}#3\par}
  \end{tcolorbox}}
\newcommand{\popcontenu}{%
  \par\noindent{\popxboldls\fontsize{7.5}{7.5}\selectfont\color{popmuted}CE COURS SIGNÉ CONTIENT}\par\vspace{7pt}
  \noindent\begin{minipage}[t]{0.235\linewidth}\poptile{popblue}{Cours}{complet, illustré, pas à pas}\end{minipage}\hfill
  \begin{minipage}[t]{0.235\linewidth}\poptile{poporange}{Méthodes}{et exercices résolus}\end{minipage}\hfill
  \begin{minipage}[t]{0.235\linewidth}\poptile{poppink}{Exercices}{type examen + corrigés}\end{minipage}\hfill
  \begin{minipage}[t]{0.235\linewidth}\poptile{popgreen}{Fiche bilan}{l'essentiel à retenir}\end{minipage}\par}

% Sommaire
\newtcolorbox{sommaire}{enhanced,colback=white,colframe=popink,boxrule=2.2pt,arc=10pt,
  drop shadow=popink,left=18pt,right=18pt,top=13pt,bottom=13pt,before skip=6pt,after skip=20pt,
  before upper={\poppill{Sommaire}{poppurple}{white}\par\vspace{9pt}\popsemi\fontsize{10.6}{14.5}\selectfont\color{popink}}}

% Signature de fin de cours — poussée tout en bas de la dernière page
\newcommand{\findecours}{%
  \par\vfill\nopagebreak
  \begin{center}\popsquiggle{poporange}\end{center}\vspace{4pt}
  \signatureonbuch
}
\AtBeginDocument{\def\llbracket{\mathopen{[\mkern-3.5mu[}}\def\rrbracket{\mathclose{]\mkern-3.5mu]}}} % intervalles d'entiers (glyphe ⟦ absent de la police)
% Glyphes absents de Fira Math : redéfinis à partir de glyphes disponibles
\AtBeginDocument{%
  \def\setminus{\mathbin{\backslash}}\let\smallsetminus\setminus
  \def\vdots{\mathord{\vbox{\baselineskip3.2pt\lineskiplimit0pt\kern4pt\hbox{.}\hbox{.}\hbox{.}}}}%
  \def\ddots{\mathinner{\mkern1mu\raise6.5pt\hbox{.}\mkern2mu\raise3.6pt\hbox{.}\mkern2mu\raise0.7pt\hbox{.}\mkern1mu}}%
  \def\triangle{\mathord{\scalebox{0.9}{$\symup{\Delta}$}}}%
  \def\nabla{\mathord{\raisebox{\depth}{\scalebox{1}[-1]{$\symup{\Delta}$}}}}%
  \def\checkmark{\popcheck{popdark}}%
  \def\star{\mathbin{\ast}}\def\bigstar{\mathord{\ast}}%
  \def\diamond{\mathbin{\scalebox{0.6}{$\Diamond$}}}%
  \def\bigcup{\mathop{\mathchoice{\raisebox{-0.3em}{\scalebox{1.7}{$\cup$}}}{\scalebox{1.25}{$\cup$}}{\cup}{\cup}}\displaylimits}%
  \def\bigcap{\mathop{\mathchoice{\raisebox{-0.3em}{\scalebox{1.7}{$\cap$}}}{\scalebox{1.25}{$\cap$}}{\cap}{\cap}}\displaylimits}%
  \def\lesseqqgtr{\mathrel{\lessgtr}}\def\bigoplus{\mathop{\oplus}\displaylimits}%
}
\providecommand{\ssi}{si et seulement si} % abréviation fréquente
\AtBeginDocument{\providecommand{\wideparen}{\overparen}} % arc de cercle

%% FIGFIX-BEGIN v5 — correctifs globaux des figures (docs/onbuch-generation/tools/figfix)
\usepackage{adjustbox}\usepackage{etoolbox}\tcbuselibrary{raster}
% toute figure TikZ/pgfplots plus large que la ligne est réduite à la largeur disponible
\BeforeBeginEnvironment{tikzpicture}{\begin{adjustbox}{max width=\linewidth}}
\AfterEndEnvironment{tikzpicture}{\end{adjustbox}}
% légendes pgfplots lisibles par-dessus les courbes
\pgfplotsset{every axis/.append style={legend style={fill=white,fill opacity=0.92,text opacity=1,draw=black!15,inner sep=3pt}}}
% étiquettes d'arêtes (syntaxe edge["..."]) avec fond blanc
\tikzset{every edge quotes/.append style={fill=white,inner sep=1.2pt}}
%% FIGFIX-END
\begin{document}

\pagegarde{La réaction chimique}{PCT}{3ème}{Module 1 : Chimie}{N04}{2 séances (1 h chacune)}{Cette leçon introduit la notion de réaction chimique, son écriture sous forme d'équation-bilan et les règles d'équilibrage traduisant la conservation de la matière. Elle étudie ensuite les combustions complètes et incomplètes, le danger du monoxyde de carbone et les gestes de prévention dans la vie quotidienne.
\vspace{4pt}
\begin{multicols}{2}\small
\begin{itemize}
\item À la fin de cette leçon, je sais distinguer une transformation chimique d'une transformation physique à partir d'indices expérimentaux.
\item À la fin de cette leçon, je sais identifier les réactifs et les produits d'une réaction chimique.
\item À la fin de cette leçon, je sais écrire l'équation-bilan d'une réaction chimique simple.
\item À la fin de cette leçon, je sais équilibrer une équation-bilan en appliquant la conservation des atomes (loi de Lavoisier).
\item À la fin de cette leçon, je sais réaliser et interpréter des expériences simples de réaction chimique (effervescence, combustion, précipitation).
\item À la fin de cette leçon, je sais distinguer une combustion complète d'une combustion incomplète et écrire leurs équations-bilans.
\end{itemize}\end{multicols}}

\begin{sommaire}
\begin{multicols}{2}
\begin{enumerate}
\item Transformation chimique : mise en évidence expérimentale
\item Conservation de la masse et loi de Lavoisier
\item Équation-bilan et équilibrage
\item Combustions complète et incomplète
\item Le monoxyde de carbone : un danger invisible
\item Sécurité et gestes responsables autour des combustions
\item Activité d'intégration
\item Méthodes et exercices résolus
\item Exercices
\item Corrigés des exercices
\item Fiche bilan
\end{enumerate}
\end{multicols}
\end{sommaire}

\begin{objectifs}
\begin{itemize}
\item À la fin de cette leçon, je sais distinguer une transformation chimique d'une transformation physique à partir d'indices expérimentaux.
\item À la fin de cette leçon, je sais identifier les réactifs et les produits d'une réaction chimique.
\item À la fin de cette leçon, je sais écrire l'équation-bilan d'une réaction chimique simple.
\item À la fin de cette leçon, je sais équilibrer une équation-bilan en appliquant la conservation des atomes (loi de Lavoisier).
\item À la fin de cette leçon, je sais réaliser et interpréter des expériences simples de réaction chimique (effervescence, combustion, précipitation).
\item À la fin de cette leçon, je sais distinguer une combustion complète d'une combustion incomplète et écrire leurs équations-bilans.
\item À la fin de cette leçon, je sais expliquer la toxicité du monoxyde de carbone et citer les mesures de prévention des intoxications domestiques.
\item À la fin de cette leçon, je sais appliquer les consignes de sécurité lors de manipulations mettant en jeu des combustions.
\end{itemize}
\end{objectifs}

\begin{prerequis}
\begin{itemize}
\item Notion d'atome, de molécule et de formule chimique (H₂O, CO₂, O₂) vues en 4ème et en début de 3ème
\item Corps purs et mélanges ; changements d'état (transformations physiques)
\item Tests d'identification des gaz : dioxygène (bûchette incandescente), dioxyde de carbone (eau de chaux), dihydrogène (détonation)
\item Notion de combustible et de comburant (triangle du feu, vu en 4ème)
\end{itemize}
\end{prerequis}

\newpage

\coursec{Leçon N04 : La réaction chimique}

\textbf{Situation d'accroche.} À Yaoundé, pendant la saison des pluies, une famille allume un brasero à charbon dans une pièce fermée pour se réchauffer. Le lendemain matin, deux personnes se réveillent avec de violents maux de tête et des nausées. Quelques jours plus tôt, au moment de la cuisson du poisson braisé, le même charbon rougeoyant produisait d'abord de petites flammes bleues, puis de grandes flammes jaunes fumeuses. Que s'est-il passé dans l'air de la pièce ? Pourquoi le charbon semble-t-il « disparaître » en brûlant alors que, en chimie, rien ne se perd ? Comment décrire et écrire la transformation qui a lieu, et surtout comment éviter ce drame domestique malheureusement fréquent au Cameroun ?

\textbf{Problématique.} Comment reconnaître qu'une transformation chimique a eu lieu ? Quelles espèces disparaissent, quelles espèces apparaissent, et que deviennent les atomes pendant la transformation ? Comment écrire la réaction sous forme d'équation-bilan ? Pourquoi certaines combustions sont-elles dangereuses ?

Pour répondre à ces questions, nous allons d'abord distinguer transformation physique et transformation chimique, puis mener des expériences d'investigation, établir la loi de conservation de la masse, apprendre à écrire et équilibrer une équation-bilan, et enfin analyser les combustions et leurs dangers.

\courssub{Rappel : la transformation physique}

\begin{definition}[Transformation physique]
Une \cle{transformation physique} est une transformation au cours de laquelle l'espèce chimique \textbf{ne change pas de nature} : seuls son état (solide, liquide, gaz), sa forme ou sa répartition dans un mélange sont modifiés.
\end{definition}

Tu connais déjà de nombreux exemples :
\begin{itemize}
\item la \textbf{fonte de la glace} : l'eau solide devient liquide, mais c'est toujours de l'eau, \ce{H2O} ;
\item la \textbf{dissolution du sel} dans l'eau de la sauce : le sel se disperse, mais il reste du chlorure de sodium, \ce{NaCl} (on peut le récupérer par évaporation) ;
\item la \textbf{vaporisation de l'eau} dans la marmite : la vapeur qui s'échappe est toujours de l'eau.
\end{itemize}

Dans tous ces cas, aucune espèce nouvelle n'apparaît : les molécules de départ sont intactes. C'est précisément ce qui distingue une transformation physique d'une transformation chimique.

\courssub{Expérience 1 : action de l'acide chlorhydrique sur la craie}

\begin{experience}[L'acide chlorhydrique attaque la craie]
\textbf{Matériel :} deux tubes à essai, un bouchon percé muni d'un tube à dégagement, un support, de la craie (carbonate de calcium \ce{CaCO3}), de l'acide chlorhydrique dilué, de l'eau de chaux.

\textbf{Protocole :}
\begin{enumerate}
\item Placer un morceau de craie au fond du premier tube à essai.
\item Verser délicatement de l'acide chlorhydrique dilué sur la craie, puis boucher aussitôt le tube avec le bouchon muni du tube à dégagement.
\item Plonger l'extrémité du tube à dégagement dans le second tube contenant de l'eau de chaux limpide.
\item Observer.
\end{enumerate}

\textbf{Observations :}
\begin{itemize}
\item Il se produit une \textbf{effervescence} : des bulles de gaz se forment à la surface de la craie, qui diminue de volume.
\item Le gaz recueilli barbote dans l'eau de chaux, qui \textbf{se trouble} (un voile blanchâtre apparaît).
\end{itemize}

\textbf{Interprétation :} le trouble de l'eau de chaux est le test caractéristique du \cle{dioxyde de carbone} \ce{CO2}. Or le dioxyde de carbone n'existait ni dans la craie seule, ni dans l'acide seul : c'est une \textbf{espèce nouvelle}, formée pendant la transformation. La craie et l'acide ont été partiellement \textbf{consommés}.

\textbf{Équation-bilan de la réaction :}
\[ \ce{CaCO3(s) + 2 HCl(aq) -> CaCl2(aq) + CO2(g) + H2O(l)} \]
\end{experience}

\begin{popfigure}
\begin{center}
\begin{tikzpicture}[scale=1.0, every node/.style={font=\small}]
% Tube 1 : craie + acide
\draw[thick, popdark] (0,0) -- (0,2.6);
\draw[thick, popdark] (1.2,0) -- (1.2,2.6);
\draw[thick, popdark] (0,0) arc (180:360:0.6);
% Liquide acide
\fill[popblueL] (0.03,0.25) -- (0.03,1.3) -- (1.17,1.3) -- (1.17,0.25) arc (180:360:0.57) -- cycle;
% Craie
\fill[popgold] (0.35,0.12) rectangle (0.85,0.5);
\draw[popdark] (0.35,0.12) rectangle (0.85,0.5);
% Bulles
\foreach \p in {(0.5,0.9),(0.7,1.1),(0.6,0.7),(0.9,0.95),(0.45,1.15)} \draw[popblue, fill=white] \p circle (0.06);
% Bouchon
\fill[popmuted!60] (0.1,2.6) rectangle (1.1,2.95);
% Tube à dégagement
\draw[thick, popdark] (0.6,2.95) -- (0.6,3.5) -- (4.6,3.5) -- (4.6,1.6);
% Tube 2 : eau de chaux
\draw[thick, popdark] (4.0,0) -- (4.0,2.6);
\draw[thick, popdark] (5.2,0) -- (5.2,2.6);
\draw[thick, popdark] (4.0,0) arc (180:360:0.6);
\fill[popmuted!35] (4.03,0.25) -- (4.03,1.5) -- (5.17,1.5) -- (5.17,0.25) arc (180:360:0.57) -- cycle;
% Bulles dans eau de chaux
\foreach \p in {(4.6,0.5),(4.6,0.8),(4.6,1.1)} \draw[popblue, fill=white] \p circle (0.06);
% Étiquettes
\node[align=center] at (0.6,-0.9) {craie \ce{CaCO3}\\ + acide chlorhydrique};
\node[align=center] at (4.6,-0.9) {eau de chaux\\ qui se trouble};
\node[popblue] at (2.6,3.85) {gaz dégagé : \ce{CO2}};
\draw[->, popblue, thick] (2.6,3.65) -- (2.6,3.55);
\end{tikzpicture}
\end{center}
\legende{Dispositif expérimental : l'acide chlorhydrique versé sur la craie provoque une effervescence. Le gaz dégagé est conduit par un tube à dégagement dans l'eau de chaux, qui se trouble : c'est la preuve de la formation de dioxyde de carbone \ce{CO2}, une espèce nouvelle.}
\end{popfigure}

\begin{exemplebox}[Le test à l'eau de chaux]
Le gaz qui trouble l'eau de chaux lors de l'action de l'acide chlorhydrique sur la craie est le \cle{dioxyde de carbone} \ce{CO2}. Ce test est très utilisé au laboratoire : si un gaz incolore trouble l'eau de chaux, alors ce gaz contient du dioxyde de carbone. C'est ainsi qu'on prouve qu'une \textbf{espèce nouvelle} est apparue : la transformation est donc chimique.
\end{exemplebox}

\courssub{Expérience 2 : combustion du ruban de magnésium}

\begin{experience}[Brûler du magnésium dans l'air]
\textbf{Matériel :} un ruban de magnésium (métal gris brillant), une pince en fer, une flamme de bec Bunsen, des lunettes de protection, une coupelle en porcelaine.

\textbf{Protocole :}
\begin{enumerate}
\item Saisir un petit ruban de magnésium avec la pince.
\item Approcher le ruban de la flamme, \textbf{sans regarder directement} la lumière produite (risque de lésion oculaire).
\item Recueillir le résidu solide sur la coupelle et l'observer.
\end{enumerate}

\textbf{Observations :}
\begin{itemize}
\item Le magnésium brûle en émettant une \textbf{lumière blanche très vive} et en dégageant de la chaleur.
\item Il reste une \textbf{cendre blanche et poudreuse}, l'oxyde de magnésium \ce{MgO}, qui ne ressemble en rien au métal gris de départ : elle ne conduit plus le courant, ne brille plus, ne se plie plus.
\end{itemize}

\textbf{Interprétation :} le magnésium a réagi avec le \cle{dioxygène} \ce{O2} de l'air. Le métal magnésium a disparu et une espèce nouvelle, l'oxyde de magnésium, est apparue. La transformation est chimique.

\textbf{Équation-bilan de la réaction :}
\[ \ce{2 Mg(s) + O2(g) -> 2 MgO(s)} \]
\end{experience}

\courssub{Les indices d'une transformation chimique}

Comment, dans la vie courante ou au laboratoire, savoir qu'une transformation chimique vient de se produire ? On se fie à des \cle{indices expérimentaux} observables.

\begin{methode}[Reconnaître une transformation chimique]
\begin{enumerate}
\item \textbf{Observer avant et après} : noter l'aspect, la couleur, l'état des espèces présentes au départ.
\item \textbf{Chercher les indices} pendant la transformation :
\begin{itemize}
\item un \textbf{dégagement gazeux} (effervescence, bulles) ;
\item un \textbf{changement de couleur} durable ;
\item une \textbf{variation de température} (le milieu chauffe ou refroidit) ;
\item l'apparition d'un \textbf{précipité} (solide qui se forme dans un liquide) ;
\item une \textbf{émission de lumière} (flamme, éclat).
\end{itemize}
\item \textbf{Identifier une espèce nouvelle} si possible, grâce à un test d'identification (eau de chaux pour \ce{CO2}, par exemple).
\item \textbf{Conclure} : si au moins une espèce nouvelle est apparue et qu'au moins une espèce de départ a disparu, la transformation est \textbf{chimique}.
\end{enumerate}
\end{methode}

\begin{attention}[Ce qui n'est PAS une transformation chimique]
La \textbf{fonte de la glace}, la \textbf{dissolution du sel} dans l'eau ou l'\textbf{évaporation de l'eau} ne sont \textbf{pas} des réactions chimiques : aucune espèce nouvelle n'apparaît. L'eau reste de l'eau, le sel reste du sel. Ce sont des transformations \textbf{physiques}. Piège classique au BEPC : un changement d'état spectaculaire (bouillonnement de l'eau qui bout) n'est pas un dégagement gazeux d'espèce nouvelle !
\end{attention}

\courssub{Définition et vocabulaire}

\begin{definition}[Transformation chimique, réactifs, produits]
Une \cle{transformation chimique} est une transformation au cours de laquelle des espèces chimiques \textbf{disparaissent} et de \textbf{nouvelles espèces} chimiques \textbf{apparaissent}.
\begin{itemize}
\item Les espèces \textbf{consommées} (qui disparaissent) sont les \cle{réactifs}.
\item Les espèces \textbf{formées} (qui apparaissent) sont les \cle{produits}.
\end{itemize}
La \cle{réaction chimique} est la modélisation de cette transformation à l'échelle microscopique : on la décrit par un schéma
\[ \text{réactifs} \longrightarrow \text{produits}. \]
\end{definition}

Dans nos deux expériences :
\begin{itemize}
\item Expérience 1 : réactifs = carbonate de calcium \ce{CaCO3} et acide chlorhydrique \ce{HCl} ; produits = dioxyde de carbone \ce{CO2}, chlorure de calcium \ce{CaCl2} et eau \ce{H2O}.
\item Expérience 2 : réactifs = magnésium \ce{Mg} et dioxygène \ce{O2} ; produit = oxyde de magnésium \ce{MgO}.
\end{itemize}

\courssub{Interprétation microscopique : les atomes se réorganisent}

Que se passe-t-il à l'échelle des atomes et des molécules ? Prenons l'exemple de la combustion du méthane \ce{CH4} (gaz des marais, proche du gaz naturel) dans le dioxygène, qui donne du dioxyde de carbone et de l'eau :
\[ \ce{CH4 + 2 O2 -> CO2 + 2 H2O} \]

Au cours d'une réaction chimique :
\begin{itemize}
\item les \textbf{liaisons} des molécules des réactifs se \textbf{cassent} ;
\item les \textbf{atomes se réorganisent} pour former les molécules des produits ;
\item les atomes eux-mêmes ne sont \textbf{ni créés, ni détruits, ni transformés} : on retrouve exactement les mêmes atomes avant et après, simplement rangés différemment.
\end{itemize}

\begin{popfigure}
\begin{center}
\begin{tikzpicture}[scale=0.85]
% Réactifs : CH4
\fill[popdark] (0,0) circle (0.30);
\node[white, font=\footnotesize] at (0,0) {C};
\foreach \a/\x/\y in {45/0.55/0.55, 135/-0.55/0.55, 225/-0.55/-0.55, 315/0.55/-0.55}{
  \fill[popmuted!50] (\x,\y) circle (0.18);
  \node[font=\scriptsize] at (\x,\y) {H};
  \draw[thick, popline] (0,0) -- (\x,\y);
}
\node[align=center, font=\small] at (0,-1.15) {\ce{CH4}};
% + 2 O2
\node[font=\Large] at (1.6,0) {$+$};
\fill[poporange] (2.6,0.35) circle (0.26); \fill[poporange] (3.3,0.35) circle (0.26);
\draw[thick, popline] (2.6,0.35) -- (3.3,0.35);
\node[white, font=\footnotesize] at (2.6,0.35) {O}; \node[white, font=\footnotesize] at (3.3,0.35) {O};
\fill[poporange] (2.6,-0.55) circle (0.26); \fill[poporange] (3.3,-0.55) circle (0.26);
\draw[thick, popline] (2.6,-0.55) -- (3.3,-0.55);
\node[white, font=\footnotesize] at (2.6,-0.55) {O}; \node[white, font=\footnotesize] at (3.3,-0.55) {O};
\node[align=center, font=\small] at (2.95,-1.15) {\ce{2 O2}};
% Flèche
\draw[->, very thick, popblue] (4.1,0) -- (5.4,0);
% Produits : CO2
\fill[poporange] (6.2,0.35) circle (0.26); \fill[popdark] (6.95,0.35) circle (0.30); \fill[poporange] (7.7,0.35) circle (0.26);
\draw[thick, popline] (6.2,0.35) -- (6.95,0.35) -- (7.7,0.35);
\node[white, font=\footnotesize] at (6.2,0.35) {O}; \node[white, font=\footnotesize] at (6.95,0.35) {C}; \node[white, font=\footnotesize] at (7.7,0.35) {O};
\node[align=center, font=\small] at (6.95,-1.15) {\ce{CO2}};
\node[font=\Large] at (8.4,0) {$+$};
% 2 H2O
\fill[poporange] (9.3,0.45) circle (0.26);
\fill[popmuted!50] (8.85,0.85) circle (0.18); \fill[popmuted!50] (9.75,0.85) circle (0.18);
\draw[thick, popline] (9.3,0.45) -- (8.85,0.85); \draw[thick, popline] (9.3,0.45) -- (9.75,0.85);
\node[white, font=\footnotesize] at (9.3,0.45) {O};
\node[font=\scriptsize] at (8.85,0.85) {H}; \node[font=\scriptsize] at (9.75,0.85) {H};
\fill[poporange] (9.3,-0.65) circle (0.26);
\fill[popmuted!50] (8.85,-0.25) circle (0.18); \fill[popmuted!50] (9.75,-0.25) circle (0.18);
\draw[thick, popline] (9.3,-0.65) -- (8.85,-0.25); \draw[thick, popline] (9.3,-0.65) -- (9.75,-0.25);
\node[white, font=\footnotesize] at (9.3,-0.65) {O};
\node[font=\scriptsize] at (8.85,-0.25) {H}; \node[font=\scriptsize] at (9.75,-0.25) {H};
\node[align=center, font=\small] at (9.3,-1.15) {\ce{2 H2O}};
\end{tikzpicture}
\end{center}
\legende{Modèle microscopique de la combustion du méthane : les molécules de \ce{CH4} et de \ce{O2} se cassent, puis les atomes (C en noir, H en gris, O en orange) se réorganisent pour former \ce{CO2} et \ce{H2O}. Compte les atomes : 1 C, 4 H et 4 O avant, et 1 C, 4 H et 4 O après. Aucun atome n'a été créé ni détruit.}
\end{popfigure}

\begin{aretenir}[L'essentiel : transformation chimique]
\begin{itemize}
\item Une transformation chimique s'accompagne de la \textbf{disparition de réactifs} et de l'\textbf{apparition de produits}.
\item On la repère par des indices : dégagement gazeux, changement de couleur, variation de température, précipité, lumière.
\item À l'échelle microscopique, les atomes des réactifs \textbf{se réorganisent} : ils sont conservés en nombre et en nature.
\end{itemize}
\end{aretenir}

\begin{exoresolu}[Physique ou chimique ?]
\textbf{Énoncé.} Pour chacune des transformations suivantes, dire s'il s'agit d'une transformation physique ou chimique, en justifiant :
\begin{enumerate}
\item On fait fondre du fer pour fabriquer une marmite.
\item Du bois brûle dans un foyer.
\item On dissout du sucre dans l'eau du « foléré ».
\item Du lait caillé se forme quand on ajoute du jus de citron au lait chaud.
\end{enumerate}
\tcblower
\textbf{Solution détaillée.}
\begin{enumerate}
\item \textbf{Physique} : le fer change d'état (solide $\to$ liquide $\to$ solide) mais reste du fer ; aucune espèce nouvelle n'apparaît.
\item \textbf{Chimique} : le bois et le dioxygène disparaissent ; des espèces nouvelles apparaissent (dioxyde de carbone, vapeur d'eau, cendres). Indices : flamme, chaleur, fumée.
\item \textbf{Physique} : le sucre se disperse dans l'eau mais reste du sucre ; on peut le récupérer par évaporation. Aucune espèce nouvelle.
\item \textbf{Chimique} : un solide nouveau (le caillé, caséine) apparaît dans le liquide, avec changement d'aspect et de goût irréversibles : une espèce nouvelle s'est formée.
\end{enumerate}
\end{exoresolu}

\begin{savaistu}[Le secret du puff-puff qui gonfle]
Le dégagement de \ce{CO2} n'est pas qu'une affaire de laboratoire ! Quand la maman prépare les beignets ou le « puff-puff », la levure chimique (bicarbonate de sodium + acide) incorporée à la pâte réagit à la chaleur et libère du \textbf{dioxyde de carbone}. Les bulles de gaz, prisonnières de la pâte, la font \textbf{lever} et donnent au beignet son moelleux. C'est une réaction chimique... que tu manges !
\end{savaistu}

\courssub{Conservation de la masse et loi de Lavoisier}

Si les atomes ne sont ni créés ni détruits, la masse totale doit rester constante. C'est ce qu'a démontré Antoine Lavoisier au XVIII\textsuperscript{e} siècle.

\begin{experience}[Vérification de la conservation de la masse (précipitation)]
\textbf{Matériel :} une balance électronique de précision (0,01 g), un erlenmeyer avec bouchon, deux petits tubes à essai, du chlorure de sodium \ce{NaCl} (sel de table), du nitrate d'argent \ce{AgNO3} (solution incolore).
\textbf{Protocole :}
\begin{enumerate}
\item Mettre la solution de \ce{NaCl} dans l'erlenmeyer et la solution de \ce{AgNO3} dans un petit tube à essai. Placer le tube dans l'erlenmeyer sans verser. Boucher l'erlenmeyer.
\item Peser l'ensemble : masse initiale $m_i$.
\item Renverser l'erlenmeyer pour mélanger les solutions. Un précipité blanc (\ce{AgCl}) se forme instantanément.
\item Peser à nouveau l'ensemble fermé : masse finale $m_f$.
\end{enumerate}
\textbf{Observation :} $m_f = m_i$ (à la précision de la balance près). La masse ne change pas malgré la formation d'un solide nouveau.
\textbf{Conclusion :} Au cours d'une transformation chimique dans un système \textbf{fermé}, la masse totale reste constante.
\end{experience}

\begin{propriete}[Loi de Lavoisier (Conservation de la masse)]
\textbf{Au cours d'une transformation chimique réalisée dans un système fermé, la masse du système ne change pas.}
\[ m_{\text{réactifs}} = m_{\text{produits}} \]
Cette loi est une conséquence directe de la conservation des atomes (même nombre, même nature, donc même masse).
\end{propriete}

\begin{attention}[Système ouvert vs fermé]
Si un gaz s'échappe (système ouvert), la masse \emph{mesurée} diminue. Exemple : la craie + acide dans un bécher non bouché : le \ce{CO2} part dans l'air, la masse finale est plus faible. La loi de Lavoisier ne s'applique qu'aux systèmes \textbf{fermés} (aucun échange de matière avec l'extérieur).
\end{attention}

\courssub{Équation-bilan et équilibrage}

L'équation-bilan est l'écriture symbolique de la réaction chimique. Elle doit respecter la conservation des atomes (loi de Lavoisier).

\begin{definition}[Équation-bilan]
Une \cle{équation-bilan} s'écrit :
\[ \text{Réactifs} \longrightarrow \text{Produits} \]
\begin{itemize}
\item Les formules chimiques des réactifs sont à gauche, celles des produits à droite.
\item L'état physique est précisé entre parenthèses : \ce{(s)} solide, \ce{(l)} liquide, \ce{(g)} gaz, \ce{(aq)} aqueux (dissous dans l'eau).
\item Les \textbf{nombres stœchiométriques} (coefficients devant les formules) indiquent les proportions en nombre de molécules (ou de moles). Ils sont ajustés pour \textbf{équilibrer} l'équation : même nombre d'atomes de chaque élément des deux côtés de la flèche.
\end{itemize}
\end{definition}

\begin{methode}[Équilibrer une équation-bilan (méthode par inspection)]
\begin{enumerate}
\item \textbf{Écrire les formules} correctes des réactifs et des produits (ne jamais modifier les indices dans les formules !).
\item \textbf{Dresser l'inventaire} des atomes : compter les atomes de chaque élément à gauche et à droite.
\item \textbf{Équilibrer} les éléments qui ne sont ni H ni O (souvent C, puis métaux, puis non-métaux).
\item \textbf{Équilibrer l'oxygène} \ce{O}.
\item \textbf{Équilibrer l'hydrogène} \ce{H}.
\item \textbf{Vérifier} l'inventaire final. Simplifier les coefficients si possible (diviser par le PGCD).
\end{enumerate}
\end{methode}

\begin{exemplebox}[Équilibrage de la combustion du butane (gaz de cuisine)]
Le butane \ce{C4H10} brûle dans le dioxygène pour donner \ce{CO2} et \ce{H2O}.
\begin{enumerate}
\item Ébauche : \ce{C4H10 + O2 -> CO2 + H2O}
\item Inventaire initial : Gauche : 4 C, 10 H, 2 O. Droite : 1 C, 2 H, 3 O.
\item Carbone : mettre 4 devant \ce{CO2}. \ce{C4H10 + O2 -> 4 CO2 + H2O}
\item Hydrogène : mettre 5 devant \ce{H2O}. \ce{C4H10 + O2 -> 4 CO2 + 5 H2O}
\item Oxygène : à droite $4\times2 + 5\times1 = 13$ atomes O. Il faut $13/2$ molécules \ce{O2} à gauche.
\item \ce{C4H10 + 13/2 O2 -> 4 CO2 + 5 H2O} $\to$ multiplier par 2 pour éviter les fractions :
\item \textbf{Équation finale :} \ce{2 C4H10 + 13 O2 -> 8 CO2 + 10 H2O}
\end{enumerate}
Vérification : Gauche : 8 C, 20 H, 26 O. Droite : 8 C, 20 H, $16+10=26$ O. \textbf{Équilibrée.}
\end{exemplebox}

\begin{exoresolu}[Équilibrage : action de l'acide sur la craie]
\textbf{Énoncé.} Écrire et équilibrer l'équation-bilan de la réaction entre le carbonate de calcium solide \ce{CaCO3} et l'acide chlorhydrique aqueux \ce{HCl}. Les produits sont le chlorure de calcium aqueux \ce{CaCl2}, le dioxyde de carbone gazeux \ce{CO2} et l'eau liquide.
\tcblower
\textbf{Solution.}
\begin{enumerate}
\item Formules : \ce{CaCO3(s) + HCl(aq) -> CaCl2(aq) + CO2(g) + H2O(l)}
\item Inventaire : Gauche : 1 Ca, 1 C, 3 O + 1 H, 1 Cl. Droite : 1 Ca, 1 C, 3 O, 2 H, 2 Cl.
\item Ca, C, O (hors H/O) : Ca et C OK. Cl : 1 à G, 2 à D $\to$ mettre 2 devant \ce{HCl}.
\item \ce{CaCO3 + 2 HCl -> CaCl2 + CO2 + H2O}
\item H : 2 à G, 2 à D (dans \ce{H2O}). O : 3+0=3 à G, 2+1+1=4 à D ? Attendre.
\item Re-vérif O : \ce{CaCO3} (3 O) + 2\ce{HCl} (0 O) = 3 O. Produits : \ce{CaCl2} (0) + \ce{CO2} (2) + \ce{H2O} (1) = 3 O. \textbf{OK.}
\item \textbf{Équation finale :} \ce{CaCO3(s) + 2 HCl(aq) -> CaCl2(aq) + CO2(g) + H2O(l)}
\end{enumerate}
\end{exoresolu}

\courssub{Combustions complète et incomplète}

La combustion est une réaction chimique rapide entre un \cle{combustible} (souvent du carbone \ce{C} et de l'hydrogène \ce{H}) et un \cle{comburant} (le dioxygène \ce{O2} de l'air), qui dégage de la \textbf{chaleur} et de la \textbf{lumière}.

\begin{definition}[Combustion complète]
Lorsque le combustible brûle dans une \textbf{quantité suffisante de dioxygène}, il se transforme totalement en \textbf{dioxyde de carbone} \ce{CO2} et \textbf{eau} \ce{H2O}. La flamme est \textbf{bleue} (gaz) ou claire, non fumante. C'est le rendement énergétique maximal.
\end{definition}

\begin{definition}[Combustion incomplète]
Lorsque le combustible brûle dans une \textbf{quantité insuffisante de dioxygène} (manque d'air), il se forme du \textbf{monoxyde de carbone} \ce{CO} (gaz incolore, inodore, très toxique) et/ou du \textbf{carbone} \ce{C} (suie, particules noires), en plus de \ce{CO2} et \ce{H2O}. La flamme est \textbf{jaune-orangée, fumeuse}. Le rendement énergétique est plus faible.
\end{definition}

\begin{center}
\begin{tabularx}{\linewidth}{|l|X|X|}\hline
\rowcolor{popblueL}
\textbf{Caractéristique} & \textbf{Combustion complète} & \textbf{Combustion incomplète} \\ \hline
\textbf{Dioxygène} & Excès ou quantité stœchiométrique & Déficient (manque d'air) \\ \hline
\textbf{Produits principaux} & \ce{CO2} + \ce{H2O} & \ce{CO} + \ce{C} (suie) + \ce{CO2} + \ce{H2O} \\ \hline
\textbf{Aspect flamme} & Bleue (gaz), claire, stable & Jaune, orangée, vacillante, fumante \\ \hline
\textbf{Exemple courant} & Gazinière bien réglée (bleu), chalumeau & Brasero, feu de bois étouffé, moteur mal réglé, groupe électrogène dans pièce close \\ \hline
\textbf{Danger} & Brûlures, incendie & \textbf{Intoxication mortelle au \ce{CO}}, encrassement \\ \hline
\end{tabularx}
\end{center}

\begin{exemplebox}[Équations de combustion du méthane \ce{CH4} (gaz naturel)]
\begin{itemize}
\item \textbf{Complète :} \ce{CH4 + 2 O2 -> CO2 + 2 H2O} \quad (Flamme bleue, gazinière bien réglée)
\item \textbf{Incomplète (manque d'O2) :} \ce{2 CH4 + 3 O2 -> 2 CO + 4 H2O} \quad (Formation de \ce{CO})
\item \textbf{Très incomplète :} \ce{CH4 + O2 -> C + 2 H2O} \quad (Formation de suie, flamme jaune)
\end{itemize}
\end{exemplebox}

\courssub{Le monoxyde de carbone : un danger invisible}

Le \cle{monoxyde de carbone} (\ce{CO}) est le produit majeur des combustions incomplètes. C'est un gaz \textbf{incolore, inodore, non irritant} et \textbf{légèrement plus léger que l'air} (il se diffuse partout dans la pièce).

\begin{propriete}[Toxicité du \ce{CO} : fixation sur l'hémoglobine]
Dans le sang, l'hémoglobine (Hb) transporte le dioxygène \ce{O2} vers les cellules. Le \ce{CO} se fixe sur l'hémoglobine \textbf{200 à 250 fois plus fortement} que l'\ce{O2}, formant de la \textbf{carboxyhémoglobine} (\ce{HbCO}).
\[ \ce{Hb + CO -> HbCO} \quad (\text{fixation quasi irréversible}) \]
\emph{Conséquence :} le sang ne peut plus transporter l'oxygène $\to$ asphyxie cellulaire $\to$ décès rapide si concentration élevée.
\end{propriete}

\begin{attention}[Symptômes et pièges]
\begin{itemize}
\item \textbf{Faible concentration :} maux de tête, nausées, vertiges, fatigue inhabituelle (souvent confondus avec grippe ou malaise vague).
\item \textbf{Forte concentration :} perte de connaissance en quelques minutes, coma, mort.
\item \textbf{Piège mortel :} pas d'odeur, pas de fumée visible (le \ce{CO} est incolore), pas de toux. La victime s'endort et ne se réveille pas.
\end{itemize}
\end{attention}

\begin{savaistu}[Tragédies évitables au Cameroun]
Chaque année, au Cameroun, des familles entières décèdent dans leur sommeil à cause d'un brasero, d'un groupe électrogène ou d'une moto-taxi laissés en marche dans une pièce close ou mal ventilée (portes/fenêtres fermées à cause de la pluie ou du froid). Le \ce{CO} remplit la pièce en silence. \textbf{Aérer sauve des vies.}
\end{savaistu}

\courssub{Sécurité et gestes responsables autour des combustions}

\begin{methode}[Prévention des accidents liés au monoxyde de carbone]
\begin{enumerate}
\item \textbf{Jamais} de brasero, fourneau à charbon, groupe électrogène, moto ou voiture en marche dans un lieu \textbf{fermé ou mal ventilé} (chambre, salon, garage, couloir). Sortir le groupe électrogène \textbf{à l'extérieur} (au moins 5 m des ouvertures).
\item \textbf{Aérer} en permanence : ouvrir grand les fenêtres et portes lors de l'utilisation de tout appareil à combustion (gazinière, chauffage, lampe à pétrole). Créer un courant d'air traversant.
\item \textbf{Entretenir} les appareils : brûleurs de gazinière propres (flamme bleue), pot d'échappement moto/voiture non bouché, révision groupe électrogène.
\item \textbf{Installer un détecteur de \ce{CO}} (norme EN 50291) dans les chambres et couloirs. Il sonne avant que le taux ne devienne mortel.
\item \textbf{En cas de suspicion d'intoxication (maux de tête soudains, nausées, plusieurs personnes touchées) :}
\begin{itemize}
\item Ouvrir \textbf{immédiatement} toutes les portes et fenêtres.
\item \textbf{Évacuer} les lieux au plus vite.
\item Appeler les secours (\textbf{112} ou \textbf{119} au Cameroun) ou emmener les victimes à l'hôpital (service d'urgence, oxygénothérapie).
\item Ne pas réintégrer le lieu avant aération complète et recherche de la source.
\end{itemize}
\end{enumerate}
\end{methode}

\begin{aretenir}[L'essentiel : Réaction, Équation, Combustion, Sécurité]
\begin{itemize}
\item \textbf{Transformation chimique} : disparition de réactifs, apparition de produits (indices : gaz, couleur, T°, précipité, lumière).
\item \textbf{Loi de Lavoisier} : masse conservée en système fermé (atomes réorganisés, non créés/détruits).
\item \textbf{Équation-bilan} : formules exactes + coefficients stœchiométriques pour équilibrer les atomes (méthode : C, puis O, puis H).
\item \textbf{Combustion complète} (excès \ce{O2}) $\to$ \ce{CO2} + \ce{H2O} (flamme bleue, max énergie).
\item \textbf{Combustion incomplète} (manque \ce{O2}) $\to$ \ce{CO} + \ce{C} + \ce{CO2} + \ce{H2O} (flamme jaune, suie, \textbf{\ce{CO} mortel}).
\item \textbf{Sécurité} : Aérer, sortir les groupes électrogènes, détecteur \ce{CO}, connaître le 112/119.
\end{itemize}
\end{aretenir}

\begin{exercice}[Entraînement BEPC - La réaction chimique]
\begin{enumerate}
\item \textbf{Équilibrage.} Équilibrer les équations-bilan suivantes :
\begin{enumerate}
\item \ce{Fe + O2 -> Fe2O3} (rouille)
\item \ce{C3H8 + O2 -> CO2 + H2O} (propane, gaz de camping)
\item \ce{Al + HCl -> AlCl3 + H2} (aluminium + acide)
\end{enumerate}
\item \textbf{Combustion.} On brûle du charbon de bois (principalement du carbone \ce{C}) dans un brasero.
\begin{enumerate}
\item Écrire l'équation-bilan de la combustion \textbf{complète} du carbone.
\item Écrire l'équation-bilan d'une combustion \textbf{incomplète} du carbone formant du \ce{CO}.
\item Expliquer pourquoi il est mortel de dormir avec un brasero allumé dans une chambre fermée, en utilisant les notions de combustion incomplète et de toxicité du \ce{CO}.
\end{enumerate}
\item \textbf{Masse.} On introduit \SI{10,0}{g} de marbre (\ce{CaCO3}) dans \SI{50,0}{g} de solution d'acide chlorhydrique dans un erlenmeyer bouché posé sur une balance. La masse initiale est \SI{60,0}{g}. Après réaction complète, on observe qu'il reste du marbre non réagi. La masse finale affichée est \SI{60,0}{g}. Expliquer pourquoi la masse n'a pas changé malgré la réaction et la formation de gaz.
\item \textbf{Identification.} On chauffe un mélange de poudre de fer et de soufre. Il se forme un solide noir, non magnétique, le sulfure de fer \ce{FeS}.
\begin{enumerate}
\item Citer les réactifs et le produit.
\item Est-ce une transformation physique ou chimique ? Justifier par deux arguments.
\end{enumerate}
\end{enumerate}
\end{exercice}

\begin{corrige}[Exercice : Entraînement BEPC]
\begin{enumerate}
\item \textbf{Équilibrage.}
\begin{enumerate}
\item Inventaire : Fe (1/2), O (2/3). Mettre 2 devant \ce{Fe} $\to$ \ce{2 Fe + O2 -> Fe2O3}. O : 2/3. PPCM 6. \ce{2 Fe + 3 O2 -> 2 Fe2O3} ? Non. \ce{4 Fe + 3 O2 -> 2 Fe2O3}. \textbf{Vérif : Fe 4/4, O 6/6.}
\item \ce{C3H8 + O2 -> CO2 + H2O}. C : 3 $\to$ \ce{3 CO2}. H : 8 $\to$ \ce{4 H2O}. O : à D $3\times2 + 4 = 10$. À G : \ce{5 O2}. \textbf{\ce{C3H8 + 5 O2 -> 3 CO2 + 4 H2O}.}
\item \ce{Al + HCl -> AlCl3 + H2}. Al : 1/1. Cl : 1/3 $\to$ \ce{3 HCl}. H : 3/2. PPCM 6. \ce{2 Al + 6 HCl -> 2 AlCl3 + 3 H2}. \textbf{Vérif : Al 2/2, Cl 6/6, H 6/6.}
\end{enumerate}
\item \textbf{Combustion du charbon.}
\begin{enumerate}
\item \ce{C + O2 -> CO2} (Complète, flamme bleue à la base, rougeoyante au sommet).
\item \ce{2 C + O2 -> 2 CO} (Incomplète, flamme bleue faible, production de \ce{CO}).
\item Dans une chambre fermée, le dioxygène diminue (consommé par la combustion et la respiration). La combustion devient \textbf{incomplète} et produit du \ce{CO}. Ce gaz incolore, inodore, se diffuse dans la pièce. Il se fixe sur l'hémoglobine du sang à la place de l'\ce{O2}, empêchant le transport de l'oxygène vers le cerveau et le cœur. Les victimes s'endorment et meurent par asphyxie cellulaire sans s'en rendre compte.
\end{enumerate}
\item \textbf{Conservation de la masse.} L'erlenmeyer est \textbf{bouché} (système fermé). Le gaz \ce{CO2} formé reste piégé dans le ballon. Selon la loi de Lavoisier, la masse totale du système fermé reste constante. La masse finale égale la masse initiale (\SI{60,0}{g}). Le marbre restant prouve que l'acide était le réactif limitant.
\item \textbf{Fer + Soufre.}
\begin{enumerate}
\item Réactifs : Fer \ce{Fe} (poudre grise, magnétique) et Soufre \ce{S} (poudre jaune). Produit : Sulfure de fer \ce{FeS} (solide noir, non magnétique).
\item \textbf{Transformation chimique.} Arguments : 1) Apparition d'une espèce nouvelle (\ce{FeS}) aux propriétés différentes (couleur noire, non magnétique). 2) Disparition des réactifs (le fer n'est plus magnétique, le soufre jaune a disparu). 3) Émission de lumière/chaleur lors de l'inflammation (réaction exothermique).
\end{enumerate}
\end{enumerate}
\end{corrige}

\coursec{Conservation de la masse et loi de Lavoisier}

Dans la section précédente, nous avons appris à reconnaître une \cle{transformation chimique} : de nouvelles espèces apparaissent (un gaz, un précipité, un changement de couleur), tandis que les espèces de départ disparaissent. Mais une question fondamentale reste en suspens : lorsque la craie réagit avec l'acide, lorsque le charbon brûle dans le brasero, \emph{que devient la matière} ? La masse change-t-elle ? C'est à cette question qu'un savant français, Antoine Laurent de \cle{Lavoisier}, a répondu il y a plus de deux siècles, en introduisant un outil révolutionnaire dans le laboratoire de chimie : la \cle{balance}. Cette section raconte cette découverte et en tire la loi la plus importante de toute la chimie.

\courssub{L'expérience historique de Lavoisier (vers 1774)}

\begin{savaistu}[Lavoisier, père de la chimie moderne]
Antoine Laurent de Lavoisier (1743-1794) était un savant français passionné de mesures précises. À une époque où les chimistes raisonnaient surtout de manière qualitative (« ça brûle », « ça fume »), lui décida de \textbf{tout peser}, avant et après chaque transformation. Grâce à cette rigueur, il découvrit le rôle du dioxygène dans les combustions et énonça la loi de conservation de la masse. On le considère comme le \textbf{père de la chimie moderne}. Destin tragique : fermier général (collecteur d'impôts) sous l'Ancien Régime, il fut arrêté pendant la Révolution française et \textbf{guillotiné le 8 mai 1794}, à 50 ans. Le mathématicien Lagrange aurait déclaré : « Il ne leur a fallu qu'un instant pour faire tomber cette tête, et cent années peut-être ne suffiront pas pour en reproduire une semblable. »
\end{savaistu}

\begin{experience}[L'expérience de la cornue de mercure (vers 1774)]
\textbf{Matériel :} une cornue en verre (ballon à long col recourbé) contenant du \cle{mercure} liquide et un volume d'air connu, reliée à une cloche renversée sur une cuve à mercure ; un fourneau pour chauffer.

\textbf{Protocole :}
\begin{enumerate}
\item Lavoisier place du mercure dans la cornue et mesure le volume d'air enfermé dans le dispositif.
\item Il chauffe le mercure à une température proche de son ébullition, pendant \textbf{douze jours}, dans le dispositif \textbf{fermé}.
\item Il observe et pèse le contenu avant, pendant et après le chauffage.
\end{enumerate}

\textbf{Observations :}
\begin{itemize}
\item À la surface du mercure apparaissent des paillettes \textbf{rouges} : une nouvelle espèce chimique s'est formée, l'\cle{oxyde de mercure} (appelé autrefois « précipité rouge »).
\item Le volume d'air dans la cloche a \textbf{diminué} d'environ un cinquième : une partie de l'air a été « consommée ».
\item L'air restant éteint une bougie et ne permet plus la respiration : c'est de l'air « appauvri » (essentiellement du diazote).
\item Surtout : la \textbf{masse totale du dispositif fermé n'a pas changé} d'un milligramme entre le début et la fin de l'expérience.
\end{itemize}

\textbf{Interprétation :} le mercure a réagi avec le dioxygène de l'air pour former l'oxyde de mercure rouge :
\[ \ce{2Hg + O2 -> 2HgO} \]
La masse gagnée par le mercure (devenu oxyde) est \emph{exactement égale} à la masse de dioxygène disparue de l'air. La matière n'a ni disparu ni été créée : elle a simplement \textbf{changé de partenaire}.
\end{experience}

\begin{popfigure}
\begin{center}
\begin{tikzpicture}[scale=1, every node/.style={font=\small}]
% Frise historique
\draw[very thick, popblue, -{Stealth}] (0,0) -- (12.5,0);
\foreach \x/\year in {1/1743, 4.5/1774, 8/1789, 11/1794} {
  \draw[popblue, thick] (\x,-0.12) -- (\x,0.12);
  \node[below, popdark] at (\x,-0.15) {\textbf{\year}};
}
% Événements
\node[above, align=center, popdark] at (1,0.25) {Naissance\\à Paris};
\fill[poporange] (4.5,0) circle (0.12);
\node[above, align=center, popdark] at (4.5,0.25) {Expérience de la cornue :\\conservation de la masse};
\fill[popgreen] (8,0) circle (0.12);
\node[above, align=center, popdark] at (8,0.25) {Publication du\\ \emph{Traité élémentaire de chimie}};
\fill[poppink] (11,0) circle (0.12);
\node[above, align=center, popdark] at (11,0.25) {Guillotiné pendant\\la Révolution};
\end{tikzpicture}
\end{center}
\legende{Frise historique : les grandes dates de la vie d'Antoine Laurent de Lavoisier (1743-1794).}
\end{popfigure}

\courssub{Expérience de classe : la pesée avant et après réaction}

Reproduisons l'esprit de la démarche de Lavoisier avec du matériel simple, que l'on trouve dans n'importe quel laboratoire de collège.

\begin{experience}[Craie et acide en bouteille fermée]
\textbf{Matériel :} une bouteille en plastique avec son bouchon, un morceau de craie, un petit flacon (ou un tube) contenant de l'acide chlorhydrique dilué, une balance électronique.

\textbf{Protocole :}
\begin{enumerate}
\item Déposer la craie au fond de la bouteille et y introduire délicatement le flacon d'acide \textbf{sans le renverser}.
\item Boucher la bouteille et poser l'ensemble sur la balance. Noter la masse : $m_1$.
\item Sans ouvrir la bouteille, l'incliner pour verser l'acide sur la craie : une \textbf{effervescence} se produit (dégagement de dioxyde de carbone \ce{CO2}).
\item Attendre la fin de la réaction et lire la nouvelle masse : $m_2$.
\end{enumerate}

\textbf{Observation :} malgré l'effervescence spectaculaire, la balance indique \textbf{exactement la même masse} : $m_2 = m_1$.

\textbf{Interprétation :} la réaction \ce{CaCO3 + 2HCl -> CaCl2 + H2O + CO2} a bien eu lieu (le gaz le prouve), mais comme la bouteille est \textbf{fermée}, aucune matière n'a pu sortir ni entrer. Le gaz produit est resté prisonnier : sa masse est donc toujours comptée par la balance.
\end{experience}

\begin{experience}[La même expérience en système ouvert]
Reprenons le même montage, mais cette fois \textbf{sans le bouchon} (ou en ouvrant la bouteille juste après avoir versé l'acide).

\textbf{Observation :} la masse affichée par la balance \textbf{diminue} au fur et à mesure de l'effervescence, puis se stabilise à une valeur inférieure à $m_1$.

\textbf{Interprétation :} le dioxyde de carbone \ce{CO2} formé \textbf{s'échappe dans l'air} de la salle. La balance ne pèse plus que ce qui reste dans la bouteille. La masse « perdue » sur l'écran est précisément la masse du gaz parti. Rien n'a disparu : la matière a simplement \textbf{changé de place}.
\end{experience}

\begin{popfigure}
\begin{center}
\begin{tikzpicture}[scale=0.9, every node/.style={font=\small}]
% ===== Balance 1 : système fermé =====
\begin{scope}[shift={(0,0)}]
% pied et colonne
\draw[fill=popmuted!60] (-0.6,0) rectangle (0.6,0.25);
\draw[fill=popmuted!60] (-0.08,0.25) rectangle (0.08,1.6);
% fléau horizontal
\draw[very thick, popdark] (-2.3,1.6) -- (2.3,1.6);
\fill[popdark] (0,1.6) circle (0.07);
% aiguille verticale (inchangée)
\draw[thick, poporange, -{Stealth}] (0,1.6) -- (0,2.15);
\node[poporange, above] at (0,2.15) {\textbf{aiguille au zéro}};
% plateau gauche
\draw[thick, popdark] (-2.3,1.6) -- (-2.3,1.1);
\draw[fill=popblueL] (-2.9,1.0) arc (180:360:0.6) -- cycle;
% bouteille fermée sur plateau gauche
\draw[fill=popblue!20, draw=popblue, thick] (-2.75,1.0) rectangle (-1.85,2.3);
\draw[fill=popblue!20, draw=popblue, thick] (-2.45,2.3) rectangle (-2.15,2.55);
\draw[fill=poppink, draw=popdark] (-2.5,2.55) rectangle (-2.1,2.7);
\node[popdark] at (-2.3,0.7) {bouteille \textbf{fermée}};
% plateau droit : masses
\draw[thick, popdark] (2.3,1.6) -- (2.3,1.1);
\draw[fill=popblueL] (1.7,1.0) arc (180:360:0.6) -- cycle;
\draw[fill=popmuted!70, draw=popdark] (2.05,1.0) rectangle (2.55,1.45);
\node[popdark] at (2.3,1.22) {\footnotesize $m$};
\node[popdark] at (2.3,0.7) {avant = après};
\node[below, popdark, align=center] at (0,-0.35) {\textbf{a) Système fermé :} la masse totale ne change pas};
\end{scope}
% ===== Balance 2 : système ouvert =====
\begin{scope}[shift={(8.5,0)}]
\draw[fill=popmuted!60] (-0.6,0) rectangle (0.6,0.25);
\draw[fill=popmuted!60] (-0.08,0.25) rectangle (0.08,1.6);
% fléau penché (plateau gauche remonté car masse perdue)
\draw[very thick, popdark] (-2.3,1.85) -- (2.3,1.35);
\fill[popdark] (0,1.6) circle (0.07);
% aiguille déviée
\draw[thick, poppink, -{Stealth}] (0,1.6) -- (0.35,2.1);
\node[poppink, above right, align=left] at (0.3,2.05) {\textbf{aiguille déviée}};
% plateau gauche remonté
\draw[thick, popdark] (-2.3,1.85) -- (-2.3,1.35);
\draw[fill=popblueL] (-2.9,1.25) arc (180:360:0.6) -- cycle;
% bouteille ouverte
\draw[fill=popblue!20, draw=popblue, thick] (-2.75,1.25) rectangle (-1.85,2.55);
\draw[fill=popblue!20, draw=popblue, thick] (-2.45,2.55) rectangle (-2.15,2.8);
% bulles de CO2 qui s'échappent
\fill[popgreen!60] (-2.3,3.0) circle (0.09);
\fill[popgreen!60] (-2.15,3.3) circle (0.07);
\fill[popgreen!60] (-2.4,3.55) circle (0.05);
\node[popgreen!60!popdark, right] at (-2.05,3.3) {\footnotesize \ce{CO2} s'échappe};
\node[popdark] at (-2.3,0.95) {bouteille \textbf{ouverte}};
% plateau droit descendu
\draw[thick, popdark] (2.3,1.35) -- (2.3,0.85);
\draw[fill=popblueL] (1.7,0.75) arc (180:360:0.6) -- cycle;
\draw[fill=popmuted!70, draw=popdark] (2.05,0.75) rectangle (2.55,1.2);
\node[popdark] at (2.3,0.97) {\footnotesize $m$};
\node[popdark] at (2.3,0.45) {masse diminuée};
\node[below, popdark, align=center] at (0,-0.35) {\textbf{b) Système ouvert :} le gaz part, la masse affichée diminue};
\end{scope}
\end{tikzpicture}
\end{center}
\legende{Pesée de la réaction craie + acide. a) En système fermé, le gaz reste prisonnier : la masse est conservée. b) En système ouvert, le \ce{CO2} s'échappe : la balance indique une masse plus faible.}
\end{popfigure}

\courssub{La loi de Lavoisier}

\begin{propriete}[Loi de Lavoisier : conservation de la masse]
Dans une \cle{transformation chimique} réalisée en \cle{système fermé} (aucun échange de matière avec l'extérieur), la \textbf{masse totale se conserve} :
\[ m_{\text{total avant}} = m_{\text{total après}} \qquad \text{ou encore} \qquad \sum m_{\text{réactifs}} = \sum m_{\text{produits}} \]
Cette loi a été \textbf{établie expérimentalement} par Lavoisier grâce à des pesées de précision ; elle est aujourd'hui admise comme un \textbf{principe fondateur de la chimie}. Aucune démonstration théorique n'est exigible au BEPC, mais tu dois savoir \emph{la justifier par une expérience de pesée} et l'utiliser dans des calculs.
\end{propriete}

\begin{attention}[Un système fermé n'est pas un système ouvert]
La loi de Lavoisier ne s'applique \textbf{que si le système est fermé} : ni matière entrante, ni matière sortante. Une bouteille bouchée, une cornue scellée, un ballon hermétique sont des systèmes fermés. Un feu de bois dans la cour, une casserole sur le feu, une bouteille débouchée sont des systèmes \textbf{ouverts} : les gaz peuvent s'échapper et la masse mesurée peut diminuer (ou augmenter si un gaz de l'air est capté). Avant d'appliquer la loi, pose-toi toujours la question : \emph{le système est-il fermé ?}
\end{attention}

\courssub{Conséquence microscopique : les atomes se conservent}

Pourquoi la masse se conserve-t-elle ? La réponse se trouve à l'échelle \cle{microscopique}, celle des atomes.

Au cours d'une transformation chimique, les molécules des réactifs se « cassent » et leurs atomes se \textbf{réarrangent} pour former de nouvelles molécules, celles des produits. Mais :
\begin{itemize}
\item \textbf{aucun atome ne disparaît} ;
\item \textbf{aucun atome n'est créé} ;
\item \textbf{aucun atome ne change de nature} (un atome de carbone reste un atome de carbone).
\end{itemize}

\begin{aretenir}[Conservation des atomes]
Lors d'une transformation chimique, il y a \textbf{conservation du nombre et de la nature des atomes} : chaque atome présent dans les réactifs se retrouve dans les produits. Comme chaque atome garde sa masse, la masse totale est nécessairement conservée. C'est l'explication microscopique de la loi de Lavoisier.
\end{aretenir}

\begin{exemplebox}[Le charbon qui « disparaît » dans le brasero]
Quand maman fait braiser le poisson au charbon de bois, le charbon semble \emph{disparaître} : il ne reste qu'un peu de cendre. Où est passée la matière ? Le charbon (essentiellement du carbone \ce{C}) réagit avec le dioxygène de l'air : \ce{C + O2 -> CO2}. Le dioxyde de carbone, \textbf{gazeux}, s'échappe dans l'atmosphère avec la chaleur. Si l'on pesait \emph{tout} — le charbon, l'air consommé, le gaz produit et les cendres — on retrouverait exactement la même masse totale. La disparition du charbon n'est qu'\textbf{apparente} : les produits sont partis en fumée ! C'est exactement ce que résumait Lavoisier : \emph{« Rien ne se perd, rien ne se crée, tout se transforme. »}
\end{exemplebox}

\courssub{Méthode et application numérique}

\begin{methode}[Interpréter une pesée avant/après réaction]
\begin{enumerate}
\item \textbf{Identifier le système} : le dispositif est-il \cle{fermé} (bouché, scellé) ou \cle{ouvert} (gaz libre de sortir ou d'entrer) ?
\item \textbf{Si le système est fermé} : appliquer la loi de Lavoisier, $m_{\text{avant}} = m_{\text{après}}$. Toute différence apparente provient d'une erreur de mesure.
\item \textbf{Si le système est ouvert} : la différence de masse correspond à la masse de matière échangée avec l'extérieur.
   \begin{itemize}
   \item Masse diminuée $\Rightarrow$ un \textbf{gaz produit s'est échappé} : $m_{\text{gaz}} = m_{\text{avant}} - m_{\text{après}}$.
   \item Masse augmentée $\Rightarrow$ un \textbf{gaz de l'air a été capté} (souvent \ce{O2}) : $m_{\text{gaz}} = m_{\text{après}} - m_{\text{avant}}$.
   \end{itemize}
\item \textbf{Conclure} en précisant que la loi de Lavoisier reste valable si l'on tient compte de \emph{toute} la matière, y compris les gaz.
\end{enumerate}
\end{methode}

\begin{exoresolu}[Craie et acide : système fermé contre système ouvert]
\textbf{Énoncé.} On place \SI{10}{g} de craie (\ce{CaCO3}) dans une bouteille contenant un flacon d'acide chlorhydrique en excès. L'ensemble, bouché, pèse $m_1 = \SI{150}{g}$. On mélange : une effervescence se produit.

\textbf{1.} Quelle masse $m_2$ la balance indique-t-elle à la fin de la réaction, la bouteille restant fermée ? Justifier.

\textbf{2.} On ouvre maintenant le bouchon et on laisse le gaz s'échapper. La balance indique alors $m_3 = \SI{145,6}{g}$. Calculer la masse de dioxyde de carbone dégagée.

\textbf{3.} Vérifier que la loi de Lavoisier reste respectée.
\tcblower
\textbf{Solution détaillée.}

\textbf{1.} La bouteille est \textbf{fermée} : aucune matière ne peut sortir, pas même le gaz \ce{CO2} formé. D'après la loi de Lavoisier, la masse totale se conserve :
\[ m_2 = m_1 = \SI{150}{g} \]
L'effervescence ne change rien à la pesée, car le gaz produit reste prisonnier dans la bouteille.

\textbf{2.} Une fois le bouchon ouvert, le dioxyde de carbone s'échappe dans l'air. La masse « perdue » par la balance est exactement la masse du gaz parti :
\[ m_{\ce{CO2}} = m_2 - m_3 = 150 - 145{,}6 = \SI{4,4}{g} \]

\textbf{3.} Vérifions la conservation en comptant \emph{toute} la matière :
\begin{align*}
m_{\text{total après}} &= m_{\text{reste dans la bouteille}} + m_{\text{gaz échappé}} \\
&= 145{,}6 + 4{,}4 \\
&= \SI{150}{g} = m_1
\end{align*}
La masse totale est bien conservée : la loi de Lavoisier est respectée, à condition de compter le gaz. La diminution affichée par la balance n'est qu'une \textbf{disparition apparente}.
\end{exoresolu}

\begin{attention}[Erreur fréquente : « le gaz ne pèse rien »]
Beaucoup d'élèves croient que la masse diminue parce que « les gaz ne pèsent rien ». C'est \textbf{faux} ! Tous les gaz ont une masse : un litre d'air pèse environ \SI{1,2}{g}, et la bonbonne de gaz domestique se vend justement \emph{au kilogramme} de butane ! Si la balance indique moins après la réaction, ce n'est pas parce que le gaz est « sans poids », mais parce qu'il a \textbf{quitté le système pesé}. Au BEPC, écrire « la masse diminue car le gaz ne pèse rien » fait perdre des points ; écris plutôt : « la masse diminue car le gaz produit s'échappe du système ouvert ».
\end{attention}

\begin{exercice}[À toi de jouer]
On chauffe \SI{2,4}{g} de magnésium dans un creuset. Le magnésium brûle avec le dioxygène de l'air et il se forme \SI{4,0}{g} d'oxyde de magnésium \ce{MgO} (poudre blanche).
\begin{enumerate}
\item Le système est-il ouvert ou fermé ? Pourquoi la masse du solide a-t-elle augmenté ?
\item Calculer la masse de dioxygène qui a réagi.
\item Citer la loi utilisée.
\end{enumerate}
\end{exercice}

\begin{corrige}[Exercice : combustion du magnésium]
\textbf{1.} Le creuset est \textbf{ouvert} sur l'air : le dioxygène de l'air peut entrer et réagir avec le magnésium. La masse du solide augmente car le solide final contient, en plus du magnésium, l'oxygène capté dans l'air.

\textbf{2.} D'après la loi de Lavoisier appliquée à l'ensemble \{magnésium + dioxygène\} :
\begin{align*}
m_{\ce{Mg}} + m_{\ce{O2}} &= m_{\ce{MgO}} \\
m_{\ce{O2}} &= 4{,}0 - 2{,}4 = \SI{1,6}{g}
\end{align*}

\textbf{3.} On a utilisé la \textbf{loi de Lavoisier} (loi de conservation de la masse) : la somme des masses des réactifs est égale à la somme des masses des produits.
\end{corrige}

\begin{aretenir}[L'essentiel de la section]
\begin{itemize}
\item \textbf{Loi de Lavoisier} : dans une transformation chimique en \textbf{système fermé}, la masse totale se conserve : $\sum m_{\text{réactifs}} = \sum m_{\text{produits}}$.
\item En \textbf{système ouvert}, la masse mesurée peut varier à cause des gaz échangés avec l'air — mais la loi reste vraie si l'on compte toute la matière.
\item Explication microscopique : les \textbf{atomes se conservent} en nombre et en nature ; ils ne font que se réarranger.
\item Les gaz \textbf{ont une masse} : une disparition apparente de matière cache toujours un départ (ou une arrivée) de gaz.
\item Devise à retenir : \emph{« Rien ne se perd, rien ne se crée, tout se transforme. »}
\end{itemize}
\end{aretenir}

Cette conservation des atomes a une conséquence pratique immédiate : pour écrire correctement la transformation que nous venons d'étudier, il faudra que l'équation-bilan respecte, atome par atome, cette loi de conservation. C'est précisément l'objet de la section suivante : l'\textbf{équation-bilan et son équilibrage}.

\coursec{Transformation chimique : mise en évidence expérimentale}

\courssub{Comment reconnaître une transformation chimique ?}

Dans la vie de tous les jours, la matière change sans cesse. Quand tu fais cuire un œuf, que tu fais brûler du bois dans un foyer \emph{trois pierres}, que tu verses du vinaigre sur de la craie ou que tu laisses une moto rouiller sous la pluie, tu observes des \textbf{transformations chimiques}. Contrairement à une transformation physique (fusion de la glace, dissolution du sucre dans l'eau), une transformation chimique \textbf{modifie l'identité des substances} : les espèces chimiques initiales (réactifs) disparaissent et de nouvelles espèces (produits) apparaissent.

\begin{experience}[Signes révélateurs d'une transformation chimique]
\textbf{Matériel :} Tube à essai, bec Bunsen, pince à creuset, poudre de cuivre (ou limaille de fer), acide chlorhydrique dilué, solution de soude, thermomètre, bois (allumette).
\textbf{Protocole et observations :}
\begin{enumerate}
    \item \textbf{Changement de couleur} : Chauffer fortement de la poudre de cuivre (\ce{Cu}, rouge-orangé) dans un tube à essai. Elle noircit : formation d'oxyde de cuivre (\ce{CuO}, noir).
    \item \textbf{Dégagement de gaz (effervescence)} : Verser de l'acide chlorhydrique (\ce{HCl}) sur de la craie (\ce{CaCO3}). Des bulles se forment : c'est du dioxyde de carbone (\ce{CO2}).
    \item \textbf{Variation de température} : Mélanger de l'acide chlorhydrique et de la soude (\ce{NaOH}) dans un bécher muni d'un thermomètre. La température monte : réaction \textbf{exothermique}. (L'inverse, endothermique, absorbe la chaleur, ex. : dissolution du nitrate d'ammonium).
    \item \textbf{Formation d'un précipité} : Mélanger une solution de nitrate de plomb \ce{Pb(NO3)2} (incolore) et une solution d'iodure de potassium \ce{KI} (incolore). Apparition immédiate d'un solide jaune (iodure de plomb \ce{PbI2}).
    \item \textbf{Apparition de lumière / flamme} : Brûler du magnésium ou du gaz de cuisine. Émission de lumière intense.
\end{enumerate}
\textbf{Interprétation :} Si au moins un de ces signes est observé, une transformation chimique a eu lieu. Les atomes ont réarrangé leurs liaisons.
\legende{Les cinq signes macroscopiques d'une transformation chimique. Au laboratoire, on en note au moins un pour l'affirmer.}
\end{experience}

\begin{definition}[Transformation chimique (réaction chimique)]
C'est une transformation au cours de laquelle \textbf{au moins une espèce chimique disparaît} (réactif) et \textbf{au moins une nouvelle espèce chimique apparaît} (produit). Les atomes se réorganisent : les liaisons se rompent et de nouvelles liaisons se forment.
\end{definition}

\begin{attention}[Ne pas confondre physique et chimique !]
\begin{itemize}
    \item \textbf{Physique} : L'espèce chimique ne change pas (ex. : \ce{H2O} liquide $\rightarrow$ \ce{H2O} gazeux). C'est réversible facilement.
    \item \textbf{Chimique} : Les espèces changent d'identité (ex. : \ce{2 H2 + O2 -> 2 H2O}). Souvent irréversible spontanément.
\end{itemize}
\textbf{Astuce BEPC :} Cherche la \textbf{formation d'une nouvelle substance} (gaz, précipité, couleur, chaleur, lumière).
\end{attention}

\coursec{Conservation de la masse et loi de Lavoisier}

\courssub{L'expérience fondatrice : le creuset fermé}

\begin{savaistu}[Antoine Lavoisier (1743--1794), le père de la chimie moderne]
Avant lui, on croyait au « phlogiston » (principe du feu qui s'échappe). Lavoisier, avec sa balance de précision (au mg près), a prouvé que \textbf{la masse ne se perd pas}. Il a nommé l'oxygène (\emph{acide générateur}) et l'hydrogène (\emph{eau génératrice}). Guillotiné pendant la Terreur : « La République n'a pas besoin de savants » (Lagrange).
\end{savaistu}

\begin{experience}[Conservation de la masse (Lavoisier) -- Démonstration]
\textbf{Matériel :} Balance électronique de précision (\SI{0,01}{g}), creuset avec couvercle, acide chlorhydrique dilué, marbre (calcaire \ce{CaCO3}), petit flacon en verre fin (ou tube à essai) pouvant entrer dans le creuset.
\textbf{Protocole :}
\begin{enumerate}
    \item Placer le marbre dans le creuset. Placer l'acide dans le petit flacon. Fermer le creuset avec son couvercle. Peser l'ensemble : $m_{\text{initial}}$.
    \item Basculer le creuset pour verser l'acide sur le marbre \textbf{sans ouvrir le couvercle}. La réaction se produit : effervescence (\ce{CO2} gazeux reste piégé).
    \item Laisser refroidir à température ambiante. Peser à nouveau : $m_{\text{final}}$.
\end{enumerate}
\textbf{Observation :} $m_{\text{initial}} = m_{\text{final}}$ (à la précision de la balance près).
\textbf{Interprétation :} Le gaz \ce{CO2} produit est resté dans le système fermé. La masse totale est conservée.
\legende{Montage de Lavoisier : le système fermé (creuset + couvercle) garantit qu'aucune matière ne s'échange avec l'extérieur.}
\end{experience}

\begin{experience}[Système ouvert : l'illusion de la perte de masse]
Recommencer l'expérience \textbf{sans le couvercle}. On observe $m_{\text{final}} < m_{\text{initial}}$. Le gaz \ce{CO2} s'est échappé dans l'air de la classe. La masse de \textbf{matériau dans le creuset} a diminué, mais la masse \textbf{totale de l'univers} (creuset + air ambiant) est conservée.
\end{experience}

\begin{propriete}[Loi de Lavoisier (Loi de conservation de la masse)]
\textbf{Au cours d'une transformation chimique réalisée dans un système fermé, la masse totale du système ne change pas.}
\[
m_{\text{réactifs}} = m_{\text{produits}}
\]
\textbf{Au niveau microscopique :} Les atomes ne sont ni créés, ni détruits, ils changent seulement de partenaires. Le nombre d'atomes de chaque élément est conservé.
\end{propriete}

\begin{aretenir}[Système ouvert vs fermé]
\begin{itemize}
    \item \textbf{Système fermé} : Aucun échange de matière avec l'extérieur (ex. : creuset fermé, seringue bouchée, Terre \emph{quasi} fermée). \textbf{Masse conservée.}
    \item \textbf{Système ouvert} : Échange de matière possible (ex. : creuset ouvert, casserole, moteur de moto). La masse \emph{apparente} peut varier si un gaz entre ou sort.
\end{itemize}
\end{aretenir}

\coursec{Équation-bilan et équilibrage}

\courssub{De la conservation de la masse à l'écriture symbolique}

Grâce à la loi de Lavoisier, nous savons que les atomes se conservent. Pour décrire ce réarrangement de manière précise, universelle et quantitative, les chimistes utilisent \textbf{l'équation-bilan}. C'est le « passeport » de la réaction chimique.

\begin{definition}[Équation-bilan d'une réaction chimique]
L'\textbf{équation-bilan} est l'écriture symbolique d'une transformation chimique.
\begin{itemize}
    \item À \textbf{gauche} de la flèche ($\rightarrow$) : les \textbf{réactifs} (espèces chimiques qui disparaissent ou diminuent).
    \item À \textbf{droite} de la flèche : les \textbf{produits} (espèces chimiques qui apparaissent ou augmentent).
    \item La flèche signifie « donne » ou « se transforme en ».
    \item Les formules chimiques (\ce{H2O}, \ce{CO2}, \ce{CH4}, etc.) ne doivent \textbf{jamais} être modifiées : elles représentent l'identité immuable des substances.
    \item On note souvent l'état physique : (s) solide, (l) liquide, (g) gaz, (aq) aqueux (dissous dans l'eau).
\end{itemize}
\end{definition}

\begin{definition}[Coefficient stœchiométrique]
Le nombre entier placé \textbf{devant} une formule chimique (ex. : le \textbf{2} dans $2\,\ce{H2O}$) s'appelle le \textbf{coefficient stœchiométrique}. Il indique le \textbf{rapport entre les nombres d'entités} (atomes, molécules, ions) qui réagissent ou se forment. S'il n'y a pas de nombre écrit, le coefficient vaut \textbf{1}.
\end{definition}

\courssub{Exemple fondateur : la combustion du carbone}

Reprenons l'expérience classique : on fait brûler du charbon (carbone \ce{C}) dans du dioxygène (\ce{O2}). On observe la formation de dioxyde de carbone (\ce{CO2}).

\begin{enumerate}
    \item \textbf{Identification des espèces} : Réactifs = \ce{C} (solide) et \ce{O2} (gaz) ; Produit = \ce{CO2} (gaz).
    \item \textbf{Écriture provisoire} : $\ce{C + O2 -> CO2}$.
    \item \textbf{Vérification atome par atome} :
    \begin{itemize}
        \item Carbone : 1 à gauche, 1 à droite $\rightarrow$ \textbf{OK}.
        \item Oxygène : 2 à gauche (dans \ce{O2}), 2 à droite (dans \ce{CO2}) $\rightarrow$ \textbf{OK}.
    \end{itemize}
    L'équation est \textbf{déjà équilibrée} : $\ce{C + O2 -> CO2}$.
\end{enumerate}

\begin{exemplebox}[Lecture quantitative de l'équation $\ce{C + O2 -> CO2}$]
Cette équation nous dit trois choses équivalentes :
\begin{itemize}
    \item \textbf{Niveau microscopique} : \textbf{1 atome} de carbone réagit avec \textbf{1 molécule} de dioxygène pour former \textbf{1 molécule} de dioxyde de carbone.
    \item \textbf{Niveau macroscopique (mole)} : \textbf{1 mole} de carbone (\SI{12,0}{g}) réagit avec \textbf{1 mole} de dioxygène (\SI{32,0}{g}) pour former \textbf{1 mole} de dioxyde de carbone (\SI{44,0}{g}). \emph{(Notion de mole approfondie en Seconde)}.
    \item \textbf{Niveau masse} : \SI{12,0}{g} de carbone + \SI{32,0}{g} de dioxygène $\rightarrow$ \SI{44,0}{g} de dioxyde de carbone. La masse est conservée (\SI{44,0}{g} = \SI{44,0}{g}).
\end{itemize}
\end{exemplebox}

\courssub{La règle d'or : on ne touche qu'aux coefficients}

C'est le point le plus important, et la source de l'erreur n°1 au BEPC.

\begin{attention}[Interdiction formelle : ne jamais modifier les indices !]
Dans une formule chimique, les \textbf{indices} (les petits chiffres en bas à droite, comme le 2 dans \ce{H2O} ou \ce{O2}) définissent \textbf{l'identité de la substance}.
\begin{itemize}
    \item \ce{H2O} = eau. \ce{H2O2} = eau oxygénée (peroxyde d'hydrogène). Ce sont deux substances totalement différentes !
    \item \ce{CO2} = dioxyde de carbone. \ce{CO} = monoxyde de carbone (toxique mortel).
\end{itemize}
\textbf{Équilibrer une équation, c'est ajouter des coefficients DEVANT les formules, jamais modifier les indices DANS les formules.}
\end{attention}

\courssub{Méthode pas à pas : l'équilibrage de la combustion du méthane}

Le méthane (\ce{CH4}) est le composant principal du gaz naturel et du gaz de bouteille (butane/propane) utilisé dans nos cuisines au Cameroun. Sa combustion complète donne du dioxyde de carbone (\ce{CO2}) et de l'eau (\ce{H2O}). Appliquons la méthode rigoureuse en 4 étapes.

\begin{methode}[Équilibrer une équation-bilan en 4 étapes]
\begin{enumerate}
    \item \textbf{Écrire les formules} correctes des réactifs (gauche) et des produits (droite). Ne rien inventer, utiliser les connaissances du cours (nomenclature, ions, molécules courantes).
    \item \textbf{Compter les atomes} de chaque élément présent, des deux côtés de la flèche. Faire un petit tableau de comptage (brouillon).
    \item \textbf{Ajuster les coefficients} un élément après l'autre.
        \begin{itemize}
            \item Commencer par l'élément le plus « complexe » (présent dans une seule espèce de chaque côté, souvent C ou métaux).
            \item Continuer par l'hydrogène (H).
            \item Finir par l'oxygène (O) qui est souvent présent partout.
        \end{itemize}
    \item \textbf{Vérifier} le bilan final atome par atome. Tous les coefficients doivent être des \textbf{entiers} et \textbf{premiers entre eux} (simplifiés au maximum).
\end{enumerate}
\end{methode}

\begin{exoresolu}[Équilibrage complet de la combustion du méthane]
\textbf{Étape 1 : Écriture des formules.}
Réactifs : Méthane \ce{CH4} + Dioxygène \ce{O2}.
Produits : Dioxyde de carbone \ce{CO2} + Eau \ce{H2O}.
Équation provisoire : $\ce{CH4 + O2 -> CO2 + H2O}$

\textbf{Étape 2 : Comptage initial (coefficients = 1 partout).}

\begin{popfigure}
\centering
\begin{tikzpicture}[
    cell/.style={minimum width=2.2cm, minimum height=0.9cm, draw=popdark, thick},
    header/.style={cell, fill=popblue, text=white, font=\bfseries},
    ok/.style={cell, fill=popgreenL},
    nok/.style={cell, fill=poppinkL},
    fix/.style={cell, fill=popgoldL},
    elem/.style={minimum width=2.5cm, minimum height=0.9cm, draw=popdark, thick, fill=popblueL, font=\bfseries, align=center}
]
% En-têtes
\node[header] (e) at (0,0) {Élément};
\node[header] (rg) at (2.2,0) {Réactifs (Gauche)};
\node[header] (pd) at (4.4,0) {Produits (Droite)};
\node[header] (eq) at (6.6,0) {Équilibré ?};

% Carbone
\node[elem] (c) at (0,-0.9) {Carbone (C)};
\node[nok] (cr) at (2.2,-0.9) {1 (dans \ce{CH4})};
\node[nok] (cp) at (4.4,-0.9) {1 (dans \ce{CO2})};
\node[ok] (ceq) at (6.6,-0.9) {\textcolor{popgreen}{\textbf{✓}} \quad OK};

% Hydrogène
\node[elem] (h) at (0,-1.8) {Hydrogène (H)};
\node[nok] (hr) at (2.2,-1.8) {4 (dans \ce{CH4})};
\node[nok] (hp) at (4.4,-1.8) {2 (dans \ce{H2O})};
\node[nok] (heq) at (6.6,-1.8) {\textbf{Non} (4 vs 2)};

% Oxygène
\node[elem] (o) at (0,-2.7) {Oxygène (O)};
\node[nok] (or) at (2.2,-2.7) {2 (dans \ce{O2})};
\node[nok] (op) at (4.4,-2.7) {2+1=3 (\ce{CO2}+\ce{H2O})};
\node[nok] (oeq) at (6.6,-2.7) {\textbf{Non} (2 vs 3)};

% Flèches d'action
\draw[poporange, very thick, ->] (6.6,-1.8) -- ++(1.5,0) node[right, align=left, font=\small] {Mettre \textbf{2} devant \ce{H2O}};
\draw[poporange, very thick, ->] (6.6,-2.7) -- ++(1.5,0) node[right, align=left, font=\small] {Mettre \textbf{2} devant \ce{O2}};
\end{tikzpicture}
\legende{Tableau de comptage atome par atome avant équilibrage. Le carbone est déjà bon. L'hydrogène et l'oxygène ne le sont pas.}
\end{popfigure}

\textbf{Étape 3 : Ajustement des coefficients.}
\begin{enumerate}
    \item \textbf{Carbone (C)} : 1 à gauche, 1 à droite $\rightarrow$ \textbf{OK}. On ne touche pas.
    \item \textbf{Hydrogène (H)} : 4 à gauche (\ce{CH4}), 2 à droite (\ce{H2O}). Pour avoir 4 à droite, je mets \textbf{2} devant \ce{H2O} : $\ce{CH4 + O2 -> CO2 + 2 H2O}$.
    \item \textbf{Oxygène (O)} : Maintenant à droite : 2 (dans \ce{CO2}) + 2$\times$1 (dans 2 \ce{H2O}) = \textbf{4 atomes O}. À gauche : 2 (dans \ce{O2}). Pour avoir 4 à gauche, je mets \textbf{2} devant \ce{O2} : $\ce{CH4 + 2 O2 -> CO2 + 2 H2O}$.
\end{enumerate}

\textbf{Étape 4 : Vérification finale.}

\begin{popfigure}
\centering
\begin{tikzpicture}[
    cell/.style={minimum width=2.2cm, minimum height=0.9cm, draw=popdark, thick},
    header/.style={cell, fill=popblue, text=white, font=\bfseries},
    ok/.style={cell, fill=popgreenL},
    elem/.style={minimum width=2.5cm, minimum height=0.9cm, draw=popdark, thick, fill=popblueL, font=\bfseries, align=center}
]
\node[header] (e) at (0,0) {Élément};
\node[header] (rg) at (2.2,0) {Réactifs (Gauche)};
\node[header] (pd) at (4.4,0) {Produits (Droite)};
\node[header] (eq) at (6.6,0) {Bilan};

\node[elem] (c) at (0,-0.9) {Carbone (C)};
\node[ok] (cr) at (2.2,-0.9) {1};
\node[ok] (cp) at (4.4,-0.9) {1};
\node[ok] (ceq) at (6.6,-0.9) {1 = 1 \textcolor{popgreen}{\textbf{✓}}};

\node[elem] (h) at (0,-1.8) {Hydrogène (H)};
\node[ok] (hr) at (2.2,-1.8) {4};
\node[ok] (hp) at (4.4,-1.8) {2$\times$2 = 4};
\node[ok] (heq) at (6.6,-1.8) {4 = 4 \textcolor{popgreen}{\textbf{✓}}};

\node[elem] (o) at (0,-2.7) {Oxygène (O)};
\node[ok] (or) at (2.2,-2.7) {2$\times$2 = 4};
\node[ok] (op) at (4.4,-2.7) {2 + 2$\times$1 = 4};
\node[ok] (oeq) at (6.6,-2.7) {4 = 4 \textcolor{popgreen}{\textbf{✓}}};
\end{tikzpicture}
\legende{Tableau de vérification final : tous les atomes sont conservés. L'équation est équilibrée.}
\end{popfigure}

\textbf{Équation-bilan finale :}
\[
\boxed{\ce{CH4 + 2 O2 -> CO2 + 2 H2O}}
\]

\textbf{Lecture quantitative :} \textbf{1 molécule} de méthane réagit avec \textbf{2 molécules} de dioxygène pour donner \textbf{1 molécule} de dioxyde de carbone et \textbf{2 molécules} d'eau.
\end{exoresolu}

\courssub{Représentation microscopique : visualiser les coefficients}

Pour bien comprendre ce que signifient les coefficients, rien ne vaut un schéma « à boules » montrant les atomes réarrangés.

\begin{popfigure}
\centering
\begin{tikzpicture}[
    atomC/.style={circle, draw=popdark, thick, fill=popblue!80, minimum size=0.7cm, font=\bfseries\small\color{white}},
    atomH/.style={circle, draw=popdark, thick, fill=poporange!80, minimum size=0.5cm, font=\bfseries\small\color{white}},
    atomO/.style={circle, draw=popdark, thick, fill=popgreen!80, minimum size=0.6cm, font=\bfseries\small\color{white}},
    coef/.style={font=\Large\bfseries, text=popdark},
    plus/.style={font=\Large\bfseries, text=popdark},
    arrow/.style={->, very thick, popdark, shorten >=2pt, shorten <=2pt}
]
% Réactifs : CH4
\node[coef] at (-5.5, 1.5) {\textbf{1}};
\node[atomC] (C) at (-5.5, 1.5) {C};
\foreach \ang in {90, 210, 330, -30} {
    \node[atomH] at ($(C) + (\ang:1.1)$) {H};
}
\node[plus] at (-4.0, 1.5) {+};

% Réactifs : 2 O2
\node[coef] at (-2.5, 1.5) {\textbf{2}};
\foreach \j in {0, 1.3} {
    \node[atomO] (O1) at (-2.5+\j, 2.2) {O};
    \node[atomO] at (-2.5+\j, 1.0) {O};
}

% Flèche
\draw[arrow] (-0.8, 1.5) -- (0.8, 1.5) node[midway, above, font=\small] {Chaleur / Flamme};

% Produits : CO2
\node[coef] at (2.8, 2.2) {\textbf{1}};
\node[atomC] (C2) at (2.8, 2.2) {C};
\node[atomO] at (2.0, 2.2) {O};
\node[atomO] at (3.6, 2.2) {O};

% Produits : 2 H2O
\node[coef] at (2.8, 0.3) {\textbf{2}};
\foreach \k in {0, 1.3} {
    \node[atomO] (Ow) at (2.8+\k, 0.8) {O};
    \node[atomH] at (2.8+\k, 1.5) {H};
    \node[atomH] at (2.8+\k, 0.1) {H};
}

% Légende atomes
\node[align=center, font=\small] at (-7.5, 1.5) {\textbf{Réactifs}};
\node[align=center, font=\small] at (5.5, 1.5) {\textbf{Produits}};
\begin{scope}[yshift=-2.5cm]
\node[atomC, label=right:\textbf{Carbone (C)}] at (0,0) {};
\node[atomH, label=right:\textbf{Hydrogène (H)}] at (0,-0.9) {};
\node[atomO, label=right:\textbf{Oxygène (O)}] at (0,-1.8) {};
\end{scope}
\end{tikzpicture}
\legende{Modélisation microscopique de $\ce{CH4 + 2 O2 -> CO2 + 2 H2O}$. On compte bien 1 C, 4 H, 4 O de chaque côté. Les atomes sont les mêmes, seules les liaisons changent.}
\end{popfigure}

\courssub{Autres exemples classiques au programme de 3ème}

Tu dois savoir équilibrer ces réactions types rencontrées au laboratoire ou dans la vie courante.

\begin{exemplebox}[Combustion du magnésium (éclatante !)]
Expérience : On tient un ruban de magnésium (\ce{Mg}) avec une pince à creuset et on l'approche de la flamme du bec Bunsen. \textbf{Attention : ne pas regarder la flamme directement (risque pour les yeux, lumière UV intense)}. Il se forme une poudre blanche : l'oxyde de magnésium (\ce{MgO}).
\begin{enumerate}
    \item Formules : \ce{Mg} + \ce{O2} $\rightarrow$ \ce{MgO}
    \item Comptage : Mg (1 vs 1) OK. O (2 vs 1) \textbf{Pas OK}.
    \item Ajustement : Mettre \textbf{2} devant \ce{MgO} $\rightarrow$ \ce{Mg} + \ce{O2} $\rightarrow$ \textbf{2} \ce{MgO}.
    \item Nouveau comptage Mg : 1 vs 2. Mettre \textbf{2} devant \ce{Mg}.
    \item \textbf{Équation finale :} $\boxed{\ce{2 Mg + O2 -> 2 MgO}}$
\end{enumerate}
\end{exemplebox}

\begin{exemplebox}[Action de l'acide chlorhydrique sur la craie (calcaire)]
C'est la réaction qui fait « mousser » le calcaire (roche, coquille d'œuf, craie scolaire) quand on verse de l'acide (détartrant, vinaigre, acide chlorhydrique dilué). Gaz dégagé : \ce{CO2} (test à l'eau de chaux trouble).
\begin{enumerate}
    \item Espèces : Craie = \ce{CaCO3} (carbonate de calcium). Acide = \ce{HCl}. Produits : Chlorure de calcium \ce{CaCl2} (soluble), Eau \ce{H2O}, Dioxyde de carbone \ce{CO2} (gaz).
    \item Provisoire : $\ce{CaCO3 + HCl -> CaCl2 + CO2 + H2O}$
    \item Comptage : Ca (1=1), C (1=1), Cl (1 vs 2) $\rightarrow$ mettre \textbf{2} devant \ce{HCl}.
    \item Vérif H : 2 (dans 2 \ce{HCl}) vs 2 (dans \ce{H2O}) OK. O : 3 (dans \ce{CaCO3}) vs 2+1=3 OK.
    \item \textbf{Équation finale :} $\boxed{\ce{CaCO3 + 2 HCl -> CaCl2 + CO2 + H2O}}$
\end{enumerate}
\end{exemplebox}

\courssub{Complément : Réactif limitant et réactif en excès (ouverture Seconde)}

\emph{Ce paragraphe n'est pas exigible au BEPC 3ème, mais il éclaire la compréhension des quantités et prépare la Seconde.}

Dans la réalité (cuisine, industrie, moteur), on ne mélange pas toujours les réactifs dans les \textbf{proportions exactes} de l'équation équilibrée (proportions stœchiométriques).

\begin{definition}[Réactif limitant et réactif en excès]
\begin{itemize}
    \item Le \textbf{réactif limitant} est celui qui sera \textbf{entièrement consommé} le premier. Il \textbf{limite} la quantité de produits formés.
    \item Le(s) \textbf{réactif(s) en excès} reste(nt) en partie \textbf{non réagis} à la fin de la transformation.
\end{itemize}
\end{definition}

\begin{exemplebox}[L'acide en excès sur la craie (détartrage)]
Pour détartrer une bouilloire, on verse beaucoup de vinaigre (acide acétique) ou d'acide chlorhydrique dilué sur le calcaire (\ce{CaCO3}).
\begin{itemize}
    \item L'acide est mis \textbf{volontairement en excès} (beaucoup plus que les 2 moles nécessaires par mole de \ce{CaCO3}).
    \item Le \textbf{calcaire est le réactif limitant} : tout le calcaire disparaît.
    \item À la fin, il reste de l'acide non réagi dans la solution (d'où la nécessité de rincer abondamment !).
    \item Avantage : la réaction est rapide et complète pour le calcaire.
\end{itemize}
\end{exemplebox}

\begin{savaistu}[Le gaz de cuisine au Cameroun : Butane et Propane]
Dans nos bouteilles bleues (butane) ou grises (propane) de la SONARA / Tradex / Gaz du Cameroun, l'hydrocarbure principal est le \textbf{butane} (\ce{C4H10}) ou le \textbf{propane} (\ce{C3H8}). Leur combustion complète s'écrit :
\begin{align*}
\textbf{Butane :} \quad & \ce{2 C4H10 + 13 O2 -> 8 CO2 + 10 H2O} \\
\textbf{Propane :} \quad & \ce{C3H8 + 5 O2 -> 3 CO2 + 4 H2O}
\end{align*}
Regarde les coefficients : pour 2 molécules de butane, il faut \textbf{13 molécules de dioxygène} ! C'est énorme. C'est pour cela qu'une cuisine mal aérée devient vite dangereuse : l'oxygène de l'air (\approx 21\%) manque, la combustion devient \textbf{incomplète} et produit du \textbf{monoxyde de carbone (\ce{CO})}, un gaz inodore, incolore et mortel. \textbf{Aère toujours ta cuisine quand tu cuisines au gaz !}
\end{savaistu}

\coursec{Combustions complète et incomplète}

\courssub{La combustion : réaction d'oxydation rapide}

Une \textbf{combustion} est une transformation chimique rapide entre un \textbf{combustible} (souvent un hydrocarbure : carbone + hydrogène) et un \textbf{comburant} (le plus souvent le dioxygène \ce{O2} de l'air), qui dégage de la \textbf{chaleur} et de la \textbf{lumière} (flamme). C'est la base de nos énergies : cuisine (gaz), transport (moto-taxi, voiture, avion), électricité (groupe électrogène, centrale thermique).

\begin{definition}[Combustion complète]
Lorsque le comburant (dioxygène) est en \textbf{quantité suffisante}, le carbone et l'hydrogène du combustible s'oxydent au maximum :
\begin{itemize}
    \item Carbone $\rightarrow$ Dioxyde de carbone \ce{CO2} (gaz incolore, non toxique à faible dose, mais gaz à effet de serre).
    \item Hydrogène $\rightarrow$ Eau \ce{H2O} (vapeur d'eau).
\end{itemize}
\textbf{Flamme bleue, non lumineuse, très chaude} (ex. : bec Bunsen air ouvert, gazinière bien réglée).
\end{definition}

\begin{definition}[Combustion incomplète]
Lorsque le comburant \textbf{manque} (espace confiné, brûleur encrassé, moteur mal réglé), l'oxydation est partielle :
\begin{itemize}
    \item Carbone $\rightarrow$ \textbf{Monoxyde de carbone \ce{CO}} (gaz \textbf{toxique mortel}, incolore, inodore) \textbf{ET/OU} Carbone \ce{C} (suie, particules fines, noir de fumée).
    \item Hydrogène $\rightarrow$ Eau \ce{H2O}.
\end{itemize}
\textbf{Flamme jaune/orange, lumineuse, moins chaude, dégage de la suie} (ex. : bec Bunsen air fermé, flamme de bougie, pot d'échappement noir).
\end{definition}

\begin{exemplebox}[Combustions du méthane : les deux cas]
\begin{itemize}
    \item \textbf{Complète (air suffisant) :} $\ce{CH4 + 2 O2 -> CO2 + 2 H2O}$ \quad (Flamme bleue, rendement max).
    \item \textbf{Incomplète (air manquant) :} Deux possibilités selon le manque d'O2 :
    \begin{align*}
        \text{Modéré :} \quad & \ce{2 CH4 + 3 O2 -> 2 CO + 4 H2O} \quad (\ce{CO} toxique) \\
        \text{Sévère :} \quad & \ce{CH4 + O2 -> C + 2 H2O} \quad (\text{Suie } \ce{C} \text{ noire})
    \end{align*}
\end{itemize}
\end{exemplebox}

\begin{attention}[Piège BEPC : Équilibrer l'incomplète]
Ne jamais écrire $\ce{CH4 + O2 -> CO + H2O}$ (non équilibré en H et O).
Toujours faire le comptage atome par atome. Pour $\ce{2 CH4 + 3 O2 -> 2 CO + 4 H2O}$ :
C: 2=2, H: 8=8, O: 6 = 2+4=6. \textbf{OK}.
\end{attention}

\coursec{Le monoxyde de carbone : un danger invisible}

\courssub{Pourquoi le CO est-il si dangereux ?}

Le monoxyde de carbone (\ce{CO}) est un gaz \textbf{incolore}, \textbf{inodore}, \textbf{inflammable} et \textbf{toxique}. Il ne se détecte pas par les sens humains. C'est la première cause de mortalité par intoxication accidentelle dans le monde (et au Cameroun : groupes électrogènes dans les chambres, foyers \emph{trois pierres} dans les cases fermées, gazinières défectueuses, pots d'échappement dans les garages).

\begin{propriete}[Mécanisme toxique : l'hémoglobine piégée]
Dans le sang, l'\textbf{hémoglobine} (\ce{Hb}) transporte l'oxygène (\ce{O2}) vers les cellules.
\begin{itemize}
    \item Normalement : $\ce{Hb + O2 <=> HbO2}$ (oxyhémoglobine, sang rouge vif).
    \item Avec le CO : $\ce{Hb + CO -> HbCO}$ (carboxyhémoglobine, sang rouge cerise).
\end{itemize}
L'affinité de l'hémoglobine pour le CO est \textbf{200 à 250 fois plus grande} que pour l'\ce{O2}. Le CO « vole » la place de l'\ce{O2} sur l'hémoglobine. Les cellules s'asphyxient (hypoxie) alors que le sang est rouge.
\end{propriete}

\begin{experience}[Détection du monoxyde de carbone (Démonstration professeur)]
\textbf{Matériel :} Tube détecteur \ce{CO} (type Dräger/Gastec) avec pompe manuelle, source de CO contrôlée (ex. : combustion incomplète dans hotte aspirante, ou ampoule calibrée).
\textbf{Protocole :} Casser les embouts du tube, l'insérer dans la pompe, aspirer le volume d'air indiqué (ex. : 100 mL).
\textbf{Observation :} Le réactif dans le tube (souvent à base de palladium ou iode) change de couleur (blanc $\rightarrow$ brun/vert/noir) proportionnellement à la concentration en \ce{CO}.
\textbf{Interprétation :} Lecture directe de la concentration en \SI{}{ppm} (parties par million) sur l'échelle graduée du tube. Seuil danger : > \SI{50}{ppm} (exposition prolongée), > \SI{200}{ppm} (maux de tête), > \SI{800}{ppm} (mortel en quelques heures).
\legende{Tube détecteur de CO : outil indispensable pour les sapeurs-pompiers et la maintenance des chaudières/générateurs.}
\end{experience}

\begin{aretenir}[Symptômes d'intoxication au CO (ordre de gravité)]
\begin{enumerate}
    \item Maux de tête, vertiges, nausées, fatigue inhabituelle (souvent confondus avec grippe/paludisme).
    \item Troubles de la vision, confusion, essoufflement, douleurs thoraciques.
    \item Perte de connaissance, coma, arrêt cardio-respiratoire, \textbf{mort}.
    \item \textbf{Séquelles neurologiques} possibles chez les survivants (troubles mémoire, Parkinsonien).
\end{enumerate}
\textbf{Particulièrement vulnérables :} Enfants, femmes enceintes (risque fœtus), personnes âgées, cardiaques, anémiques.
\end{aretenir}

\coursec{Sécurité et gestes responsables autour des combustions}

\courssub{Prévenir le risque : la ventilation avant tout}

\begin{methode}[Les 5 règles d'or au Cameroun (et partout)]
\begin{enumerate}
    \item \textbf{Aérer en permanence :} Ouvrir les fenêtres/portes quand on cuisine au gaz/bois/charbon. Ne \textbf{jamais} boucher les grilles d'aération hautes et basses (obligatoires dans les cuisines). Un groupe électrogène \textbf{ne fonctionne JAMAIS} dans une pièce fermée, un couloir, un garage attenant à une chambre. \textbf{TOUJOURS À L'EXTÉRIEUR}, loin des ouvertures.
    \item \textbf{Entretenir les appareils :} Vérifier les tuyaux de gaz (date de péremption, fissures), les brûleurs (flamme bleue = bon, flamme jaune = encrassé/manque d'air), le pot d'échappement de la moto/voiture (pas de fuite dans l'habitacle).
    \item \textbf{Installer un détecteur de CO :} Norme EN 50291. À placer dans les pièces à risque (chambre près du garage, cuisine au gaz, salon avec cheminée/poêle). Tester la pile chaque mois.
    \item \textbf{Ne pas utiliser d'appareils non prévus :} Pas de réchaud à gaz « camping » pour chauffer une chambre. Pas de brasero/charbon dans la case pour chauffer. Pas de moteur tournant au ralenti dans un garage fermé.
    \item \textbf{Connaître la conduite à tenir :} Voir ci-dessous.
\end{enumerate}
\end{methode}

\courssub{Conduite à tenir en cas de suspicion d'intoxication au CO}

\begin{attention}[URGENCE VITALE : Réagir vite et bien]
\begin{enumerate}
    \item \textbf{Aérer IMMÉDIATEMENT :} Ouvrir grandes portes et fenêtres.
    \item \textbf{Évacuer les lieux :} Sortir tout le monde (ne pas chercher les affaires). Ne pas rallumer d'interrupteur (étincelle = explosion si fuite gaz).
    \item \textbf{Appeler les secours :} \textbf{112} (numéro d'urgence unique Cameroun) ou \textbf{119} (Sapeurs-Pompiers). Dire : « Intoxication au monoxyde de carbone, adresse, nombre de victimes, état conscient/inconscient ».
    \item \textbf{Ne pas réintégrer} tant qu'un professionnel (pompiers, technicien gaz) n'a pas mesuré l'air et réparé la cause.
\end{enumerate}
\textbf{Règle d'or :} Un détecteur qui sonne = \textbf{On sort, on appelle, on ne rentre pas.}
\end{attention}

\begin{savaistu}[Le triangle du feu et l'extinction]
Pour qu'il y ait combustion (feu), il faut 3 éléments : \textbf{Combustible} + \textbf{Comburant (O2)} + \textbf{Énergie d'activation (Chaleur)}. Retirer \textbf{un seul} des trois éteint le feu.
\begin{itemize}
    \item Retirer combustible : Couper l'arrivée de gaz, éloigner le bois.
    \item Retirer comburant : Étouffer (couvercle sur casserole, couverture anti-feu, extincteur \ce{CO2}/poudre).
    \item Retirer chaleur : Refroidir à grande eau (feu de classe A : bois, papier, tissu). \textbf{JAMAIS d'eau sur feu d'huile/électrique/gaz} (éclaboussures brûlantes, conduction électrique, explosion).
\end{itemize}
\textbf{Extincteur poudre ABC} : polyvalent (solides, liquides, gaz, électricité). À avoir à la maison et dans la voiture. Vérifier la pression (aiguille dans le vert) chaque année.
\end{savaistu}

\courssub{Exercices d'entraînement (niveau BEPC)}

\begin{exercice}[Équilibrage de combustions]
Équilibrer les équations-bilan suivantes (combustions complètes) :
\begin{enumerate}
    \item Hydrogène : $\ce{H2 + O2 -> H2O}$
    \item Monoxyde de carbone : $\ce{CO + O2 -> CO2}$
    \item Éthane (gaz) : $\ce{C2H6 + O2 -> CO2 + H2O}$
    \item Pentane (essence) : $\ce{C5H12 + O2 -> CO2 + H2O}$
\end{enumerate}
\end{exercice}

\begin{exercice}[Identification et équilibrage]
On fait réagir du zinc (\ce{Zn}) avec de l'acide chlorhydrique (\ce{HCl}). Il se forme du chlorure de zinc (\ce{ZnCl2}) et un gaz qui « pop » à la flamme (dihydrogène \ce{H2}).
\begin{enumerate}
    \item Écrire l'équation-bilan provisoire.
    \item L'équilibrer.
    \item Dire quel est le réactif limitant si on met 1 g de zinc dans 100 mL d'acide chlorhydrique \SI{1}{mol/L} (approximatif : \ce{Zn} limitant).
\end{enumerate}
\end{exercice}

\begin{exercice}[Combustion incomplète : danger au garage]
Un groupe électrogène tourne dans un garage fermé. Le carburant est de l'essence (modélisée par l'octane \ce{C8H18}). Faute d'air, la combustion est incomplète et produit du \ce{CO} et de l'eau.
\begin{enumerate}
    \item Écrire l'équation-bilan provisoire de cette combustion incomplète.
    \item L'équilibrer.
    \item Pourquoi est-ce mortel de dormir dans la pièce au-dessus du garage ?
\end{enumerate}
\end{exercice}

\begin{corrige}[Exercice 1]
\begin{enumerate}
    \item $\ce{2 H2 + O2 -> 2 H2O}$ (H : 4=4, O : 2=2)
    \item $\ce{2 CO + O2 -> 2 CO2}$ (C : 2=2, O : 2+2=4 vs 4 OK)
    \item $\ce{2 C2H6 + 7 O2 -> 4 CO2 + 6 H2O}$ (C : 4=4, H : 12=12, O : 14 = 8+6=14)
    \item $\ce{C5H12 + 8 O2 -> 5 CO2 + 6 H2O}$ (C : 5=5, H : 12=12, O : 16 = 10+6=16)
\end{enumerate}
\end{corrige}

\begin{corrige}[Exercice 2]
\begin{enumerate}
    \item $\ce{Zn + HCl -> ZnCl2 + H2}$
    \item Comptage : Zn (1=1), Cl (1 vs 2) $\rightarrow$ \textbf{2} devant \ce{HCl}. H (2 vs 2) OK.
    \item $\boxed{\ce{Zn + 2 HCl -> ZnCl2 + H2}}$
    \item Masse molaire \ce{Zn} = \SI{65,4}{g/mol}. Quantité de zinc = \SI{1}{g} / \SI{65,4}{g/mol} $\approx$ \SI{0,0153}{mol}. Quantité \ce{HCl} = \SI{0,100}{L} $\times$ \SI{1,0}{mol/L} = \SI{0,100}{mol}. Il faut 2 mol \ce{HCl} par mol \ce{Zn} $\rightarrow$ besoin \SI{0,0306}{mol} \ce{HCl}. On a \SI{0,100}{mol} \ce{HCl} $\rightarrow$ \ce{HCl} en excès, \ce{Zn} limitant.
\end{enumerate}
\end{corrige}

\begin{corrige}[Exercice 3]
\begin{enumerate}
    \item $\ce{C8H18 + O2 -> CO + H2O}$
    \item Comptage : C : 8 $\rightarrow$ 8 \ce{CO}. H : 18 $\rightarrow$ 9 \ce{H2O}. O : à droite = 8 (dans 8 \ce{CO}) + 9 (dans 9 \ce{H2O}) = 17 atomes O. À gauche : \ce{O2} $\rightarrow$ coefficient 17/2 = 8,5. Multiplier par 2 pour entiers :
    $\boxed{\ce{2 C8H18 + 17 O2 -> 16 CO + 18 H2O}}$
    \item Le \ce{CO} est un gaz incolore, inodore, plus léger que l'air (densité 0,97). Il traverse les planchers, les cloisons, monte à l'étage. Il se fixe sur l'hémoglobine des dormants $\rightarrow$ asphyxie silencieuse, mort dans le sommeil. \textbf{JAMAIS de groupe électrogène dans un garage attenant à l'habitation.}
\end{enumerate}
\end{corrige}

\begin{aretenir}[Résumé : La réaction chimique, de l'expérience à la sécurité]
\begin{enumerate}
    \item \textbf{Reconnaître} une transformation chimique : changement de couleur, gaz, précipité, température, lumière. Nouvelles substances formées.
    \item \textbf{Loi de Lavoisier} : Dans un système fermé, la masse est conservée. Les atomes se réorganisent.
    \item \textbf{Équation-bilan} : Réactifs $\rightarrow$ Produits. \textbf{Seuls les coefficients changent} (devant). Indices intouchables.
    \item \textbf{Méthode d'équilibrage} : Formules $\rightarrow$ Comptage $\rightarrow$ Coefficients (C, H, O, autres) $\rightarrow$ Vérification.
    \item \textbf{Combustion complète} (O2 suffisant) $\rightarrow$ \ce{CO2} + \ce{H2O} (bleue, max énergie).
    \item \textbf{Combustion incomplète} (O2 manquant) $\rightarrow$ \ce{CO} (mortel) + \ce{C} (suie) + \ce{H2O} (jaune, sale).
    \item \textbf{Sécurité CO} : Ventilation permanente, entretien, détecteur \ce{CO}, conduite à tenir (Aérer, Sortir, Appeler 112/119).
\end{enumerate}
\end{aretenir}

\coursec{Combustions complète et incomplète}

Dans la section précédente, nous avons appris à écrire et à équilibrer l'équation-bilan d'une réaction chimique. Nous savons maintenant que les atomes se conservent et se réarrangent. Appliquons cette rigueur à une réaction familière : la combustion. Pourquoi la flamme de ta gazinière est-elle bleue alors que celle d'un feu de bois ou d'une lampe à pétrole est jaune et fumeuse ? La réponse tient en un mot : \cle{dioxygène}.

\courssub{Rappel : c'est quoi, une combustion ?}

Une \cle{combustion} est une réaction chimique rapide entre un \cle{combustible} (qui brûle) et un \cle{comburant} (qui entretient la flamme), déclenchée par une \cle{énergie d'activation} (étincelle, flamme, chaleur). Dans notre vie quotidienne au Cameroun, le comburant est presque toujours le \cle{dioxygène} $\ce{O2}$ contenu dans l'air (environ 21 \%). Les combustibles sont variés : le \cle{carbone} $C$ (charbon de bois), le \cle{butane} $\ce{C4H10}$ (bonbonne bleue), l'\cle{essence} et le \cle{gasoil} (station SONARA), le \cle{bois} (cuisine traditionnelle), l'\cle{alcool} (lampe à pétrole ou gel hydroalcoolique).

\begin{experience}[Combustion du carbone dans le dioxygène pur]
\textbf{Matériel :} creuset, charbon de bois (carbone), cuillère à déflagration, flacon à large col rempli de dioxygène $\ce{O2}$ (recueilli par déplacement de l'eau), eau de chaux.
\textbf{Protocole :}
\begin{enumerate}
\item Enflammer le charbon dans la cuillère à déflagration (énergie d'activation).
\item Plonger la cuillère dans le flacon de $\ce{O2}$.
\item Observer la combustion.
\item Après refroidissement, verser de l'eau de chaux dans le flacon, boucher et agiter.
\end{enumerate}
\textbf{Observations :}
\begin{itemize}
\item Le charbon brûle avec une \textbf{vive incandescence} (blanc éclatant), bien plus intense qu'à l'air libre.
\item L'eau de chaux \textbf{devient trouble} (blanchiment).
\end{itemize}
\textbf{Interprétation :} Le trouble de l'eau de chaux prouve la formation de \cle{dioxyde de carbone} $\ce{CO2}$. Le carbone a réagi avec le dioxygène pour former du $\ce{CO2}$.
\[
\ce{C + O2 -> CO2}
\]
\textbf{Conclusion :} En présence de dioxygène \textbf{en excès} (pur ici), le carbone brûle \textbf{totalement} en dioxyde de carbone. C'est une \textbf{combustion complète}.
\end{experience}

\courssub{Combustion complète : le dioxygène est en quantité suffisante}

\begin{definition}[Combustion complète]
Une \textbf{combustion complète} est une combustion où le comburant (dioxygène) est en \textbf{quantité suffisante} pour oxyder \textbf{totalement} tous les atomes du combustible.
Pour un combustible carboné (contenant $C$ et $H$), les produits sont \textbf{uniquement} du \cle{dioxyde de carbone} $\ce{CO2}$ et de l'\cle{eau} $\ce{H2O}$.
\end{definition}

\begin{exemplebox}[Gazinière bien réglée : combustion complète du butane]
Le butane $\ce{C4H10}$ de ta bonbonne réagit avec l'oxygène de l'air.
Équation-bilan équilibrée :
\[
\ce{2 C4H10 + 13 O2 -> 8 CO2 + 10 H2O}
\]
\textbf{Signes visibles :} Flamme \textbf{bleue}, stable, silencieuse, \textbf{peu ou pas de fumée}. C'est le signe d'un bon réglage de l'arrivée d'air (aération primaire) sur ta gazinière.
\end{exemplebox}

\courssub{Combustion incomplète : le dioxygène manque}

\begin{definition}[Combustion incomplète]
Une \textbf{combustion incomplète} est une combustion où le comburant (dioxygène) est en \textbf{quantité insuffisante} pour oxyder totalement le combustible.
Pour un combustible carboné, il se forme un \textbf{mélange} de \cle{dioxyde de carbone} $\ce{CO2}$, de \cle{monoxyde de carbone} $\ce{CO}$ (gaz incolore, inodore, très toxique) et/ou de \cle{carbone} $C$ sous forme de \cle{suie} (poudre noire).
\end{definition}

\begin{experience}[La flamme du bec Bunsen : complète vs incomplète]
\textbf{Matériel :} Bec Bunsen (ou bec de gaz de TP), arrivée de gaz, arrivée d'air réglable (collerette).
\textbf{Protocole :}
\begin{enumerate}
\item Ouvrir le gaz, allumer. Fermer complètement l'arrivée d'air (collerette basse). Observer.
\item Ouvrir progressivement l'arrivée d'air au maximum. Observer.
\end{enumerate}
\textbf{Observations et interprétation :}
\end{experience}
\begin{center}
\begin{popfigure}
\begin{tikzpicture}[scale=0.9]
% Flamme INCOMPLETE (jaune) - Air fermé
\begin{scope}[xshift=-4cm]
% Base
\draw[popdark, thick] (-0.5,0) -- (-0.5,1) -- (0.5,1) -- (0.5,0) -- cycle;
\draw[popdark, thick] (-0.8,1) -- (-0.8,1.2) -- (0.8,1.2) -- (0.8,1);
% Flamme jaune fumeuse
\draw[poporange!80!popdark, thick, decorate, decoration={random steps,segment length=4pt,amplitude=2pt}] 
(-0.4,1.2) .. controls (-0.6,2) and (0.6,2.5) .. (0,3.5) .. controls (-0.6,2.5) and (-0.4,2) .. (-0.4,1.2);
\fill[poporange!60!popdark, opacity=0.6] (-0.4,1.2) .. controls (-0.6,2) and (0.6,2.5) .. (0,3.5) .. controls (-0.6,2.5) and (-0.4,2) .. (-0.4,1.2);
% Particules de suie
\foreach \i in {1,...,15} {
    \pgfmathsetmacro{\x}{rand*0.3}
    \pgfmathsetmacro{\y}{1.2+rand*2.3}
    \fill[popdark] (\x,\y) circle (0.03);
}
% Flèches air (peu)
\draw[popblue, ->, thick] (-1.5,0.5) -- (-0.9,0.5) node[midway, above, font=\tiny] {Peu d'air};
\node[align=center, font=\small\bfseries] at (0,-0.5) {Arrivée d'air \textbf{FERMÉE}};
\node[align=center, font=\small, poporange] at (0,4) {Flamme \textbf{JAUNE}\\\textbf{Fumeuse} (suie)};
\end{scope}

% Flamme COMPLETE (bleue) - Air ouvert
\begin{scope}[xshift=4cm]
% Base
\draw[popdark, thick] (-0.5,0) -- (-0.5,1) -- (0.5,1) -- (0.5,0) -- cycle;
\draw[popdark, thick] (-0.8,1) -- (-0.8,1.2) -- (0.8,1.2) -- (0.8,1);
% Flamme bleue (deux cônes)
% Cône interne (bleu foncé)
\draw[popblue!80!popdark, thick] (-0.25,1.2) -- (0,2.8) -- (0.25,1.2) -- cycle;
\fill[popblue!60!popdark, opacity=0.8] (-0.25,1.2) -- (0,2.8) -- (0.25,1.2) -- cycle;
% Cône externe (bleu clair)
\draw[popblue!40!popgreen, thick, decorate, decoration={random steps,segment length=3pt,amplitude=1pt}] 
(-0.4,1.2) .. controls (-0.5,2) and (0.5,2.5) .. (0,3.2) .. controls (-0.5,2.5) and (-0.4,2) .. (-0.4,1.2);
\fill[popblue!30!popgreen, opacity=0.4] (-0.4,1.2) .. controls (-0.5,2) and (0.5,2.5) .. (0,3.2) .. controls (-0.5,2.5) and (-0.4,2) .. (-0.4,1.2);
% Flèches air (beaucoup)
\draw[popblue, ->, thick] (-1.5,0.5) -- (-0.9,0.5) node[midway, above, font=\tiny] {Beaucoup d'air};
\draw[popblue, ->, thick] (1.5,0.5) -- (0.9,0.5) node[midway, above, font=\tiny] {Beaucoup d'air};
\node[align=center, font=\small\bfseries] at (0,-0.5) {Arrivée d'air \textbf{OUVERTE}};
\node[align=center, font=\small, popblue] at (0,4) {Flamme \textbf{BLEUE}\\\textbf{Non fumeuse}};
\end{scope}
\end{tikzpicture}
\legende{Comparaison des flammes du bec Bunsen : arrivée d'air fermée (combustion incomplète, jaune, suie) vs arrivée d'air ouverte (combustion complète, bleue, propre).}
\end{popfigure}
\end{center}

\begin{propriete}[Produits selon la quantité de dioxygène]
La nature des produits d'une combustion d'un combustible carboné dépend \textbf{uniquement} de la quantité de dioxygène $\ce{O2}$ disponible.
\begin{itemize}
\item \textbf{Dioxygène en excès (ou suffisant)} $\rightarrow$ Combustion \textbf{complète} : $\ce{CO2} + \ce{H2O}$.
\item \textbf{Dioxygène insuffisant (manque d'air)} $\rightarrow$ Combustion \textbf{incomplète} : $\ce{CO2} + \ce{CO} + \ce{C}$ (suie) $+ \ce{H2O}$.
\end{itemize}
Cette propriété est établie par l'expérience du bec Bunsen : en variant l'arrivée d'air, on passe d'une flamme jaune (incomplète) à une flamme bleue (complète).
\end{propriete}

\courssub{Équations-bilan de la combustion incomplète du carbone}

Pour le carbone pur (charbon), deux équations incomplètes principales sont possibles selon le degré de manque d'oxygène.

\begin{propriete}[Équations de la combustion incomplète du carbone]
\begin{enumerate}
\item \textbf{Formation de monoxyde de carbone} (manque modéré d'$\ce{O2}$) :
\[
\ce{2 C + O2 -> 2 CO}
\]
\item \textbf{Formation de suie (carbone non brûlé)} (manque sévère d'$\ce{O2}$) :
Le carbone en excès ne réagit pas et se dépose sous forme de poudre noire ($C$).
\end{enumerate}
\textbf{Attention :} Dans la réalité (feu de bois, moteur de moto-taxi mal entretenu), les trois réactions ont lieu \textbf{simultanément} : on obtient un mélange de $\ce{CO2}$, $\ce{CO}$ et de suie $C$.
\end{propriete}

\begin{methode}[Choisir l'équation-bilan correcte]
Pour écrire l'équation-bilan d'une combustion :
\begin{enumerate}
\item \textbf{Identifier le combustible} (formule brute : $C$, $\ce{CH4}$, $\ce{C4H10}$, etc.).
\item \textbf{Identifier le type de combustion} d'après l'énoncé ou les indices (flamme bleue = complète ; flamme jaune/fumeuse, manque d'air, espace confiné = incomplète).
\item \textbf{Écrire les produits} :
\begin{itemize}
\item Complète : $\ce{CO2}$ et $\ce{H2O}$ (s'il y a H).
\item Incomplète : $\ce{CO2}$, $\ce{CO}$, $\ce{C}$ (suie) et $\ce{H2O}$.
\end{itemize}
\item \textbf{Équilibrer} (atomes $C$, puis $H$, puis $O$).
\end{enumerate}
\end{methode}

\begin{exoresolu}[Équilibrage combustion incomplète du butane]
\textbf{Énoncé :} Une bonbonne de butane $\ce{C4H10}$ brûle dans une pièce mal aérée (manque d'oxygène). On observe une flamme jaune et le dépôt de suie. Écrire l'équation-bilan d'une combustion incomplète possible produisant du $\ce{CO2}$, du $\ce{CO}$ et du $C$.

\textbf{Solution détaillée :}
\begin{enumerate}
\item \textbf{Réactifs :} $\ce{C4H10} + \ce{O2}$.
\item \textbf{Produits (incomplète) :} $\ce{CO2} + \ce{CO} + \ce{C} + \ce{H2O}$.
\item \textbf{Équation squelette :} $\ce{C4H10 + O2 -> a CO2 + b CO + c C + d H2O}$.
\item \textbf{Bilan Carbone ($C$) :} $4 = a + b + c$. \quad (On choisit par ex. $a=1, b=2, c=1$).
\item \textbf{Bilan Hydrogène ($H$) :} $10 = 2d \Rightarrow d = 5$.
\item \textbf{Bilan Oxygène ($O$) :} $2 = 2a + b + d = 2(1) + 2 + 5 = 9$. \quad Il faut $9/2 = 4,5$ molécules $\ce{O2}$.
\item \textbf{Équation équilibrée (coefficients entiers) :} On multiplie par 2.
\[
\ce{2 C4H10 + 9 O2 -> 2 CO2 + 4 CO + 2 C + 10 H2O}
\]
\textbf{Vérification :} $C: 8=2+4+2=8$ ; $H: 20=20$ ; $O: 18=4+4+10=18$. \checkmark
\end{enumerate}
\end{exoresolu}

\courssub{Le monoxyde de carbone : un danger invisible}

\begin{savaistu}[Le monoxyde de carbone (CO), tueur silencieux]
Le \cle{monoxyde de carbone} $\ce{CO}$ est un gaz \textbf{incolore}, \textbf{inodore} et \textbf{non irritant}. On ne le sent pas, on ne le voit pas, il ne pique pas les yeux. Pourtant, c'est un poison violent.
\begin{itemize}
\item \textbf{Mécanisme :} $\ce{CO}$ se fixe sur l'hémoglobine du sang (à la place de l'$\ce{O2}$) avec une affinité 200 à 250 fois plus grande. Il forme de la \cle{carboxyhémoglobine}. Le sang ne peut plus transporter l'oxygène $\rightarrow$ asphyxie cellulaire.
\item \textbf{Symptômes :} Maux de tête, nausées, vertiges, confusion, perte de connaissance, \textbf{mort} en quelques minutes si concentration élevée.
\item \textbf{Situations à risque au Cameroun :} Brasero ou feu de charbon dans une chambre close (portes/fenêtres fermées), groupe électrogène placé dans un couloir ou près d'une fenêtre ouverte, gazinière mal réglée (flamme jaune) dans une cuisine non aérée, moto-taxi (bensikin) dans les embouteillages (inhalation des gaz d'échappement).
\end{itemize}
\end{savaistu}

\courssub{Sécurité et gestes responsables autour des combustions}

\begin{aretenir}[Comparatif complet et prévention]
\begin{center}
\begin{popfigure}
\begin{tikzpicture}
\node[inner sep=0pt, align=center] (table){
\begin{tabularx}{\linewidth}{|>{\centering\arraybackslash}X|>{\centering\arraybackslash}X|>{\centering\arraybackslash}X|}
\hline
\rowcolor{popblue!20} 
\textbf{Critère} & \textbf{Combustion COMPLÈTE} & \textbf{Combustion INCOMPLÈTE} \\ \hline
\textbf{Dioxygène $\ce{O2}$} & Quantité suffisante (excès) & Quantité insuffisante (manque d'air) \\ \hline
\textbf{Produits (combustible $C_xH_y$)} & $\ce{CO2} + \ce{H2O}$ \textbf{uniquement} & Mélange $\ce{CO2} + \ce{CO} + \ce{C (suie)} + \ce{H2O}$ \\ \hline
\textbf{Couleur flamme} & \textcolor{popblue}{\textbf{Bleue}} (deux cônes nets) & \textcolor{poporange}{\textbf{Jaune / Orange}} (fumeuse, instable) \\ \hline
\textbf{Fumée / Suie} & Aucune (propre) & \textbf{Abondante} (suie noire, encrassement) \\ \hline
\textbf{Énergie dégagée} & \textbf{Maximale} (rendement optimal) & \textbf{Moins élevée} (perte d'énergie dans CO et C) \\ \hline
\textbf{Dangers principaux} & $\ce{CO2}$ (gaz à effet de serre), brûlures & \textbf{CO (mortel)}, suie (cancérigène, encrasse moteurs), $\ce{CO2}$ \\ \hline
\textbf{Exemples au Cameroun} & Gazinière bien réglée, moto-taxi bien entretenue, groupe électrogène récent & Feu de bois (marmite noircie), brasero en chambre close, vieux moteur 2-temps qui fume, gazinière mal réglée \\ \hline
\end{tabularx}
};
\end{tikzpicture}
\legende{Tableau comparatif : combustion complète vs incomplète.}
\end{popfigure}
\end{center}
\end{aretenir}

\begin{savaistu}[Le feu de bois et la marmite noircie]
Quand tu cuisines sur un feu de bois (trois pierres) ou un brasero en terre cuite, la marmite noircit vite. Ce dépôt noir, c'est de la \textbf{suie} (carbone $C$ non brûlé). C'est la preuve visuelle d'une \textbf{combustion incomplète} : l'air ne circule pas assez bien entre les bûches pour apporter l'oxygène nécessaire.
\begin{itemize}
\item \textbf{Perte d'argent :} Le bois (ou le charbon) n'est pas brûlé totalement $\rightarrow$ gaspillage de combustible.
\item \textbf{Santé :} La fumée contient du $\ce{CO}$ et des particules fines (suie) qui irritent les poumons et les yeux (maladies respiratoires fréquentes chez les femmes cuisinant au feu de bois).
\item \textbf{Environnement :} La suie (carbone noir) est un puissant réchauffant climatique (absorbe la chaleur du soleil).
\end{itemize}
Les \textbf{foyers améliorés} (type « foyers ronds » ou « foyers à économie d'énergie » diffusés par le MINEE et les ONG) forcent l'air à passer dans le feu $\rightarrow$ combustion plus complète $\rightarrow$ moins de bois, moins de fumée, marmite propre !
\end{savaistu}

\begin{attention}[Pièges fréquents au BEPC]
\begin{itemize}
\item \textbf{Erreur 1 :} « Combustion incomplète = production de CO \textbf{uniquement} ». \textbf{Faux !} C'est presque toujours un \textbf{mélange} $\ce{CO2} + \ce{CO} + \ce{C}$. L'équation $\ce{2C + O2 -> 2CO}$ est un \textbf{cas particulier} (modèle simplifié), pas la réalité générale.
\item \textbf{Erreur 2 :} Confondre \textbf{comburant} et \textbf{combustible}. Mnémo : le \textbf{Comburant} \textbf{C}ontient l'\textbf{O}xygène ($\ce{O2}$) ; le \textbf{Combustible} \textbf{C}ontient le \textbf{C}arbone (et/ou H).
\item \textbf{Erreur 3 :} Oublier l'eau $\ce{H2O}$ dans les produits si le combustible contient de l'hydrogène (butane, essence, bois, alcool). Même en combustion incomplète, l'hydrogène brûle presque toujours en $\ce{H2O}$.
\item \textbf{Erreur 4 :} Écrire $\ce{C + O2 -> CO}$ (non équilibré : 1 C à gauche, 1 C à droite mais 2 O à gauche, 1 O à droite). La bonne équation pour le CO est $\ce{2C + O2 -> 2CO}$.
\end{itemize}
\end{attention}

\courssub{Exemples concrets : la gazinière et le moto-taxi}

\begin{exemplebox}[Gazinière : régler l'air pour économiser et vivre]
Sur ta gazinière (butane $\ce{C4H10}$), il y a une \textbf{collerette} (bague métallique) à la base du brûleur.
\begin{itemize}
\item \textbf{Collerette FERMÉE} : Peu d'air $\rightarrow$ Flamme \textbf{jaune}, haute, fumeuse $\rightarrow$ Combustion \textbf{incomplète} $\rightarrow$ Marmite noircie (suie), \textbf{gaz gaspillé}, risque de $\ce{CO}$ dans la cuisine.
\item \textbf{Collerette OUVERTE} : Beaucoup d'air $\rightarrow$ Flamme \textbf{bleue}, courte, silencieuse $\rightarrow$ Combustion \textbf{complète} $\rightarrow$ Marmite propre, \textbf{cuisson plus rapide}, gaz économisé, air sain.
\end{itemize}
\textbf{Geste responsable :} Vérifie et ouvre la collerette à chaque utilisation ! Aère la cuisine (ouvre les fenêtres).
\end{exemplebox}

\begin{exemplebox}[Moto-taxi (bensikin) et vieux moteurs]
Tu as sûrement vu ces motos qui crachent une \textbf{fumée bleue/noire épaisse} à l'échappement.
\begin{itemize}
\item \textbf{Fumée noire} = \textbf{Suie} (carbone $C$) $\rightarrow$ Combustion incomplète (mélange trop riche en essence, filtre à air bouché, bougie usée).
\item \textbf{Fumée bleue} = \textbf{Huile brûlée} (moteur usé) + combustion incomplète.
\end{itemize}
Conséquences : \textbf{Moteur encrassé} (perte de puissance, surchauffe), \textbf{essence gaspillée} (le conducteur perd de l'argent), \textbf{pollution} grave pour les piétons et le conducteur lui-même (inhalation de $\ce{CO}$ et particules). \textbf{Entretenir son moteur (vidange, filtre à air, bougie), c'est faire une combustion plus complète !}
\end{exemplebox}

\begin{exercice}[Entraînement : Identifier et équilibrer]
\textbf{1.} Complète le tableau ci-dessous pour la combustion du méthane $\ce{CH4}$.
\begin{center}
\begin{tabularx}{\linewidth}{|>{\centering\arraybackslash}X|>{\centering\arraybackslash}X|>{\centering\arraybackslash}X|}
\hline
\textbf{Type de combustion} & \textbf{Produits (formules)} & \textbf{Équation-bilan équilibrée} \\ \hline
Complète & & \\ \hline
Incomplète (avec CO et C) & & \\ \hline
\end{tabularx}
\end{center}
\textbf{2.} Un élève écrit pour la combustion incomplète du carbone : $\ce{C + O2 -> CO}$. 
\begin{enumerate}
\item Cette équation est-elle équilibrée ? Justifie.
\item Propose l'équation correcte pour la formation de $\ce{CO}$.
\end{enumerate}
\textbf{3.} Pourquoi la flamme d'une bougie est-elle jaune et celle d'un bec Bunsen (air ouvert) est-elle bleue ? Explique en utilisant les mots \cle{suie}, \cle{incandescence}, \cle{dioxygène}.
\end{exercice}

\begin{corrige}[Exercice n°4 - Combustions complète et incomplète]
\textbf{1.}
\begin{center}
\begin{tabularx}{\linewidth}{|>{\centering\arraybackslash}X|>{\centering\arraybackslash}X|>{\centering\arraybackslash}X|}
\hline
\textbf{Type} & \textbf{Produits} & \textbf{Équation équilibrée} \\ \hline
Complète & $\ce{CO2}$, $\ce{H2O}$ & $\ce{CH4 + 2 O2 -> CO2 + 2 H2O}$ \\ \hline
Incomplète & $\ce{CO2}$, $\ce{CO}$, $\ce{C}$, $\ce{H2O}$ & Exemple : $\ce{6 CH4 + 9 O2 -> 2 CO2 + 2 CO + 2 C + 12 H2O}$ \\ \hline
\end{tabularx}
\end{center}
\textit{(Note : plusieurs équations incomplètes sont possibles selon les proportions de $\ce{CO2}$/$\ce{CO}$/$C$ choisies, du moment qu'elles sont équilibrées. L'exemple ci-dessus produit les quatre espèces demandées.)}

\textbf{2.}
\begin{enumerate}
\item \textbf{Non.} Bilan atomes O : 2 à gauche (dans $\ce{O2}$), 1 à droite (dans $\ce{CO}$). Bilan atomes C : 1=1 (OK). L'oxygène n'est pas conservé.
\item \textbf{Équation correcte :} $\ce{2 C + O2 -> 2 CO}$. (2 C, 2 O de chaque côté).
\end{enumerate}

\textbf{3.} 
\begin{itemize}
\item \textbf{Bougie (flamme jaune) :} L'arrivée d'air se fait seulement par diffusion au pourtour de la flamme (pas d'arrivée d'air forcée). Le dioxygène \textbf{manque} au cœur de la flamme $\rightarrow$ \textbf{combustion incomplète} $\rightarrow$ formation de particules de \textbf{suie} (carbone $C$). Ces particules très chaudes deviennent \textbf{incandescentes} (jaune/orange) $\rightarrow$ lumière jaune.
\item \textbf{Bec Bunsen (air ouvert, flamme bleue) :} L'arrivée d'air est \textbf{forcée et pré-mélangée} au gaz (aération primaire). Le dioxygène est \textbf{en excès} dès le départ $\rightarrow$ \textbf{combustion complète} $\rightarrow$ pas de suie, pas d'incandescence de particules solides. La couleur bleue vient de l'émission de lumière par les molécules excitées ($\ce{CH*}$, $\ce{C2*}$) dans la zone de réaction (chimiluminescence).
\end{itemize}
\end{corrige}

\coursec{Le monoxyde de carbone : un danger invisible}

\courssub{Transition : de la flamme au danger invisible}

Dans la section précédente, nous avons vu qu'une \cle{combustion incomplète} se produit lorsque le dioxygène manque. Elle produit alors du \cle{dioxyde de carbone} ($\ce{CO2}$) mais aussi, et surtout, du \cle{monoxyde de carbone} ($\ce{CO}$). Contrairement au $\ce{CO2}$ qui est un produit « normal » de la respiration et de la combustion complète, le $\ce{CO}$ est un \textbf{poison violent}. C'est le sujet de cette section : comprendre pourquoi ce gaz est si dangereux, comment il agit sur notre organisme, et surtout comment s'en protéger au quotidien, ici au Cameroun.

\courssub{Qu'est-ce que le monoxyde de carbone ?}

\begin{definition}[Monoxyde de carbone (CO)]
Le \textbf{monoxyde de carbone} est un gaz de formule brute \ce{CO}. Il est \textbf{incolore}, \textbf{inodore} (sans odeur), \textbf{sans saveur} et \textbf{non irritant} pour les yeux ou la gorge. Il se mélange parfaitement à l'air (légèrement plus léger, il se diffuse rapidement).
\end{definition}

\begin{attention}[Le piège des sens]
\emph{« Je ne sens rien, donc il n'y a pas de danger. »} \textbf{FAUX.} C'est précisément l'absence d'odeur, de couleur et d'irritation qui fait du $\ce{CO}$ un \textbf{tueur invisible}. Nos sens (odorat, vue) sont totalement incapables de le détecter. Seule une alarme spécifique (détecteur de CO) ou l'apparition de symptômes peut nous alerter, souvent trop tard.
\end{attention}

\begin{propriete}[Origine domestique du CO]
Le monoxyde de carbone provient \textbf{uniquement} de \cle{combustions incomplètes} d'origine carbonée (bois, charbon, essence, gaz butane/propane, gasoil, papier, plastiques).
Au Cameroun, les sources fréquentes sont :
\begin{itemize}
    \item Les \textbf{braseros} (feu de bois ou charbon de bois) utilisés pour chauffer la pièce ou faire griller du poisson/maïs dans le salon ou la chambre.
    \item Les \textbf{groupes électrogènes} (essence ou gasoil) placés dans une véranda, un garage, un couloir ou une pièce mal ventilée pendant les coupures \textsc{eneo}.
    \item Les \textbf{gazinières} ou fours à gaz mal réglés (flamme jaune/orange au lieu de bleue) ou utilisés pour se chauffer porte ouverte.
    \item Les \textbf{moteurs thermiques} (moto-taxi, voiture, pompe à eau) laissés tourner au ralenti dans un garage fermé ou une cour close.
    \item Les \textbf{chauffe-eau} à gaz non évacués vers l'extérieur (cheminée bouchée ou inexistante).
\end{itemize}
\end{propriete}

\begin{savaistu}[Contexte camerounais : la saison fraîche et les délestages]
Chaque année, pendant la saison fraîche (novembre--février) et lors des délestages électriques, les urgences des hôpitaux (Laquintinie, Hôpital Central, Hôpital Général de Yaoundé, etc.) reçoivent des familles entières intoxiquées au monoxyde de carbone. L'utilisation de braseros pour « chasser le froid » ou de groupes électrogènes placés « juste à côté de la fenêtre » (mais à l'intérieur) sont les causes numéro un. La prévention est notre meilleure arme.
\end{savaistu}

\courssub{Comment le CO nous tue-t-il ? Le mécanisme de la toxicité}

Nous ne pouvons pas « voir » le CO agir dans le sang au collège, mais la médecine a établi le mécanisme exact. C'est un \textbf{savoir admis} au niveau 3ème, fondé sur des observations cliniques et biochimiques rigoureuses.

\begin{experience}[Modélisation de la fixation du CO sur l'hémoglobine (Démarche d'investigation)]
\textbf{Question de départ :} Pourquoi l'inhalation d'un gaz inodore empêche-t-elle le sang de transporter l'oxygène ?

\textbf{Document d'observation :} Analyse de sang d'une victime d'intoxication au CO.
\begin{itemize}
    \item Sang sain : rouge vif (artériel) / rouge foncé (veineux).
    \item Sang intoxiqué au CO : \textbf{rouge cerise vif} (couleur « sang de homard »), même dans les veines.
\end{itemize}

\textbf{Hypothèse :} Le CO modifie la couleur de l'hémoglobine et l'empêche de relâcher l'oxygène (ou de le prendre).

\textbf{Interprétation scientifique (savoir établi) :}
L'hémoglobine (protéine des globules rouges) possède 4 sites de fixation pour l'oxygène ($\ce{O2}$). Le monoxyde de carbone ($\ce{CO}$) a une \textbf{affinité pour l'hémoglobine environ 200 à 250 fois plus grande} que celle du dioxygène.
\begin{enumerate}
    \item Quand on respire de l'air contenant du $\ce{CO}$, celui-ci traverse les poumons et arrive dans le sang.
    \item Il « vole » la place de l'$\ce{O2}$ sur l'hémoglobine : $\ce{Hb + CO <=> HbCO}$ (carboxyhémoglobine).
    \item L'hémoglobine ainsi bloquée (\ce{HbCO}) ne peut plus transporter l'$\ce{O2}$ vers les organes (cerveau, cœur, muscles).
    \item De plus, la présence de $\ce{CO}$ sur une partie de l'hémoglobine \textbf{renforce la fixation de l'$\ce{O2}$ sur les sites restants} : l'hémoglobine ne lâche plus l'oxygène là où les cellules en ont besoin (effet Bohr inversé).
\end{enumerate}

\textbf{Conclusion :} Le sang circule, mais il ne livre plus l'oxygène. C'est une \textbf{asphyxie cellulaire} (hypoxie toxique) alors que la respiration pulmonaire fonctionne normalement.
\end{experience}

\begin{popfigure}
\begin{tikzpicture}[scale=0.9, every node/.style={font=\small}]
    % Couleurs
    \colorlet{rbc}{popred!80!popdark}
    \colorlet{o2col}{popblue}
    \colorlet{cocol}{poporange}
    \colorlet{hbcol}{poppurple!70!popdark}
    
    % Légende haut
    \node[align=center] at (0,5.5) {\textbf{Globule rouge (érythrocyte) -- Vue schématique}};
    
    % --- CAS 1 : Normal ---
    \begin{scope}[xshift=-6cm]
        \node[align=center] at (0,4.2) {\textbf{Normal : Transport de l'O$_2$}};
        % Globule rouge
        \draw[rbc, thick, fill=rbc!20] (0,0) ellipse (2.2cm and 1.2cm);
        % Hémoglobines (4 groupes)
        \foreach \angle/\r in {45/1.2, 135/1.2, 225/1.2, 315/1.2} {
            \node[hbcol, circle, draw=hbcol, thick, minimum size=0.6cm] (hb\angle) at (\angle:\r) {\ce{Hb}};
        }
        % O2 fixés
        \foreach \angle/\r in {45/1.2, 135/1.2, 225/1.2, 315/1.2} {
            \node[o2col, circle, draw=o2col, fill=o2col!30, minimum size=0.35cm] at (\angle:\r+0.7) {\ce{O2}};
            \draw[o2col, ->, thick] (\angle:\r+0.35) -- (\angle:\r+0.65);
        }
        % Flèche apport O2
        \draw[o2col, ->, thick] (-3.5, -1.5) -- (-1.5, -0.5) node[midway, above, sloped] {Apport \ce{O2} (poumons)};
        % Flèche livraison O2
        \draw[o2col, ->, thick] (1.5, -0.5) -- (3.5, -1.5) node[midway, above, sloped] {Livraison \ce{O2} (cellules)};
    \end{scope}

    % --- CAS 2 : Intoxication CO ---
    \begin{scope}[xshift=6cm]
        \node[align=center] at (0,4.2) {\textbf{Intoxication : Blocage par le CO}};
        % Globule rouge
        \draw[rbc, thick, fill=rbc!20] (0,0) ellipse (2.2cm and 1.2cm);
        % Hémoglobines bloquées par CO
        \foreach \angle/\r in {45/1.2, 135/1.2, 225/1.2, 315/1.2} {
            \node[hbcol, circle, draw=hbcol, thick, minimum size=0.6cm] (hb\angle) at (\angle:\r) {\ce{Hb}};
        }
        % CO fixés (majorité)
        \foreach \angle/\r in {45/1.2, 135/1.2, 225/1.2} {
            \node[cocol, circle, draw=cocol, fill=cocol!30, minimum size=0.35cm] at (\angle:\r+0.7) {\ce{CO}};
            \draw[cocol, ->, thick] (\angle:\r+0.35) -- (\angle:\r+0.65);
            % Croix rouge sur O2
            \draw[popred, thick] (\angle:\r+0.7) ++(-0.2,-0.2) -- ++(0.4,0.4);
            \draw[popred, thick] (\angle:\r+0.7) ++(0.2,-0.2) -- ++(-0.4,0.4);
        }
        % Un seul O2 coincé (ne part pas)
        \node[o2col, circle, draw=o2col, fill=o2col!30, minimum size=0.35cm] at (315:1.9) {\ce{O2}};
        \draw[o2col, thick] (315:1.55) -- (315:1.85); % pas de flèche sortante
        \node[o2col, align=center] at (315:2.5) {\ce{O2} \emph{piégé}};
        
        % Flèche apport CO
        \draw[cocol, ->, thick] (-3.5, -1.5) -- (-1.5, -0.5) node[midway, above, sloped] {Inhalation \ce{CO}};
        % Pas de livraison
        \draw[popred, thick, dashed, ->] (1.5, -0.5) -- (3.5, -1.5) node[midway, above, sloped] {\textbf{Aucune livraison \ce{O2}}};
        
        % Légende affinité
        \node[align=center, fill=popcream, draw=popline, rounded corners, inner sep=3pt] at (0, -3.2) {
            Affinité \ce{Hb--CO} $\approx$ \textbf{200--250$\times$} Affinité \ce{Hb--O2}
        };
    \end{scope}
\end{tikzpicture}
\legende{Schéma comparatif : à gauche, l'hémoglobine (Hb) charge et décharge normalement le dioxygène ($\ce{O2}$). À droite, le monoxyde de carbone ($\ce{CO}$) occupe les sites de fixation (affinité 200$\times$ supérieure), piégeant le peu d'$\ce{O2}$ restant et empêchant sa libération dans les tissus.}
\end{popfigure}

\begin{propriete}[Mécanisme de toxicité du CO (Savoir admis BEPC)]
Le monoxyde de carbone ($\ce{CO}$) se fixe sur le \textbf{fer de l'hémoglobine} (\ce{Hb}) à la place du dioxygène ($\ce{O2}$) pour former de la \textbf{carboxyhémoglobine} (\ce{HbCO}).
\begin{itemize}
    \item \textbf{Affinité :} $K_{\ce{CO}} / K_{\ce{O2}} \approx 200\text{--}250$. Le $\ce{CO}$ « gagne » la compétition même à très faible concentration dans l'air.
    \item \textbf{Conséquence 1 :} Diminution de la capacité de transport de l'$\ce{O2}$ (anémie fonctionnelle).
    \item \textbf{Conséquence 2 :} Déplacement de la courbe de dissociation de l'hémoglobine vers la gauche : l'$\ce{O2}$ restant est « collé » à l'hémoglobine et n'est pas libéré au niveau des capillaires tissulaires.
    \item \textbf{Résultat :} \textbf{Hypoxie tissulaire sévère} (cerveau et cœur en premier) $\rightarrow$ coma, séquelles neurologiques, mort.
\end{itemize}
\end{propriete}

\courssub{Symptômes : reconnaître l'intoxication}

Les symptômes dépendent de la concentration en $\ce{CO}$ dans l'air (ppm -- parties par million) et de la durée d'exposition. Ils sont \textbf{non spécifiques} (ressemblent à une grippe, une gastro-entérite, une fatigue, une ivresse), ce qui retarde le diagnostic.

\begin{center}
\begin{tabularx}{\linewidth}{|l|X|X|}\hline
\rowcolor{popblueL}
\textbf{Niveau d'exposition} & \textbf{Symptômes principaux} & \textbf{Gravité / Conduite} \\ \hline
\textbf{Faible / Chronique} & Maux de tête légers, fatigue inhabituelle, nausées, vertiges, troubles de la vision, irritabilité. & Souvent confondus avec grippe/paludisme. \textbf{Aérer, sortir, consulter.} \\ \hline
\textbf{Moyen / Aigu} & Maux de tête violents, confusion, vomissements, essoufflement, tachycardie, troubles de l'équilibre, \textbf{perte de conscience}. & \textbf{Urgence vitale.} Évacuation immédiate, oxygénothérapie. \\ \hline
\textbf{Élevé / Mortel} & Coma profond, convulsions, arrêt cardio-respiratoire, \textbf{mort} en quelques minutes. & \textbf{Décès rapide} sans réanimation spécialisée (caisson hyperbare). \\ \hline
\end{tabularx}
\end{center}

\begin{attention}[Populations à risque accru]
Les \textbf{enfants}, les \textbf{femmes enceintes} (le $\ce{CO}$ traverse le placenta et fixe l'hémoglobine fœtale avec une affinité encore plus grande), les \textbf{personnes âgées} et les \textbf{malades cardiaques ou respiratoires} sont touchés plus vite et plus gravement. Un nourrisson peut mourir pendant que l'adulte n'a qu'un mal de tête.
\end{attention}

\courssub{Conduite à tenir : les 4 réflexes qui sauvent}

\begin{methode}[Réflexes de survie face à une suspicion d'intoxication au CO]
Si vous suspectez une intoxication (plusieurs personnes malades en même temps, symptômes qui disparaissent en sortant de la pièce, groupe électrogène/brazero allumé) :
\begin{enumerate}
    \item \textbf{AÉRER IMMÉDIATEMENT :} Ouvrez \textbf{en grand} toutes les portes et fenêtres. Créez un courant d'air. Ne perdez pas de temps à chercher la fuite.
    \item \textbf{ÉVACUER LES VICTIMES :} Sortez tout le monde (en priorité enfants, personnes âgées, femmes enceintes) \textbf{à l'air libre}, loin de la source. Ne les laissez pas s'endormir.
    \item \textbf{COUPER LA SOURCE :} Si c'est possible \textbf{sans vous mettre en danger}, éteignez le brasero, coupez le groupe électrogène, fermez la bouteille de gaz. Sinon, fuyez et appelez les secours.
    \item \textbf{ALERTER LES SECOURS :} Appelez les numéros d'urgence camerounais :
    \begin{itemize}
        \item \textbf{118} (SAMU -- Service d'Aide Médicale Urgente)
        \item \textbf{119} (Pompiers -- Sécurité Civile)
        \item \textbf{112} (Numéro d'urgence européen, fonctionne aussi sur GSM au Cameroun)
    \end{itemize}
    Précisez : « \emph{Intoxication probable au monoxyde de carbone, adresse, nombre de victimes, état de conscience.} »
\end{enumerate}
\end{methode}

\begin{exemplebox}[Cas réel type (Fait divers récurrent au Cameroun)]
\textbf{Situation :} En janvier, à Douala, une famille de 5 personnes (parents + 3 enfants) dort dans le salon. Un brasero au charbon de bois est laissé « pour chauffer la pièce » pendant la nuit. La porte et les fenêtres sont fermées hermétiquement (moustiquaires baissées).
\textbf{Déroulement :} Le charbon consume lentement $\rightarrow$ combustion incomplète $\rightarrow$ production massive de $\ce{CO}$. Le gaz remplit la pièce. Les parents se réveillent avec de violents maux de tête, nausées, impossibilité de se lever. Les enfants ne se réveillent pas.
\textbf{Issue :} Un voisin, alerté par le silence inhabituel, force la porte, ouvre tout, sort la famille. Appel au 118. Les parents et l'aîné survivent avec séquelles (troubles mémoire). Les deux plus jeunes décèdent à l'hôpital.
\textbf{Leçon :} \textbf{Jamais de brasero, jamais de groupe électrogène dans un lieu de vie fermé.} La ventilation « par la fente de la porte » est \textbf{insuffisante}.
\end{exemplebox}

\courssub{Prévention : gestes responsables au quotidien}

\begin{popfigure}
\begin{tikzpicture}[scale=0.85, every node/.style={font=\small, transform shape}]
    % Styles
    \tikzset{
        pictogram/.style={circle, draw=popdark, thick, minimum size=2.8cm, fill=popcream},
        cross/.style={draw=popred, line width=4pt, line cap=round},
        check/.style={draw=popgreen, line width=4pt, line cap=round},
        label/.style={align=center, text width=3cm, font=\small\bfseries}
    }
    
    % --- INTERDITS ---
    \node[pictogram] (inter1) at (-5, 0) {};
    \node[label] at (-5, -2) {\textcolor{popred}{INTERDIT}};
    % Brazero inside
    \begin{scope}[shift={(-5,0)}]
        \draw[popdark, thick] (-0.8,-0.6) rectangle (0.8,0.2); % Room floor/wall
        \draw[popdark, thick] (-0.8,-0.6) -- (-0.8,0.8); % Wall
        \draw[popdark, thick] (0.8,-0.6) -- (0.8,0.8); % Wall
        \draw[poporange, thick, fill=poporange!30] (-0.3,-0.4) rectangle (0.3,0.1); % Brazero
        \draw[poporange] (-0.2,-0.4) -- (-0.1,-0.6) (0,-0.4) -- (0.1,-0.6) (0.2,-0.4) -- (0.3,-0.6); % Legs
        \draw[popgray, decorate, decoration={snake, amplitude=1pt, segment length=4pt}] (-0.3,0.1) -- (-0.1,0.6) (0.3,0.1) -- (0.1,0.6); % Smoke
        \node[popred] at (0,0.8) {\Huge $\times$};
    \end{scope}
    
    \node[pictogram] (inter2) at (0, 0) {};
    \node[label] at (0, -2) {\textcolor{popred}{INTERDIT}};
    % Generator inside
    \begin{scope}[shift={(0,0)}]
        \draw[popdark, thick] (-0.8,-0.6) rectangle (0.8,0.2);
        \draw[popdark, thick] (-0.8,-0.6) -- (-0.8,0.8);
        \draw[popdark, thick] (0.8,-0.6) -- (0.8,0.8);
        \draw[popmuted, thick, fill=popmuted!30, rounded corners=2pt] (-0.5,-0.4) rectangle (0.5,0.2); % Gen
        \draw[popgray, ->, thick] (0.5,-0.1) -- (1.0,-0.1); % Exhaust
        \node[popred] at (0,0.8) {\Huge $\times$};
    \end{scope}
    
    \node[pictogram] (inter3) at (5, 0) {};
    \node[label] at (5, -2) {\textcolor{popred}{INTERDIT}};
    % Car in closed garage
    \begin{scope}[shift={(5,0)}]
        \draw[popdark, thick] (-1.0,-0.6) rectangle (1.0,0.8); % Garage
        \draw[popdark, thick] (-0.2,-0.6) -- (-0.2,0.8); % Car shape simple
        \draw[popdark, thick] (0.2,-0.6) -- (0.2,0.8);
        \draw[popdark, thick, fill=popblue!30] (-0.5,-0.5) rectangle (0.5,-0.1); % Car body
        \draw[popgray, decorate, decoration={snake, amplitude=1.5pt, segment length=5pt}] (0,-0.1) -- (0,0.5); % Exhaust
        \node[popred] at (0,0.5) {\Huge $\times$};
    \end{scope}
    
    % --- OBLIGATOIRES / BONNES PRATIQUES ---
    \node[pictogram] (ok1) at (-5, -4) {};
    \node[label] at (-5, -6) {\textcolor{popgreen}{OBLIGATOIRE}};
    \begin{scope}[shift={(-5,-4)}]
        \draw[popdark, thick] (-0.8,-0.6) rectangle (0.8,0.2);
        \draw[popdark, thick] (-0.8,-0.6) -- (-0.8,1.0);
        \draw[popdark, thick] (0.8,-0.6) -- (0.8,1.0);
        \draw[popblue, thick, ->] (-0.8,0.2) -- (-1.3,0.2) node[left] {Fenêtre \textbf{OUVERTE}};
        \draw[popgreen, thick, ->] (0,0.5) -- (0,1.0) node[above] {Air frais};
    \end{scope}
    
    \node[pictogram] (ok2) at (0, -4) {};
    \node[label] at (0, -6) {\textcolor{popgreen}{OBLIGATOIRE}};
    \begin{scope}[shift={(0,-4)}]
        \draw[popdark, thick] (-0.8,-0.6) rectangle (0.8,0.2);
        \draw[popdark, thick] (-0.8,-0.6) -- (-0.8,1.0);
        \draw[popdark, thick] (0.8,-0.6) -- (0.8,1.0);
        \draw[popmuted, thick, fill=popmuted!30, rounded corners=2pt] (1.0,-0.4) rectangle (2.0,0.2); % Gen OUTSIDE
        \draw[popgray, ->, thick] (1.5,-0.1) -- (2.2,-0.1); % Exhaust away
        \node[popgreen] at (1.5,0.8) {\Huge $\checkmark$};
        \node[align=center] at (1.5,0.5) {Groupe \textbf{DEHORS}};
    \end{scope}
    
    \node[pictogram] (ok3) at (5, -4) {};
    \node[label] at (5, -6) {\textcolor{popgreen}{RECOMMANDÉ}};
    \begin{scope}[shift={(5,-4)}]
        \draw[popdark, thick] (-0.8,-0.6) rectangle (0.8,1.0); % Room
        \draw[popblue, thick, fill=popblue!20, rounded corners=2pt] (0.2,0.2) rectangle (0.6,0.6); % Detector
        \node[popblue] at (0.4,0.8) {\Huge $\checkmark$};
        \node[align=center] at (0.4,-0.2) {Détecteur \ce{CO}};
    \end{scope}
\end{tikzpicture}
\legende{Pictogrammes de prévention : en haut, les situations \textbf{interdites} (brasero dedans, groupe électrogène dedans, voiture au ralenti garage fermé). En bas, les \textbf{gestes qui sauvent} : aérer en permanence, placer le groupe électrogène \textbf{strictement à l'extérieur} (loin des fenêtres/portes), installer un détecteur de \ce{CO} (norme EN 50291).}
\end{popfigure}

\begin{aretenir}[Règles d'or de prévention au Cameroun]
\begin{enumerate}
    \item \textbf{JAMAIS} de brasero, fourneau à charbon/bois, réchaud à pétrole \textbf{dans une pièce fermée} (chambre, salon, cuisine sans fenêtre ouverte). Si cuisson au charbon : \textbf{porte/fenêtre grande ouverte} + sortie de la pièce dès que possible.
    \item \textbf{JAMAIS} de groupe électrogène dans la maison, la véranda, le garage, le couloir, sous la fenêtre. \textbf{TOUJOURS À L'EXTÉRIEUR}, à au moins 2--3 mètres des ouvertures, échappement dirigé vers l'extérieur.
    \item \textbf{ENTRETENIR} ses appareils : gazinière (flamme bleue = bon réglage ; flamme jaune/orange = manque d'air = danger $\ce{CO}$), chauffe-eau (ramonage cheminée, vérification tuyau d'évacuation), groupe électrogène (filtre à air, huile, pot d'échappement non bouché).
    \item \textbf{AÉRER} quotidiennement (10 min matin/soir) même par temps frais. Ne pas boucher les grilles de ventilation (souvent faites pour évacuer le $\ce{CO}$ et le $\ce{CO2}$).
    \item \textbf{NE JAMAIS} laisser tourner un moteur (moto, voiture, pompe) dans un espace clos (garage, hangar, cour close).
    \item \textbf{ÉQUIPER} son domicile d'un \textbf{détecteur de monoxyde de carbone} (norme EN 50291, pile longue durée). Il sonne \textbf{avant} que le taux ne devienne mortel. Coût : ~5 000--15 000 FCFA. Une vie n'a pas de prix.
\end{enumerate}
\end{aretenir}

\begin{savaistu}[Le détecteur de CO : un ange gardien électronique]
Dans de nombreux pays (France, Belgique, Canada, USA...), les détecteurs de monoxyde de carbone sont \textbf{obligatoires} dans les logements équipés d'appareils à combustion (chaudière, poêle, garage attenant). Au Cameroun, ils ne sont pas encore obligatoires mais se trouvent en quincaillerie et sur les sites de e-commerce (Jumia, Glotelho, etc.). Choisissez un modèle \textbf{certifié EN 50291}, à affichage numérique (ppm) et alarme sonore > 85 dB. Placez-le à hauteur de respiration (1,5--1,8 m), près des chambres, pas juste au-dessus de la gazinière (fausses alertes). Testez-le chaque mois (bouton Test). Changez la pile chaque année (ou pile 10 ans scellée).
\end{savaistu}

\courssub{Exercice résolu : Analyser une situation à risque}

\begin{exoresolu}[Analyse d'un accident domestique]
\textbf{Énoncé :} Lors d'une coupure \textsc{eneo} de 4 h du matin, M. \textsc{Kamga} démarre son groupe électrogène essence \SI{3}{kVA} placé dans le \textbf{garage attenant à la chambre parentale}, porte du garage fermée, fenêtre du garage fermée. La porte communicante garage/chambre est entrouverte pour passer les rallonges. À 6 h, sa femme se réveille avec de violents maux de tête et des nausées. M. \textsc{Kamga} ne se réveille pas. La femme appelle le 118.

\textbf{Questions :}
\begin{enumerate}
    \item Identifier l'erreur majeure de M. \textsc{Kamga}.
    \item Expliquer, au niveau microscopique (sang), pourquoi la femme a des maux de tête et nausées, et pourquoi le mari est en coma.
    \item Citer les 4 gestes que la femme aurait dû faire immédiatement en se réveillant (si elle en avait la force).
    \item Proposer 2 mesures de prévention pour que cela ne se reproduise plus.
\end{enumerate}

\tcblower
\textbf{Solution détaillée :}
\begin{enumerate}
    \item \textbf{Erreur majeure :} Placer le groupe électrogène \textbf{à l'intérieur du garage} (espace clos/confine) avec les portes/fenêtres fermées. Le garage n'est \textbf{pas} un espace extérieur. Les gaz d'échappement (riches en $\ce{CO}$ par combustion incomplète à bas régime) se sont accumulés dans le garage et ont passé la porte entrouverte pour envahir la chambre.
    \item \textbf{Mécanisme microscopique :}
    \begin{itemize}
        \item L'air de la chambre contenait du $\ce{CO}$. En respirant, le $\ce{CO}$ a traversé les alvéoles pulmonaires et s'est fixé sur l'hémoglobine (\ce{Hb}) des globules rouges pour former de la \ce{HbCO} (carboxyhémoglobine).
        \item L'affinité $\ce{Hb--CO}$ étant $\approx 200\times$ supérieure à $\ce{Hb--O2}$, le $\ce{CO}$ a « volé » la place de l'$\ce{O2}$.
        \item Le sang ne transporte plus assez d'$\ce{O2}$ vers le cerveau $\rightarrow$ \textbf{hypoxie cérébrale}.
        \item \textbf{Femme :} Taux de \ce{HbCO} modéré (ex. 20--30 \%) $\rightarrow$ maux de tête, nausées, vertiges (cerveau souffre mais conscient).
        \item \textbf{Mari :} Taux de \ce{HbCO} élevé (ex. > 50 \%) $\rightarrow$ confusion, perte de conscience, coma (cerveau gravement privé d'$\ce{O2}$). Risque de mort ou séquelles neurologiques définitives.
    \end{itemize}
    \item \textbf{Les 4 réflexes (Méthode) :}
    \begin{enumerate}
        \item \textbf{Aérer :} Ouvrir grande la fenêtre de la chambre ET la porte du garage.
        \item \textbf{Évacuer :} Se sortir elle-même, tenter de sortir le mari (ou le traîner) vers l'air libre (cour, rue).
        \item \textbf{Couper la source :} Couper le groupe électrogène (si accessible sans risque) ou fermer la vanne essence.
        \item \textbf{Alerter :} Appeler le \textbf{118} (SAMU) ou \textbf{119} (Pompiers) en précisant « Intoxication \ce{CO}, groupe électrogène dans garage, 2 victimes dont 1 inconsciente ».
    \end{enumerate}
    \item \textbf{Prévention :}
    \begin{enumerate}
        \item Installer le groupe électrogène \textbf{strictement à l'extérieur} (sur un appui stable, à l'abri de la pluie mais ventilé), à \textbf{> 3 m} de toute porte/fenêtre/prise d'air. Passer les câbles par un trou percé (gaine étanche) ou par une fenêtre entrouverte \textbf{côté opposé} au groupe.
        \item Installer un \textbf{détecteur de \ce{CO}} dans le couloir menant aux chambres et/ou dans la chambre parentale. Vérifier son fonctionnement mensuellement.
    \end{enumerate}
\end{enumerate}
\end{exoresolu}

\courssub{Bilan de la section}

\begin{aretenir}[Essentiel à retenir sur le monoxyde de carbone]
\begin{itemize}
    \item \textbf{Nature :} Gaz \ce{CO}, incolore, inodore, insipide, non irritant = \textbf{indétectable par l'homme}.
    \item \textbf{Origine :} Toute \textbf{combustion incomplète} (brasero, groupe électrogène, gazinière mal réglée, moteur au ralenti en lieu clos).
    \item \textbf{Mécanisme :} Se fixe sur l'hémoglobine (\ce{Hb}) avec une affinité \textbf{200--250 fois} plus forte que l'$\ce{O2}$ $\rightarrow$ formation de \ce{HbCO} (carboxyhémoglobine) $\rightarrow$ \textbf{blocage du transport et de la libération de l'$\ce{O2}$} $\rightarrow$ \textbf{asphyxie cellulaire}.
    \item \textbf{Symptômes :} Maux de tête, nausées, vertiges $\rightarrow$ confusion, coma $\rightarrow$ mort. \textbf{Non spécifiques !}
    \item \textbf{Conduite à tenir :} \textbf{1. Aérer -- 2. Évacuer -- 3. Couper source -- 4. Alerter (118 / 119).}
    \item \textbf{Prévention :} \textbf{Zéro combustion en espace clos.} Groupe électrogène \textbf{dehors}. Aération permanente. Entretien appareils. \textbf{Détecteur \ce{CO} recommandé.}
\end{itemize}
\end{aretenir}

\begin{exercice}[Entraînement -- Prévention et réaction]
\begin{enumerate}
    \item Complétez : Le monoxyde de carbone est un gaz \_\_\_\_\_\_, \_\_\_\_\_\_ et \_\_\_\_\_\_. Il provient des \_\_\_\_\_\_ \_\_\_\_\_\_.
    \item Pourquoi dit-on que le CO a une « affinité 200 fois supérieure » à celle de l'O$_2$ pour l'hémoglobine ? Quelle en est la conséquence directe sur le sang ?
    \item Votre voisin fait tourner son groupe électrogène dans son garage fermé, porte mitoyenne avec votre chambre entrouverte. Vous vous réveillez avec un violent mal de tête. Citez dans l'ordre les 3 premières actions à faire.
    \item Vrai ou Faux ? Justifiez.
    \begin{itemize}
        \item[a)] « La flamme de ma gazinière est bleue, donc il n'y a pas de risque de CO. »
        \item[b)] « J'ai ouvert la fenêtre de 5 cm, le CO peut sortir, je suis en sécurité. »
        \item[c)] « Le CO sent l'essence ou le gaz, on le sentirait s'il y en avait. »
    \end{itemize}
    \item Dessinez un pictogramme « Interdiction » et un pictogramme « Obligation » pour la sécurité vis-à-vis du CO (style panneau de signalisation).
\end{enumerate}
\end{exercice}

\begin{corrige}[Exercice 1]
\begin{enumerate}
    \item Incolore, inodore, sans saveur. Combustions incomplètes.
    \item Cela signifie que la constante d'équilibre de fixation du CO sur l'hémoglobine est 200 fois plus grande. Conséquence : même à très faible concentration dans l'air, le CO « gagne » la place sur l'Hb, formant de l'HbCO. Le sang perd sa capacité à transporter l'O$_2$ (anémie fonctionnelle) et l'O$_2$ restant n'est pas libéré aux tissus.
    \item 1. Aérer (ouvrir grande fenêtre/porte). 2. Évacuer (sortir de la chambre, aller dehors). 3. Couper la source (aller couper le groupe si possible sans danger, sinon fuir et appeler). 4. Alerter (118/119).
    \item 
    \begin{itemize}
        \item[a)] \textbf{Faux.} Flamme bleue = combustion \emph{quasi} complète, mais un défaut d'entretien, un tirage bouché ou un manque d'air ambiant peut faire basculer en combustion incomplète localement. De plus, un appareil peut produire du CO au démarrage/arrêt.
        \item[b)] \textbf{Faux.} Une fente de 5 cm est \textbf{très insuffisante} pour évacuer les volumes de CO produits par un groupe électrogène ou un brasero. Il faut une ventilation traversante massive (porte + fenêtre grandes ouvertes).
        \item[c)] \textbf{Faux.} C'est l'erreur la plus dangereuse. Le CO \textbf{n'a aucune odeur}. L'odeur d'essence/gaz vient d'autres composants (additifs, mercaptans pour le gaz), pas du CO lui-même.
    \end{itemize}
    \item \textit{(Réponse graphique attendue : cercle rouge barré sur brasero/groupe dedans ; cercle vert avec coche sur fenêtre ouverte / groupe dehors / détecteur).}
\end{enumerate}
\end{corrige}

\coursec{Sécurité et gestes responsables autour des combustions}

Après avoir étudié les dangers invisibles du monoxyde de carbone, nous devons maintenant apprendre à maîtriser le feu lui-même. Une combustion est une réaction chimique puissante, utile pour cuisiner, se chauffer ou faire avancer un véhicule, mais elle devient destructrice si on ne la contrôle pas. Cette section te donne les réflexes vitaux pour prévenir les accidents au laboratoire, à la maison et pour comprendre l'impact de nos feux sur la planète.

\courssub{Le triangle du feu : comprendre pour agir}

Toute combustion a besoin de trois éléments réunis simultanément. C'est le \cle{triangle du feu} :
\begin{itemize}
    \item le \textbf{combustible} (ce qui brûle : bois, gaz, essence, huile, papier) ;
    \item le \textbf{comburant} (celui qui fait brûler, presque toujours le dioxygène $\ce{O2}$ de l'air) ;
    \item l'\textbf{énergie d'activation} (étincelle, flamme, surface très chaude, choc mécanique) qui déclenche la réaction.
\end{itemize}
Si l'un des trois côtés manque, la combustion s'arrête immédiatement. C'est le principe fondamental de toute extinction.

\begin{popfigure}
\begin{tikzpicture}[scale=1.1, every node/.style={font=\small}]
    % Triangle vertices
    \coordinate (A) at (0,0); % Combustible
    \coordinate (B) at (5,0); % Comburant
    \coordinate (C) at (2.5,4.33); % Energie
    
    % Triangle edges with labels
    \draw[popdark, thick] (A) -- node[midway, below=2pt] {\textbf{Combustible}} (B);
    \draw[popdark, thick] (B) -- node[midway, right=2pt, rotate=-60, anchor=south] {\textbf{Comburant ($\ce{O2}$)}} (C);
    \draw[popdark, thick] (C) -- node[midway, left=2pt, rotate=60, anchor=south] {\textbf{Énergie d'activation}} (A);
    
    % Center "FEU"
    \node[fill=poporange!80, text=white, rounded corners=4pt, inner sep=4pt, font=\bfseries\large] at (2.5, 1.8) {FEU};
    
    % Extinction methods outside triangle
    % Remove Combustible
    \node[align=center, popblue, font=\bfseries\small] at (-1.5, 0) {Retirer le\\combustible\\(couper gaz,\\éloigner bois)};
    \draw[popblue, ->, thick] (-1, 0) -- (-0.2, 0);
    
    % Remove Comburant
    \node[align=center, popgreen, font=\bfseries\small] at (6.5, 0) {Supprimer le\\comburant\\(étouffer :\\couvercle, sable,\\couverture, CO2)};
    \draw[popgreen, ->, thick] (5.8, 0) -- (5.2, 0);
    
    % Remove Energy
    \node[align=center, poppurple, font=\bfseries\small] at (2.5, 5.5) {Refroidir /\\Baisser l'énergie\\(eau sur feu\\de bois/classe A)};
    \draw[poppurple, ->, thick] (2.5, 4.8) -- (2.5, 4.4);
\end{tikzpicture}
\legende{Le triangle du feu et les trois stratégies d'extinction associées.}
\end{popfigure}

\begin{methode}[Éteindre un feu naissant en agissant sur le triangle]
\emph{Objectif :} Choisir le bon geste selon la nature du feu.
\begin{enumerate}
    \item \textbf{Identifier le combustible} (solide : bois, papier / liquide : essence, huile / gaz : butane / métal / électrique).
    \item \textbf{Choisir l'action sur le triangle} :
    \begin{itemize}
        \item \textbf{Supprimer le comburant (étouffer) :} couvercle sur une casserole, couverture anti-feu, sable, extincteur à $\ce{CO2}$ ou poudre. \emph{Idéal pour feux de liquides (friteuse) et gaz.}
        \item \textbf{Refroidir (baisser l'énergie) :} eau en grande quantité. \emph{Uniquement sur feux de solides (bois, papier, tissu) — Classe A.}
        \item \textbf{Supprimer le combustible :} fermer le robinet de gaz, couper l'électricité, éloigner les objets inflammables.
    \end{itemize}
    \item \textbf{Ne jamais tourner le dos à un feu éteint} : risque de reprise.
    \item \textbf{Appeler les secours (18 ou 112)} si le feu dépasse la taille d'une poubelle ou s'il y a de la fumée épaisse.
\end{enumerate}
\end{methode}

\courssub{Consignes de sécurité au laboratoire}

Au collège, nous manipulons des produits chimiques et des sources de chaleur. La rigueur protège tout le monde.

\begin{itemize}
    \item \textbf{Équipement de protection individuelle (EPI) :} \cle{lunettes de protection} obligatoires (pas de lunettes de vue simples), \cle{blouse} en coton boutonnée, cheveux longs attachés, chaussures fermées (pas de sandales).
    \item \textbf{Le bec Bunsen :} on l'allume \emph{après} avoir ouvert le gaz, allumette ou briquet déjà en main. La \cle{flamme de chauffe} (bleue, non lumineuse, silencieuse, deux cônes nets) est la seule autorisée pour chauffer. La flamme de sécurité (jaune, lumineuse, instable) ne sert qu'à l'allumage ou à l'attente.
    \item \textbf{Produits inflammables :} \textbf{jamais} de flamme nue près de l'éthanol, de l'acétone, de l'éther, de l'essence. Utiliser un bain-marie électrique ou une plaque chauffante pour les chauffer.
    \item \textbf{Extinction correcte :} fermer le robinet de gaz \textbf{avant} de fermer l'arrivée d'air du bec (sinon flamme remontante dans le tuyau).
    \item \textbf{Verre chaud :} un verre chaud ne se distingue pas d'un verre froid. Poser toujours sur un support en grillage ou en bois, jamais directement sur la paillasse froide (risque d'éclatement) et ne pas le toucher à mains nues.
\end{itemize}

\begin{attention}[Solvants et flamme : le piège mortel]
Les vapeurs d'éthanol, d'acétone ou d'essence sont plus lourdes que l'air : elles rampent sur la paillasse et peuvent s'enflammer à distance (retour de flamme) si un bec Bunsen est allumé à 2 mètres. \textbf{Règle d'or :} « Pas de flamme nue si un flacon de solvant est ouvert dans la pièce. » Utilise un bain-marie ou une plaque chauffante.
\end{attention}

\courssub{Sécurité domestique : gaz, bois et électricité}

Au Cameroun, la cuisine est le lieu le plus à risque. Le gaz butane (bonbonnes bleues/rouges), le bois de feu et le charbon de bois sont nos principaux combustibles.

\begin{itemize}
    \item \textbf{Stockage des bouteilles de gaz :} toujours à \textbf{l'air libre} (balcon, cour) ou dans un local \textbf{bien ventilé} (grilles hautes et basses). \textbf{Jamais} dans une cave, un placard fermé ou sous l'évier : le butane est plus lourd que l'air, il s'accumule au sol et forme un mélange explosif.
    \item \textbf{Tuyau et détendeur :} vérifier la date de péremption sur le tuyau (généralement 5 ans). Le détendeur doit être adapté (basse pression 28-30 mbar pour le butane).
    \item \textbf{Test d'étanchéité (eau savonneuse) :} après chaque changement de bouteille ou si tu sens l'odeur (mercaptan ajouté exprès), badigeonne les raccords avec de l'eau savonneuse. Des \textbf{bulles} qui grossissent = fuite. \textbf{Jamais} tester avec une allumette !
    \item \textbf{Feu de bois / charbon :} installer le foyer loin des murs, des nattes, des bidons d'huile ou de carburant. Surélever le foyer si possible. \textbf{Éloigner les enfants} : les brûlures par chute dans le feu sont fréquentes. Éteindre complètement (arroser les braises) avant de dormir ou de sortir.
    \item \textbf{Électricité :} ne pas surcharger les rallonges (multiprises « pieuvre »). Un fil qui chauffe signale un danger. Débrancher le fer à repasser, la bouilloire après usage.
\end{itemize}

\begin{exemplebox}[Vérifier l'étanchéité d'une bouteille de gaz avec de l'eau savonneuse]
\textbf{Situation :} Tu viens de raccorder une nouvelle bouteille de butane à ta cuisinière.
\textbf{Manipulation :}
\begin{enumerate}
    \item Ouvre le robinet de la bouteille (détendeur fermé).
    \item Avec un pinceau ou une éponge, applique de l'eau savonneuse (liquide vaisselle + eau) sur le raccord bouteille/détendeur et sur le raccord détendeur/tuyau.
    \item Observe pendant 30 secondes.
\end{enumerate}
\textbf{Interprétation :}
\begin{itemize}
    \item \textbf{Pas de bulles :} l'installation est étanche. Tu peux ouvrir le détendeur et allumer.
    \item \textbf{Bulles qui se forment et grossissent :} il y a une fuite. \textbf{Ferme immédiatement le robinet de la bouteille.} Aère la pièce (ouvre grandes les portes et fenêtres). Ne touche à aucun interrupteur électrique (étincelle). Resserre le raccord ou change le joint/tuyau. Refais le test.
\end{itemize}
\end{exemplebox}

\begin{popfigure}
\begin{tikzpicture}[scale=0.9, every node/.style={font=\small}]
    % Kitchen outline
    \draw[popdark, thick] (0,0) rectangle (10,6);
    % Window
    \draw[popblue, thick, fill=popblueL!30] (7.5, 3.5) rectangle (9.5, 5.5);
    \node[popblue, align=center] at (8.5, 4.5) {Fenêtre\\OUVERTE};
    % Gas bottle outside/ventilated
    \draw[poporange, thick, fill=poporangeL!30] (0.5, 0.5) rectangle (2.5, 2.5);
    \node[poporange, align=center] at (1.5, 1.5) {Bouteille GAZ\\(extérieur/ventilé)};
    \draw[poporange, ->, thick] (2.5, 1.5) -- (3.5, 1.5);
    % Stove
    \draw[popdark, fill=popmuted!20] (4, 0.5) rectangle (7, 2);
    \node[popdark] at (5.5, 1.25) {Cuisinière};
    % Extinguisher
    \draw[poporange, thick, fill=poporange!10] (8, 0.5) rectangle (9.2, 2);
    \node[poporange, align=center, rotate=90] at (8.6, 1.25) {EXTINCTEUR\\Poudre ABC};
    % Soap water bottle near bottle
    \draw[popgreen, fill=popgreenL!30] (0.5, 2.8) rectangle (2, 3.5);
    \node[popgreen, align=center] at (1.25, 3.15) {Eau savonneuse\\pour test};
    % Child safety zone
    \draw[poppink, dashed, thick] (3, 2.5) rectangle (7, 4.5);
    \node[poppink, align=center] at (5, 3.5) {Zone interdite\\aux enfants};
    % Ventilation grilles
    \draw[popmuted, thick] (0, 5.5) -- (10, 5.5);
    \foreach \x in {1,3,5,7,9} \draw[popmuted, thick] (\x, 5.5) -- (\x, 5.8);
    \node[popmuted, above] at (5, 5.8) {Grilles de ventilation hautes};
    \foreach \x in {1,3,5,7,9} \draw[popmuted, thick] (\x, 0) -- (\x, -0.3);
    \node[popmuted, below] at (5, -0.3) {Grilles de ventilation basses};
\end{tikzpicture}
\legende{Exemple de cuisine sécurisée : bouteille ventilée, fenêtre ouverte, extincteur accessible, zone enfants délimitée, grilles de ventilation.}
\end{popfigure}

\courssub{Moyens d'extinction adaptés à chaque type de feu}

Tous les feux ne s'éteignent pas de la même façon. Utiliser le mauvais agent aggrave l'incendie.

\begin{center}
\begin{tabularx}{\linewidth}{|l|X|X|X|}\hline
\rowcolor{popblueL}
\textbf{Classe} & \textbf{Combustibles typiques} & \textbf{Agent extincteur EFFICACE} & \textbf{INTERDIT / DANGEREUX} \\ \hline
\textbf{A} (Solides) & Bois, papier, tissu, plastique, pneus & \textbf{Eau} (refroidit), Poudre ABC, Mousse & -- \\ \hline
\textbf{B} (Liquides) & Essence, huile, alcool, solvants, graisse de friture & \textbf{Étouffement :} Couvercle, Couverture anti-feu, Sable, Poudre ABC, $\ce{CO2}$, Mousse & \textbf{EAU} (projection brûlante, propagation) \\ \hline
\textbf{C} (Gaz) & Butane, propane, méthane, acétylène & \textbf{Couper l'arrivée de gaz} (robinet), Poudre ABC, $\ce{CO2}$ & Eau (inutile sur le gaz lui-même) \\ \hline
\textbf{E} (Électrique) & Tableau, moteur, prise, rallonge & \textbf{Couper le courant}, Poudre ABC, $\ce{CO2}$ & \textbf{EAU} (conduction = électrocution) \\ \hline
\textbf{F} (Huiles/Grasses cuisson) & Friteuse, huile de cuisson en feu & \textbf{Couvercle}, Couverture anti-feu, Extincteur spécial Classe F & \textbf{EAU} (explosion de vapeur), Poudre standard (reprise) \\ \hline
\end{tabularx}
\end{center}

\begin{attention}[Jamais d'eau sur une friteuse ou de l'huile en feu !]
L'huile brûle à plus de \SI{300}{\celsius}. L'eau, plus dense, tombe au fond, se vaporise instantanément (1 L d'eau $\to$ \SI{1600}{L} de vapeur). Cette explosion de vapeur projette l'huile enflammée en boule de feu (phénomène de \emph{boilover}). \textbf{Geste vital :} Pose un \textbf{couvercle} hermétique (ou une grande plaque, une couverture anti-feu) pour étouffer. Coupe le gaz/l'électricité. Laisse refroidir \textbf{au moins 30 min} avant d'ouvrir. Si pas de couvercle : extincteur poudre ABC ou $\ce{CO2}$ à distance, ou sable. \textbf{Jamais d'eau.}
\end{attention}

\courssub{Impact environnemental : combustions et climat}

Chaque combustion de carbone (bois, charbon, essence, gaz, gasoil, déchets plastiques) rejette du dioxyde de carbone $\ce{CO2}$.
\[ \ce{C + O2 -> CO2} \quad \text{ou} \quad \ce{C_xH_y + O2 -> CO2 + H2O} \]

\begin{propriete}[Le $\ce{CO2}$, gaz à effet de serre]
Le dioxyde de carbone est transparent à la lumière visible du Soleil mais absorbe le rayonnement infrarouge émis par la Terre chauffée. Il piège la chaleur dans l'atmosphère : c'est l'\textbf{effet de serre}. L'augmentation de sa concentration (due aux activités humaines : transport, industrie, déforestation, chauffage) entraîne le \textbf{réchauffement climatique global}.
\end{propriete}

Au Cameroun, l'impact est double :
\begin{itemize}
    \item \textbf{Émissions locales :} les millions de motos-taxis, les groupes électrogènes, les voitures anciennes, la cuisson au bois/charbon rejettent $\ce{CO2}$ et polluants (particules, $\ce{NOx}$, $\ce{CO}$).
    \item \textbf{Déforestation pour le charbon de bois :} pour produire 1 kg de charbon, il faut environ 5 à 10 kg de bois vert. La demande urbaine massive (Yaoundé, Douala, Garoua) détruit les savanes et forêts du Nord et de l'Est. Moins d'arbres = moins de $\ce{CO2}$ capté + érosion des sols + perte de biodiversité.
\end{itemize}

\begin{savaistu}[Le charbon de bois et la déforestation au Nord-Cameroun]
Dans les régions de l'Adamaoua, du Nord et de l'Extrême-Nord, la production artisanale de charbon de bois (carbonisation en meules) est une source de revenu majeure mais une catastrophe écologique. Les arbres (acacias, combretums) sont coupés plus vite qu'ils ne repoussent. Le sol nu subit l'érosion éolienne et hydrique. Des programmes de reboisement et de foyers améliorés (foyers « améliorés » qui consomment 2 à 3 fois moins de bois/charbon) sont développés par le MINEPDED et des ONG pour casser ce cercle vicieux : pauvreté $\to$ coupe bois $\to$ désertification $\to$ pauvreté accrue.
\end{savaistu}

\courssub{Bilan de la leçon : de la réaction à la prévention}

Nous avons construit une chaîne logique rigoureuse tout au long de cette leçon :

\begin{enumerate}
    \item \textbf{Observation :} Une transformation chimique modifie la matière (couleur, gaz, chaleur, précipité). Les atomes se réarrangent.
    \item \textbf{Loi de conservation (Lavoisier) :} La masse totale ne change pas. Les atomes sont conservés.
    \item \textbf{Modélisation :} L'\textbf{équation-bilan} traduit ce réarrangement. L'\textbf{équilibrage} (coefficients stœchiométriques) garantit le respect de la loi de Lavoisier atome par atome.
    \item \textbf{Application majeure :} La \textbf{combustion} (réaction rapide $\text{Combustible} + \ce{O2} \rightarrow \text{Produits} + \text{Énergie}$).
    \item \textbf{Analyse des produits :} Combustion complète ($\ce{CO2}$, $\ce{H2O}$) vs incomplète ($\ce{CO}$, $\ce{C}$, $\ce{H2O}$). Le \textbf{monoxyde de carbone $\ce{CO}$} est un tueur silencieux (affinité hémoglobine 200-250$\times$ $\ce{O2}$).
    \item \textbf{Prévention :} Comprendre le \textbf{triangle du feu} permet d'éteindre (supprimer un côté). Respecter les \textbf{consignes} (laboratoire, maison : gaz, bois, électricité). Choisir le \textbf{bon extincteur} (jamais d'eau sur huile/électricité).
    \item \textbf{Responsabilité citoyenne :} Limiter les combustions inutiles, entretenir son véhicule, utiliser des foyers améliorés, planter des arbres $\to$ lutter contre le réchauffement et la déforestation.
\end{enumerate}

\begin{exoresolu}[Scénario : Incendie de friteuse à la maison]
\textbf{Énoncé :} En faisant frire du poisson, l'huile de ta friteuse prend feu. La flamme monte à 30 cm. Tu as à portée de main : une carafe d'eau, le couvercle de la friteuse, une couverture en laine, un extincteur à poudre ABC (vérifié il y a 6 mois), le robinet de gaz de la cuisinière.
\textbf{Question :} Dans quel ordre exact vas-tu agir ? Justifie chaque geste par le triangle du feu.

\tcblower

\textbf{Solution détaillée :}
\begin{enumerate}
    \item \textbf{Coupe le gaz (robinet cuisinière).} \emph{Justification :} Supprime le \textbf{combustible} (gaz) qui alimente la flamme sous la friteuse. Priorité absolue pour ne pas aggraver le feu.
    \item \textbf{Pose le couvercle sur la friteuse (ou la couverture en laine si pas de couvercle).} \emph{Justification :} Supprime le \textbf{comburant} ($\ce{O2}$) au contact de l'huile en feu. Étouffe le feu de classe B/F. \textbf{Ne verse JAMAIS l'eau} (risque d'explosion \emph{boilover}).
    \item \textbf{Laisse le couvercle en place au moins 30 minutes.} \emph{Justification :} L'huile reste à température d'auto-inflammation. L'ouvrir trop tôt réintroduit le \textbf{comburant} $\to$ reprise immédiate.
    \item \textbf{Si le feu s'étend au meuble/rideau (classe A) :} utilise l'extincteur poudre ABC (efficace sur A, B, C) en visant la base des flammes. \emph{Justification :} La poudre isole le combustible du comburant et inhibe la réaction en chaîne (action chimique).
    \item \textbf{Appelle les secours (18/112)} si le feu a dépassé la friteuse.
\end{enumerate}
\textbf{Bilan :} Tu as agi sur les trois côtés du triangle (gaz = combustible, couvercle = comburant, temps = énergie) sans jamais ajouter d'eau.
\end{exoresolu}

\begin{exercice}[Entraînement : Réflexes sécurité]
\begin{enumerate}
    \item Au laboratoire, tu dois chauffer \SI{20}{mL} d'éthanol pour une distillation. Quelle source de chaleur utilises-tu ? Pourquoi le bec Bunsen est-il interdit ?
    \item Tu sens une forte odeur de gaz en entrant dans la cuisine. La bouteille est dans le placard sous l'évier. Liste les 4 actions immédiates dans l'ordre, en expliquant le lien avec le triangle du feu ou la prévention $\ce{CO}$/explosion.
    \item Pourquoi la production massive de charbon de bois au Nord-Cameroun aggrave-t-elle le réchauffement climatique de deux façons différentes ?
    \item Complète : Un extincteur à $\ce{CO2}$ éteint en supprimant le côté \underline{\hspace{3cm}} du triangle. Il est idéal pour les feux de classe \underline{\hspace{2cm}} et \underline{\hspace{2cm}} car il ne laisse \underline{\hspace{3cm}}.
\end{enumerate}
\end{exercice}

\begin{corrige}[Exercice n° Entraînement : Réflexes sécurité]
\begin{enumerate}
    \item \textbf{Bain-marie électrique ou plaque chauffante.} L'éthanol est très inflammable (vapeurs lourdes). La flamme nue du bec Bunsen provoquerait un retour de flamme ou l'inflammation des vapeurs rampantes.
    \item \textbf{1. Ouvre grandes portes et fenêtres (ventilation) $\to$ dilue le comburant/mélange explosif.} \textbf{2. Ferme le robinet de la bouteille $\to$ supprime le combustible.} \textbf{3. Ne touche à AUCUN interrupteur/prise/téléphone $\to$ évite l'énergie d'activation (étincelle).} \textbf{4. Sors, appelle les secours de l'extérieur.} (Ne pas chercher la fuite avec une flamme !).
    \item \textbf{1. Émission directe :} La combustion du charbon (cuisson) rejette $\ce{CO2}$ (gaz à effet de serre). \textbf{2. Puits de carbone détruit :} La coupe des arbres pour faire le charbon supprime la photosynthèse qui captait le $\ce{CO2}$. Double peine.
    \item \textbf{Comburant} (il chasse l'$\ce{O2}$ par dilution/refroidissement) ; \textbf{B} (liquides) et \textbf{E} (électrique) ; \textbf{aucun résidu} (gaz qui se disperse).
\end{enumerate}
\end{corrige}

\newpage

\coursec{Activité d'intégration}

\coursec{Leçon 4 : La réaction chimique — Enquête à Nkolbisson}

\courssub{Objectifs de l'activité}
Cette activité d'intégration vise à mobiliser l'ensemble des compétences acquises au cours de la leçon « La réaction chimique » à travers l'analyse d'un fait divers réel et tragique. À l'issue de ce travail, tu seras capable de :
\begin{itemize}
    \item Identifier une transformation chimique (combustion) à partir d'observations macroscopiques et de résultats expérimentaux.
    \item Appliquer la loi de conservation de la masse (loi de Lavoisier) pour vérifier qu'aucun atome n'a disparu lors de la réaction.
    \item Écrire et équilibrer les équations-bilan des deux types de combustion du carbone (complète et incomplète).
    \item Expliquer l'influence de la quantité de dioxygène disponible sur le type de combustion.
    \item Relier la formation de monoxyde de carbone (CO) à son mécanisme toxique (fixation sur l'hémoglobine) et aux symptômes cliniques.
    \item Concevoir un message de prévention clair, illustré et adapté au contexte camerounais pour lutter contre les intoxications domestiques au CO.
\end{itemize}

\courssub{Situation : Le drame du brasero de Nkolbisson}
\emph{Yaoundé, quartier Nkolbisson, nuit du 12 au 13 janvier 2024.}

Un article du journal \emph{Mutations} relate le drame suivant :
\begin{quote}
\textit{« Une famille de quatre personnes (les deux parents et leurs deux enfants âgés de 3 et 6 ans) a été retrouvée inanimée dans leur chambre unique de 12 m², close hermétiquement pour se protéger du froid et des moustiques. La source de chaleur : un brasero métallique (diamètre 30 cm) contenant environ \SI{2,0}{kg} de charbon de bois en combustion lente. Le père, rentré tardivement, a découvert sa femme et ses enfants inconscients. Lui-même a ressenti de violents maux de tête, des nausées et une confusion mentale avant d'ouvrir la porte en grand et d'appeler les secours. Les pompiers, arrivés rapidement, ont mesuré un taux de monoxyde de carbone (CO) dans l'air de la chambre supérieur à \num{500}~ppm (parties par million). La mère et le plus jeune enfant, dans un état critique, ont été évacués vers l'Hôpital Central de Yaoundé pour une oxygénothérapie hyperbare. Le diagnostic médical confirme une intoxication aiguë au monoxyde de carbone. »}
\end{quote}

Pour comprendre scientifiquement ce drame et en tirer des leçons, ton groupe dispose du dossier documentaire suivant :

\begin{experience}[Simulation de la combustion en milieu confiné (Document 1)]
Un élève de 3ème a réalisé l'expérience suivante sous la surveillance de son professeur :
\begin{enumerate}
    \item Il place \SI{10,0}{g} de charbon de bois (constitué quasi exclusivement de carbone C) dans une coupelle en porcelaine.
    \item Il met la coupelle sous une cloche en verre de volume \SI{2,0e-3}{\cubic\meter} (soit \SI{2,0}{L}), posée sur un plateau muni d'un robinet. La cloche contient initialement de l'air atmosphérique (environ \SI{21}{\percent} de dioxygène O₂ en volume).
    \item Il ferme le système, allume le charbon à l'aide d'un fil chauffant (sans apport d'air extérieur), et laisse la combustion se faire jusqu'à extinction complète de la braise (environ 45 minutes).
    \item Il laisse refroidir à température ambiante (20 °C), puis pèse l'ensemble \textbf{cloche + plateau + coupelle + cendres + gaz}. La masse totale est \textbf{inchangée} par rapport à la masse initiale.
    \item Il ouvre le robinet et fait bubbler les gaz de la cloche dans un tube à essai contenant \SI{20}{mL} d'eau de chaux (solution saturée d'hydroxyde de calcium Ca(OH)₂). L'eau de chaux devient \textbf{trouble} (blanchiment).
    \item Il effectue un test colorimétrique spécifique au monoxyde de carbone (tube détecteur à l'iodure de palladium(II) ou papier au chlorure de palladium) sur un échantillon des gaz : le réactif vire au \textbf{noir/brun foncé}, confirmant la présence de CO.
\end{enumerate}
\end{experience}

\paragraph{Document 2 : Données sur l'atmosphère de la chambre.}
\begin{itemize}
    \item Volume de la chambre : \SI{12}{\square\meter} au sol $\times$ \SI{2,5}{m} de hauteur = \SI{30}{\cubic\meter} d'air.
    \item Masse volumique de l'air : \SI{1,2}{kg/m^3}. Masse d'air totale dans la chambre $\approx$ \SI{36}{kg}.
    \item Masse de dioxygène (O₂) initialement disponible : \SI{21}{\percent} de \SI{36}{kg} $\approx$ \SI{7,6}{kg} d'O₂.
    \item Charbon brûlé estimé (brasero) : \SI{2,0}{kg} de carbone (C).
    \item Mesure du CO dans l'air par les pompiers : \textbf{\num{500}~ppm} en volume (soit \SI{500}{L} de CO pour \SI{1\,000\,000}{L} d'air, soit \num{0,05}\%).
\end{itemize}

\paragraph{Document 3 : Rappel physiologique (extrait manuel SVT).}
L'hémoglobine (Hb) est une protéine du sang qui transporte le dioxygène (O₂) des poumons vers les cellules. Le monoxyde de carbone (CO) se fixe sur l'hémoglobine à la place de l'O₂ pour former de la \textbf{carboxyhémoglobine} (HbCO). L'affinité de l'hémoglobine pour le CO est \textbf{200 à 250 fois plus grande} que pour l'O₂. Même à faible concentration dans l'air inhalé, le CO « vole » la place de l'O₂ sur l'hémoglobine, empêchant l'oxygénation des organes (cerveau, cœur). Les symptômes apparaissent dès \SI{10}{\percent} d'HbCO (maux de tête), s'aggravent à \SI{30}{\percent} (nausées, confusion) et deviennent mortels au-delà de \SI{50}{\percent}.

\courssub{Consignes}
Travaillant par groupe de 3 ou 4, réponds aux questions suivantes en t'appuyant sur les documents et sur tes connaissances du cours. Rédige tes réponses de manière structurée.

\begin{enumerate}
    \item \textbf{Identification de la réaction.}
    \begin{enumerate}
        \item D'après le Document 1, quelles observations macroscopiques prouvent qu'une transformation chimique a eu lieu sous la cloche ?
        \item Nomme les réactifs initiaux et les produits finaux détectés. Quelle est la nature chimique du charbon utilisé ?
    \end{enumerate}
    
    \item \textbf{Conservation de la masse et équations-bilan.}
    \begin{enumerate}
        \item L'expérience montre que la masse du système fermé ne change pas. Quelle loi fondamentale cela illustre-t-il ? Énonce-la précisément.
        \item Écris l'équation-bilan de la \textbf{combustion complète} du carbone (formation exclusive de CO₂). Équilibre-la.
        \item Écris l'équation-bilan de la \textbf{combustion incomplète} du carbone (formation exclusive de CO). Équilibre-la.
        \item L'eau de chaux se trouble (Document 1, étape 5) et le test au CO est positif (étape 6). Que concluis-tu sur le \textbf{type de combustion} qui a réellement eu lieu sous la cloche ? Justifie par les deux équations.
    \end{enumerate}
    
    \item \textbf{Analyse du drame de Nkolbisson (quantitatif).}
    \begin{enumerate}
        \item Calcule la masse de dioxygène (O₂) nécessaire pour brûler \textbf{complètement} les \SI{2,0}{kg} de carbone du brasero (équation de la question 2.b). Exprime le résultat en kg.
        \item Calcule la masse de dioxygène (O₂) nécessaire pour brûler \textbf{incomplètement} les \SI{2,0}{kg} de carbone (équation de la question 2.c). Exprime le résultat en kg.
        \item Compare ces deux valeurs à la masse d'O₂ réellement disponible dans la chambre (\SI{7,6}{kg}, Document 2). Quel type de combustion la quantité d'air disponible a-t-elle favorisé ? Explique pourquoi la combustion a été \emph{forcément} incomplète.
        \item Vérification d'ordre de grandeur : La mesure des pompiers donne \num{500}~ppm de CO dans \SI{30}{\cubic\meter} d'air. Calcule le volume de CO formé (en litres). En supposant que tout ce CO provient de la combustion incomplète de 2,0 kg de C, quel pourcentage du carbone a-t-il brûlé en CO ? (Donnée : au cours, on admet que 1 mol de gaz occupe environ \SI{24}{L} à 20 °C ; M(CO) = \SI{28}{g/mol}).
    \end{enumerate}
    
    \item \textbf{Mécanisme toxique et symptômes.}
    \begin{enumerate}
        \item À l'aide du Document 3, explique pourquoi le père a ressenti des maux de tête et des nausées alors qu'il n'est resté que peu de temps dans la pièce.
        \item Pourquoi l'ouverture brutale de la porte a-t-elle permis sa survie, alors que sa femme et son jeune enfant, restés plus longtemps, ont nécessité une oxygénothérapie hyperbare ?
    \end{enumerate}
    
    \item \textbf{Production : Affiche de prévention « Protégez votre famille du tueur invisible ».}
    Conçois, sur une feuille A4 (ou format diapositive), une affiche destinée aux familles des quartiers populaires de Yaoundé (Nkolbisson, Mvog-Ada, Essos, etc.). Elle doit contenir :
    \begin{itemize}
        \item Un titre accrocheur.
        \item Au moins \textbf{quatre consignes de sécurité illustrées} (pictogrammes ou dessins simples) : aération, interdiction du brasero en chambre close, entretien des conduits, détecteur CO.
        \item L'explication simple du danger : « Le charbon + manque d'air = CO (gaz mortel, sans odeur, sans couleur) ».
        \item Les numéros d'urgence au Cameroun : \textbf{112} (pompiers/SAMU) et \textbf{119} (police).
    \end{itemize}
\end{enumerate}

\courssub{Ressources à mobiliser}
\begin{aretenir}[Formules et notions clés de la leçon]
    \begin{itemize}
        \item \textbf{Loi de Lavoisier :} « Rien ne se perd, rien ne se crée, tout se transforme. » Lors d'une réaction chimique, la masse totale des réactifs = masse totale des produits. Les atomes se conservent.
        \item \textbf{Équilibrage d'une équation-bilan :} On place les formules des réactifs (à gauche) et des produits (à droite) séparés par une flèche $\rightarrow$. On ajoute des \textbf{coefficients stœchiométriques} (chiffres devant les formules) pour avoir le même nombre d'atomes de chaque élément des deux côtés de la flèche. \textbf{On ne modifie jamais les formules chimiques.}
        \item \textbf{Combustion complète du carbone :} $\ce{C + O2 -> CO2}$ (nécessite beaucoup d'O₂).
        \item \textbf{Combustion incomplète du carbone :} $\ce{2C + O2 -> 2CO}$ (se produit en manque d'O₂).
        \item \textbf{Masse molaire :} $M(\ce{C}) = \SI{12}{g/mol}$ ; $M(\ce{O2}) = \SI{32}{g/mol}$ ; $M(\ce{CO2}) = \SI{44}{g/mol}$ ; $M(\ce{CO}) = \SI{28}{g/mol}$.
        \item \textbf{Lien masse -- quantité de matière :} $n = \frac{m}{M}$.
        \item \textbf{Toxicité du CO :} Affinité Hb/CO $\approx$ 250 $\times$ Affinité Hb/O₂. Formation de carboxyhémoglobine (HbCO) $\rightarrow$ hypoxie tissulaire.
    \end{itemize}
\end{aretenir}

\begin{methode}[Équilibrer une équation-bilan (rappel)]
\begin{enumerate}
    \item Écrire les formules des réactifs et des produits (ébauche).
    \item Faire le bilan des atomes pour chaque élément des deux côtés de la flèche.
    \item Ajuster les \textbf{coefficients stœchiométriques} (chiffres entiers devant les formules) en commençant par l'élément le plus complexe (souvent C, puis H, puis O).
    \item Vérifier le bilan final : même nombre d'atomes de chaque élément à gauche et à droite.
    \item \textbf{Interdiction formelle :} Ne jamais changer les indices dans les formules chimiques (ex : passer de O₂ à O est interdit).
\end{enumerate}
\end{methode}

\courssub{Démarche et corrigé commenté}
\begin{exoresolu}[Corrigé détaillé de l'enquête de Nkolbisson]
\textbf{1. Identification de la réaction.}
\begin{enumerate}
    \item[a)] \textbf{Observations prouvant la transformation chimique :}
    \begin{itemize}
        \item Disparition de la braise rouge (charbon) $\rightarrow$ formation de cendres (résidu solide différent).
        \item Troublissement de l'eau de chaux $\rightarrow$ preuve de formation de dioxyde de carbone ($\ce{CO2}$).
        \item Test colorimétrique positif au CO $\rightarrow$ preuve de formation de monoxyde de carbone ($\ce{CO}$).
        \item Échauffement de la cloche (réaction exothermique).
    \end{itemize}
    \item[b)] \textbf{Réactifs :} Carbone (C, constituant principal du charbon) et Dioxygène (O₂, de l'air enfermé dans la cloche).\\
    \textbf{Produits détectés :} Dioxyde de carbone ($\ce{CO2}$) et Monoxyde de carbone ($\ce{CO}$). Il reste probablement du carbone non réagi (cendres) et du diazote (N₂, inerte).
\end{enumerate}

\begin{popfigure}
\begin{tikzpicture}[node distance=1.5cm, >=Stealth, font=\small]
\node[draw, rounded corners, fill=popblueL, minimum width=2.5cm, minimum height=1cm] (react) {C (s) + O₂ (g)};
\node[draw, rounded corners, fill=popgreenL, minimum width=2.5cm, minimum height=1cm, right=of react, align=center] (complete){Combustion complète\\Excès d'O₂};
\node[draw, rounded corners, fill=poporangeL, minimum width=2.5cm, minimum height=1cm, below=of complete, align=center] (incomplete){Combustion incomplète\\Manque d'O₂};
\node[draw, rounded corners, fill=poppurpleL, minimum width=2cm, minimum height=1cm, right=of complete] (prodC) {CO₂ (g)};
\node[draw, rounded corners, fill=poppinkL, minimum width=2cm, minimum height=1cm, right=of incomplete] (prodI) {CO (g)};
\draw[->, thick, popgreen] (react) -- (complete);
\draw[->, thick, poporange] (react) -- (incomplete);
\draw[->, thick, popgreen] (complete) -- (prodC);
\draw[->, thick, poporange] (incomplete) -- (prodI);
\end{tikzpicture}
\legende{Les deux voies de la combustion du carbone selon la disponibilité en dioxygène.}
\end{popfigure}

\textbf{2. Conservation de la masse et équations-bilan.}
\begin{enumerate}
    \item[a)] \textbf{Loi de Lavoisier (conservation de la masse) :} « Lors d'une transformation chimique, la masse totale des réactifs est égale à la masse totale des produits. » La masse du système fermé (cloche + contenu) reste constante car aucun atome ne sort ni n'entre ; ils se réorganisent seulement.
    \item[b)] \textbf{Combustion complète :}
    \begin{itemize}
        \item Ébauche : $\ce{C + O2 -> CO2}$
        \item Bilan atomes : C : 1=1 ; O : 2=2. \textbf{Équilibrée.}
        \item \cle{Équation finale :} $\ce{C + O2 -> CO2}$
    \end{itemize}
    \item[c)] \textbf{Combustion incomplète :}
    \begin{itemize}
        \item Ébauche : $\ce{C + O2 -> CO}$ (pas équilibrée en O : 2 à gauche, 1 à droite).
        \item Pour équilibrer l'oxygène, on met 2 devant CO : $\ce{C + O2 -> 2CO}$.
        \item Bilan C : 1 à gauche, 2 à droite. On met 2 devant C : $\ce{2C + O2 -> 2CO}$.
        \item Vérification : C : 2=2 ; O : 2=2. \textbf{Équilibrée.}
        \item \cle{Équation finale :} $\ce{2C + O2 -> 2CO}$
    \end{itemize}
    \item[d)] \textbf{Conclusion sur le type de combustion :} Les \textbf{deux} tests sont positifs (eau de chaux trouble $\rightarrow$ $\ce{CO2}$ présent ; test CO positif $\rightarrow$ $\ce{CO}$ présent). La combustion sous la cloche a été \textbf{mixte} : une partie du carbone a brûlé complètement, l'autre incomplètement. Cela s'explique par la consommation progressive de l'O₂ dans le volume fini de la cloche : au début, il y a assez d'O₂ pour faire du $\ce{CO2}$ ; à la fin, le manque d'O₂ force la réaction vers du $\ce{CO}$.
\end{enumerate}

\textbf{3. Analyse quantitative du drame.}
\begin{enumerate}
    \item[a)] \textbf{O₂ nécessaire pour combustion complète de 2,0 kg de C.}
    \begin{align*}
        n(\ce{C}) &= \frac{m(\ce{C})}{M(\ce{C})} = \frac{2000\ \text{g}}{12\ \text{g/mol}} \approx 166,7\ \text{mol} \\
        \text{Équation : } \ce{C + O2 -> CO2} &\Rightarrow \frac{n(\ce{O2})}{n(\ce{C})} = 1 \\
        n(\ce{O2})_{\text{nécessaire}} &= 166,7\ \text{mol} \\
        m(\ce{O2})_{\text{nécessaire}} &= n \times M(\ce{O2}) = 166,7 \times 32 \approx 5333\ \text{g} = \mathbf{5,33\ kg}
    \end{align*}
    \item[b)] \textbf{O₂ nécessaire pour combustion incomplète de 2,0 kg de C.}
    \begin{align*}
        \text{Équation : } \ce{2C + O2 -> 2CO} &\Rightarrow \frac{n(\ce{O2})}{n(\ce{C})} = \frac{1}{2} = 0,5 \\
        n(\ce{O2})_{\text{nécessaire}} &= 0,5 \times 166,7 \approx 83,3\ \text{mol} \\
        m(\ce{O2})_{\text{nécessaire}} &= 83,3 \times 32 \approx 2667\ \text{g} = \mathbf{2,67\ kg}
    \end{align*}
    \item[c)] \textbf{Comparaison et explication.}
    \begin{itemize}
        \item O₂ disponible dans la chambre : \textbf{\SI{7,6}{kg}}.
        \item O₂ pour combustion complète : \textbf{\SI{5,33}{kg}}.
        \item O₂ pour combustion incomplète : \textbf{\SI{2,67}{kg}}.
    \end{itemize}
    \emph{À première vue}, l'air de la chambre contient assez d'O₂ (\SI{7,6}{kg}) pour une combustion complète (\SI{5,33}{kg}). \textbf{MAIS} la chambre n'est pas un réacteur parfait : l'air n'est pas brassé, le brasero est au sol, l'O₂ se consomme localement plus vite qu'il ne diffuse. La concentration en O₂ chute rapidement autour du foyer. Dès qu'elle descend sous un seuil critique (environ \SI{14}{\percent}), la combustion bascule en mode incomplet. De plus, le CO formé consomme encore de l'O₂ s'il brûle ensuite ($\ce{2CO + O2 -> 2CO2}$), mais en milieu confiné, il s'accumule. \textbf{La combustion a été forcée à l'incomplète par la diffusion lente de l'O₂ vers le foyer.}
    \item[d)] \textbf{Vérification du volume de CO mesuré.}
    \begin{align*}
        V_{\text{chambre}} &= \SI{30}{\cubic\meter} = 30\,000\ \text{L} \\
        V(\ce{CO}) &= \num{500}~\text{ppm} \times 30\,000\ \text{L} = \frac{500}{1\,000\,000} \times 30\,000 = \mathbf{15\ L\ de\ CO} \\
        n(\ce{CO}) &= \frac{V(\ce{CO})}{V_m} = \frac{15\ \text{L}}{24\ \text{L/mol}} = 0,625\ \text{mol} \\
        m(\ce{CO}) &= n \times M(\ce{CO}) = 0,625 \times 28 = 17,5\ \text{g} \\
        \text{D'après } \ce{2C + O2 -> 2CO} &: n(\ce{C})_{\text{brûlé en CO}} = n(\ce{CO}) = 0,625\ \text{mol} \\
        m(\ce{C})_{\text{brûlé en CO}} &= 0,625 \times 12 = 7,5\ \text{g} \\
        \text{Pourcentage du charbon (2000 g) ayant donné du CO} &= \frac{7,5}{2000} \times 100 \approx \mathbf{0,38\ \%}
    \end{align*}
    \emph{Interprétation :} Seule une toute petite fraction du charbon (moins de 1 \%) a brûlé en CO pour produire \num{500}~ppm. Mais \num{500}~ppm est une concentration \textbf{très toxique} (seuil dangerosité immédiate \num{400}~ppm). Le reste du charbon a brûlé en $\ce{CO2}$ (consommant l'O₂) ou n'a pas brûlé. Cela montre l'extrême dangerosité : peu de CO suffit à tuer.
\end{enumerate}

\textbf{4. Mécanisme toxique et symptômes.}
\begin{enumerate}
    \item[a)] Le père a inhalé de l'air contenant \num{500}~ppm de CO. Même à cette « faible » concentration, l'affinité de l'hémoglobine pour le CO étant 250 fois plus forte que pour l'O₂, une fraction significative de son hémoglobine s'est transformée en carboxyhémoglobine (HbCO) en quelques minutes. Ses symptômes (maux de tête violents, nausées, confusion) correspondent à un taux d'HbCO estimé entre \SI{20}{\percent} et \SI{30}{\percent}. Le cerveau, gros consommateur d'O₂, est le premier organe touché par l'hypoxie.
    \item[b)] L'ouverture de la porte a provoqué un \textbf{renouvellement brutal de l'air} : entrée d'air frais riche en O₂ (\SI{21}{\percent}) et sortie du CO. Le père a cessé d'inhaler du CO et a ré-oxygéné son sang. Sa femme et son enfant, restés plus longtemps dans l'atmosphère viciée \textbf{avant} l'arrivée des secours, ont atteint un taux d'HbCO critique (> \SI{40-50}{\percent}), entraînant le coma. L'oxygénothérapie hyperbare (O₂ à haute pression) est le seul moyen de « chasser » rapidement le CO de l'hémoglobine pour sauver les vies en jeu.
\end{enumerate}

\begin{popfigure}
\begin{tikzpicture}[>=Stealth, node distance=1.2cm, font=\small]
\node[draw, circle, fill=popredL, minimum width=1.8cm] (Hb) {Hb};
\node[draw, circle, fill=popblueL, right=of Hb] (O2) {O₂};
\node[draw, circle, fill=poporangeL, below=of O2] (CO) {CO};
\draw[<->, thick, popblue] (Hb) -- node[midway, above, sloped] {Faible affinité} (O2);
\draw[<->, thick, poporange] (Hb) -- node[midway, left] {Affinité $\times 250$} (CO);
\node[draw, rounded corners, fill=popgreenL, right=of O2, align=center] (HbO2){HbO₂\\Transport O₂};
\node[draw, rounded corners, fill=popredL, right=of CO, align=center] (HbCO){HbCO\\Blocage O₂};
\draw[->, thick, popblue] (O2) -- (HbO2);
\draw[->, thick, poporange] (CO) -- (HbCO);
\end{tikzpicture}
\legende{Compétition entre O₂ et CO pour la fixation sur l'hémoglobine. Le CO « vole » la place de l'O₂.}
\end{popfigure}

\textbf{5. Affiche de prévention (Proposition de contenu pour l'enseignant).}
\begin{center}
\begin{tabularx}{\linewidth}{|X|X|}\hline
\rowcolor{popblueL}
\textbf{Élément de l'affiche} & \textbf{Contenu suggéré (texte + pictogramme)} \\ \hline
\textbf{Titre} & \textbf{⚠ BRASERO EN CHAMBRE FERMÉE = MORT ⚠} \\ \hline
\textbf{Message central} & \textbf{Le charbon + le manque d'air = CO (Gaz tueur, invisible, sans odeur)} \\ \hline
\textbf{Consigne 1} & \textbf{AÉREZ TOUJOURS !} \newline Ouvrez grande la porte et la fenêtre si vous utilisez du charbon/bois/gaz. \newline \textit{(Pictogramme : fenêtre ouverte + flèches d'air)} \\ \hline
\textbf{Consigne 2} & \textbf{INTERDIT : Brasero, réchaud à charbon, groupe électrogène DANS la chambre.} \newline Placez-les TOUJOURS à l'extérieur (cour, balcon aéré). \newline \textit{(Pictogramme : brasero barré en rouge dans une maison)} \\ \hline
\textbf{Consigne 3} & \textbf{SURVEILLEZ LA FLAMME.} \newline Flamme bleue = bonne combustion (gaz). Flamme jaune/orange/fumée = danger CO (charbon/bois/kérosène mal réglé). \newline \textit{(Pictogramme : flamme bleue ✓ vs flamme jaune/fumée ✗)} \\ \hline
\textbf{Consigne 4} & \textbf{INSTALLEZ UN DÉTECTEUR DE CO.} \newline Il sonne avant que vous ne sentiez le danger. Prix : \SIrange{5000}{10000}{FCFA} (marché Mokolo, supermarchés). \newline \textit{(Pictogramme : détecteur rond avec onde sonore)} \\ \hline
\textbf{Urgences} & \textbf{INTOXICATION ? (Mal de tête brutal, nausées, vertiges, confusion) :} \newline 1. Sortez TOUT LE MONDE immédiatement. \newline 2. Appelez \textbf{112} (Pompiers/SAMU) ou \textbf{119} (Police). \newline 3. N'y retournez pas avant l'avis des secours. \\ \hline
\end{tabularx}
\end{center}
\end{exoresolu}

\courssub{Bilan de l'activité}
Cette enquête t'a permis de mettre en œuvre, dans une situation réelle et dramatique, l'ensemble des savoir-faire de la leçon :
\begin{itemize}
    \item Tu as \textbf{identifié} une réaction chimique (combustion du carbone) à des indices macroscopiques (troublissement eau de chaux, test CO, masse constante).
    \item Tu as \textbf{appliqué la loi de Lavoisier} pour comprendre que les atomes de carbone et d'oxygène se réorganisent sans disparaître.
    \item Tu as \textbf{écrit et équilibré} les deux équations-bilan possibles ($\ce{C + O2 -> CO2}$ et $\ce{2C + O2 -> 2CO}$) et tu as compris que \textbf{la quantité de dioxygène disponible impose le produit formé}.
    \item Tu as \textbf{calculé} les masses d'O₂ nécessaires et comparé à la réalité d'une chambre close, prouvant que le confinement crée inévitablement un manque d'O₂ local, forçant la combustion incomplète.
    \item Tu as \textbf{lié la chimie à la biologie} : la formation de quelques grammes de CO (0,38 \% du charbon) suffit à saturer l'hémoglobine (affinité 250x) et à provoquer le coma ou la mort.
    \item Tu as \textbf{produit un message citoyen} (l'affiche) transformant la connaissance scientifique en action de prévention concrète pour ton quartier.
\end{itemize}

\begin{attention}[Pièges à éviter pour le BEPC]
\begin{itemize}
    \item \textbf{Ne pas confondre « équation » et « équation-bilan ».} L'équation-bilan utilise des formules chimiques et des coefficients stœchiométriques. Elle ne comporte pas de signe $=$ mais une flèche $\rightarrow$.
    \item \textbf{Ne jamais modifier les formules chimiques pour équilibrer.} Exemple interdit : $\ce{C + O2 -> CO}$ (pas équilibré) $\rightarrow$ on n'écrit pas $\ce{C + O -> CO}$. On met des \textbf{coefficients} : $\ce{2C + O2 -> 2CO}$.
    \item \textbf{Combustion incomplète $\neq$ « pas de réaction ».} C'est une réaction chimique complète qui forme un produit différent (CO au lieu de $\ce{CO2}$) par manque de réactif (O₂).
    \item \textbf{Unités :} Masse en kg ou g (cohérence), quantité de matière en mol, volume en L ou m³. Attention aux conversions (\SI{1}{\cubic\meter} = \SI{1000}{L}).
    \item \textbf{Contexte :} Au Cameroun, le brasero (charbon de bois) et le réchaud à pétrole (kérosène) sont les causes principales d'intoxication au CO en saison sèche. Le groupe électrogène dans le couloir ou la véranda fermée est aussi un tueur silencieux.
\end{itemize}
\end{attention}

\begin{savaistu}[Le saviez-vous ?]
\begin{itemize}
    \item Le monoxyde de carbone (CO) est aussi un \textbf{polluant atmosphérique} majeur émis par les moteurs thermiques (motos, taxis, voitures) dans les embouteillages de Yaoundé ou Douala. C'est pour cela que les normes d'émission (contrôle technique) limitent le CO à l'échappement.
    \item Les \textbf{détecteurs de CO} domestiques utilisent souvent une pastille de gel de silice imprégnée de sels de palladium ou de molybdène qui change de couleur (jaune $\rightarrow$ brun/noir) au contact du CO. Ils ne détectent \textbf{pas} la fumée (ce sont les détecteurs de fumée/incendie) ni le gaz butane/propane (détecteurs de gaz explosifs). Il faut le bon détecteur pour le bon danger !
    \item L'\textbf{oxygénothérapie hyperbare} (caisson pressurisé à 2-3 bars d'O₂ pur) réduit la demi-vie de la carboxyhémoglobine de \SI{320}{min} (air ambiant) à \SI{23}{min}. C'est une urgence vitale disponible à l'Hôpital Central et à l'Hôpital Général de Yaoundé.
\end{itemize}
\end{savaistu}

\newpage

\coursec{Méthodes et exercices résolus}

\coursec{La réaction chimique}

\courssub{Transformation chimique : mise en évidence expérimentale}

\begin{definition}[Transformation chimique]
Une \textbf{transformation chimique} est un processus au cours duquel une ou plusieurs espèces chimiques (les \textbf{réactifs}) disparaissent pour former une ou plusieurs nouvelles espèces chimiques (les \textbf{produits}). Les atomes se réarrangent : les liaisons chimiques se rompent et se reforment différemment.
\end{definition}

\begin{experience}[Comment reconnaître une transformation chimique ?]
\textbf{Matériel :} Tubes à essai, bec Bunsen, balance, acide chlorhydrique dilué ($\ce{HCl}$), zinc métallique ($\ce{Zn}$), sulfate de cuivre ($\ce{CuSO4}$), soude ($\ce{NaOH}$), chlorure de sodium ($\ce{NaCl}$), bâtonnet incandescent, papier pH.
\textbf{Protocole et observations :}
\begin{enumerate}
    \item \textbf{Expérience A (Zn + HCl) :} Verser $\ce{HCl}$ sur des granules de $\ce{Zn}$. \textit{Observation :} Bulles abondantes, tube chauffe. Approcher une flamme : \textbf{« pop » sec} (détonation).
    \item \textbf{Expérience B ($\ce{CuSO4}$ + $\ce{NaOH}$) :} Verser $\ce{NaOH}$ dans la solution bleue de $\ce{CuSO4}$. \textit{Observation :} Formation immédiate d'un \textbf{solide bleu clair} (précipité), solution surnageante incolore.
    \item \textbf{Expérience C (Chauffage $\ce{NaCl}$) :} Chauffer fortement du sel de cuisine. \textit{Observation :} Fusion (liquide incolore), puis solidification à l'identique en refroidissant. Pas de gaz, pas de couleur nouvelle.
\end{enumerate}
\textbf{Interprétation :} Expériences A et B : nouvelles espèces formées ($\ce{H2}$, $\ce{Cu(OH)2}$) $\rightarrow$ \textbf{transformation chimique}. Expérience C : seul l'état change ($\ce{NaCl(s)} \rightleftharpoons \ce{NaCl(l)}$) $\rightarrow$ \textbf{transformation physique}.
\legende{Trois tests classiques : dégagement de gaz inflammable ($\ce{H2}$), formation d'un précipité coloré ($\ce{Cu(OH)2}$), changement d'état pur ($\ce{NaCl}$).}
\end{experience}

\begin{aretenir}[Signes macroscopiques d'une transformation chimique]
Au laboratoire, on conclut à une transformation chimique si l'on observe \textbf{au moins un} de ces signes visibles :
\begin{itemize}
    \item \textbf{Dégagement de gaz} (bulles, mousse, bruit) : ex. $\ce{H2}$, $\ce{CO2}$, $\ce{O2}$.
    \item \textbf{Apparition/disparition d'un solide} (précipité, dépôt, dissolution d'un métal).
    \item \textbf{Changement de couleur} franc de la solution ou du solide (pas simple dilution).
    \item \textbf{Variation de température} spontanée (exothermique : chauffe / endothermique : refroidit).
    \item \textbf{Modification de propriétés} : inflammabilité (test flamme), pH, conductivité.
\end{itemize}
\textbf{Attention :} Un mélange (ex. eau + sirop), une dissolution sans réaction (sucre, sel), un changement d'état (glace $\to$ eau) ne sont \textbf{PAS} des transformations chimiques.
\end{aretenir}

\begin{methode}[Identifier une transformation chimique au laboratoire]
\begin{enumerate}
    \item \textbf{Observer} soigneusement avant, pendant, après (couleur, état, gaz, température).
    \item \textbf{Noter} le ou les signes observés (vocabulaire précis : « bulles », « précipité bleu », « flamme verte », etc.).
    \item \textbf{Conclure} : « J'observe \emph{[signe]} ce qui prouve la formation de \emph{[nouvelle espèce]} $\rightarrow$ c'est une transformation chimique. »
    \item \textbf{Écrire} si possible l'équation-bilan (voir §3).
\end{enumerate}
\end{methode}

\begin{exemplebox}[Exemples courants au Cameroun]
\begin{itemize}
    \item \textbf{Cuisson (haricots, viande, beignets) :} Brunissement (réaction de Maillard), dégagement d'odeurs, changement texture $\rightarrow$ transformations chimiques complexes.
    \item \textbf{Rouille sur une moto / grille de fenêtre :} $\ce{Fe} + \ce{O2} + \ce{H2O} \to \ce{Fe2O3.nH2O}$ (solide orangé friable) $\rightarrow$ transformation chimique lente.
    \item \textbf{Pile usagée (lampe torche) :} Réaction redox interne $\rightarrow$ transformation chimique produisant du courant.
    \item \textbf{Fabrication du savon (saponification) :} Huile (palme, coton) + Soude $\to$ Savon + Glycérine.
\end{itemize}
\end{exemplebox}

\courssub{Conservation de la masse et loi de Lavoisier}

\begin{experience}[Loi de Lavoisier : système fermé vs ouvert]
\textbf{Matériel :} Balance électronique (précision 0,1 g), bécher 100 mL, erlenmeyer avec bouchon, acide chlorhydrique ($\ce{HCl}$ 1 mol/L), carbonate de calcium ($\ce{CaCO3}$), acide acétique ($\ce{CH3COOH}$), bicarbonate de sodium ($\ce{NaHCO3}$).
\textbf{Protocole :}
\begin{enumerate}
    \item \textbf{Montage 1 (Système OUVERT) :} Bécher sur balance $\to$ Tare. Ajouter $\SI{50}{mL}$ $\ce{HCl}$ + $\SI{5}{g}$ $\ce{CaCO3}$. Noter masse initiale. Laisser réagir (bulles $\ce{CO2}$). Noter masse finale.
    \item \textbf{Montage 2 (Système FERMÉ) :} Erlenmeyer + bouchon (percé pour tuyau vers tube à essai d'eau de chaux) sur balance $\to$ Tare. Même réaction. Noter masses.
\end{enumerate}
\textbf{Résultats types :}
\begin{center}
\begin{tabularx}{\linewidth}{|l|X|X|}\hline
\textbf{Montage} & \textbf{Masse initiale (g)} & \textbf{Masse finale (g)} \\ \hline
Ouvert (bécher) & 61,0 & 57,4 (\textbf{perte 3,6 g}) \\ \hline
Fermé (erlenmeyer bouché) & 61,0 & 61,0 (\textbf{inchangée}) \\ \hline
\end{tabularx}
\end{center}
\textbf{Interprétation :} Dans le système \textbf{fermé}, rien ne sort, la masse totale ne change pas. Dans le système \textbf{ouvert}, le gaz $\ce{CO2}$ s'échappe $\rightarrow$ perte de masse égale à la masse de gaz perdu.
\legende{Vérification de la loi de Lavoisier : la masse ne se perd pas, elle se redistribue.}
\end{experience}

\begin{propriete}[Loi de Lavoisier (Conservation de la masse)]
Lors d'une transformation chimique dans un \textbf{système fermé} (aucun échange de matière avec l'extérieur), la \textbf{masse totale du système ne change pas} :
\[ m_{\text{total initial (réactifs)}} = m_{\text{total final (produits)}} \]
\textbf{Conséquence :} Les atomes ne sont ni créés, ni détruits, ils se réarrangent. Le nombre d'atomes de chaque élément est conservé.
\end{propriete}

\begin{methode}[Appliquer la loi de Lavoisier (calcul de masse manquante)]
\begin{enumerate}
    \item \textbf{Définir le système} : Fermé (masse constante) ou Ouvert (gaz s'échappe $\rightarrow$ perte de masse).
    \item \textbf{Écrire l'équation-bilan équilibrée} (voir §3).
    \item \textbf{Appliquer la conservation} : $\sum m_{\text{réactifs initiaux}} = \sum m_{\text{produits finaux}}$ (système fermé).
    \item \textbf{Calculer l'inconnue} : Soustraction directe (masse manquante = $m_{\text{initial}} - m_{\text{final mesuré}}$) \textbf{OU} calcul stœchiométrique ($n = m/M$, tableau d'avancement).
\end{enumerate}
\end{methode}

\begin{exoresolu}[Réaction acide-carbonate : perte de masse (Système ouvert)]
\textbf{Énoncé :} On verse \SI{50,0}{mL} d'une solution d'acide chlorhydrique ($\ce{HCl}$) de masse \SI{51,0}{g} sur \SI{10,0}{g} de carbonate de calcium solide ($\ce{CaCO3}$) dans un bécher posé sur une balance. La réaction produit du chlorure de calcium ($\ce{CaCl2}$), de l'eau et du dioxyde de carbone gazeux ($\ce{CO2}$) qui s'échappe. La masse finale indiquée est \SI{57,4}{g}.
\begin{enumerate}
    \item Écrire l'équation-bilan.
    \item La balance mesure-t-elle un système fermé ou ouvert ? Justifier.
    \item Calculer la masse de $\ce{CO2}$ échappée.
    \item En déduire la quantité de matière de $\ce{CO2}$ formée ($M(\ce{CO2})=\SI{44}{g/mol}$).
\end{enumerate}
\tcblower
\textbf{Solution détaillée :}
\begin{enumerate}
    \item \textbf{Équation-bilan :} $\ce{CaCO3(s) + 2 HCl(aq) -> CaCl2(aq) + H2O(l) + CO2(g)}$
    (Vérif : Ca:1, C:1, O:5, H:2, Cl:2 des deux côtés).
    \item \textbf{Système OUVERT.} Le bécher n'est pas bouché. Le gaz $\ce{CO2}$ formé s'échappe dans l'air. La balance ne pèse que le contenu liquide/solide restant.
    \item \textbf{Masse $\ce{CO2}$ :} $m_{\text{initial}} = 51,0 + 10,0 = \SI{61,0}{g}$. $m_{\text{final}} = \SI{57,4}{g}$.
    $m(\ce{CO2}) = 61,0 - 57,4 = \SI{3,6}{g}$.
    \item \textbf{Quantité de matière :} $n(\ce{CO2}) = \frac{m}{M} = \frac{3,6}{44} \approx \SI{0,0818}{mol}$.
\end{enumerate}
\end{exoresolu}

\courssub{Équation-bilan et équilibrage}

\begin{definition}[Équation-bilan d'une transformation chimique]
L'\textbf{équation-bilan} (ou équation chimique) représente symboliquement la transformation :
\[ \text{Réactifs} \xrightarrow{\text{conditions}} \text{Produits} \]
Elle respecte deux règles fondamentales :
\begin{enumerate}
    \item \textbf{Conservation des éléments :} Même nombre d'atomes de chaque élément de part et d'autre de la flèche (loi de Lavoisier).
    \item \textbf{Conservation de la charge :} Somme des charges égales des deux côtés (pour les équations ioniques).
\end{enumerate}
Les coefficients stœchiométriques (nombres entiers devant les formules) indiquent les proportions molaires. On \textbf{ne modifie jamais les indices} dans les formules chimiques ($\ce{H2O}$ pas $\ce{H2O2}$ pour équilibrer l'O).
\end{definition}

\begin{methode}[Équilibrer une équation-bilan (Méthode par inspection)]
\begin{enumerate}
    \item \textbf{Identifier réactifs et produits} (noms $\to$ formules brutes $\ce{}$). Écrire le squelette (flèche, pas de coeff).
    \item \textbf{Équilibrer les éléments « rares »} (métaux, C, S, P, etc.) un par un.
    \item \textbf{Équilibrer l'hydrogène (H)} puis \textbf{l'oxygène (O)} en dernier (souvent présents dans plusieurs espèces).
    \item \textbf{Vérifier} chaque élément. Si coefficients fractionnaires $\to$ multiplier toute l'équation par le dénominateur commun pour obtenir des entiers.
    \item \textbf{Préciser les états physiques} si demandés : $(s)$ solide, $(l)$ liquide, $(g)$ gaz, $(aq)$ aqueux (dissous dans l'eau).
\end{enumerate}
\textbf{Astuce :} Pour hydrocarbure $\ce{C_xH_y}$ combustion complète : $\ce{C_xH_y + (x + y/4) O2 -> x CO2 + (y/2) H2O}$.
\end{methode}

\begin{exoresolu}[Équilibrage : combustion du butane (bonbonne gaz)]
\textbf{Énoncé :} Le butane $\ce{C4H10}$ brûle complètement dans l'air ($\ce{O2}$) pour donner $\ce{CO2}$ et $\ce{H2O}$.
\begin{enumerate}
    \item Écrire l'équation-bilan équilibrée (coeff entiers).
    \item Calculer la masse de $\ce{O2}$ nécessaire pour \SI{58}{g} de butane ($M_{\ce{C}}=12, M_{\ce{H}}=1, M_{\ce{O}}=16$).
\end{enumerate}
\tcblower
\textbf{Solution :}
\begin{enumerate}
    \item Squelette : $\ce{C4H10 + O2 -> CO2 + H2O}$.
    C : 4 $\to$ $\ce{4 CO2}$. H : 10 $\to$ $\ce{5 H2O}$.
    O à droite : $4\times2 + 5 = 13$ atomes $\to$ $6,5~\ce{O2}$.
    $\times 2$ pour entiers : $\boxed{\ce{2 C4H10(g) + 13 O2(g) -> 8 CO2(g) + 10 H2O(g)}}$.
    \item $M(\ce{C4H10}) = 4\times12 + 10\times1 = \SI{58}{g/mol}$. $n(\ce{C4H10}) = \frac{58}{58} = \SI{1,0}{mol}$.
    Eq. coeff 1 : $\ce{C4H10 + 6,5 O2 -> 4 CO2 + 5 H2O}$.
    $n(\ce{O2}) = 6,5 \times 1,0 = \SI{6,5}{mol}$. $M(\ce{O2}) = \SI{32}{g/mol}$.
    $m(\ce{O2}) = 6,5 \times 32 = \boxed{\SI{208}{g}}$.
\end{enumerate}
\end{exoresolu}

\courssub{Combustions complète et incomplète}

\begin{propriete}[Combustion complète vs incomplète]
La combustion est la réaction rapide d'un \textbf{combustible} (souvent hydrocarbure $\ce{C_xH_y}$) avec le \textbf{comburant} (dioxygène $\ce{O2}$ de l'air), dégageant chaleur et lumière.
\begin{itemize}
    \item \textbf{Complète} : $\ce{O2}$ en \textbf{excès} (ou quantité stœchiométrique). Produits stables : \textbf{Dioxyde de carbone $\ce{CO2}$} + \textbf{Eau $\ce{H2O}$}. Flamme \textbf{bleue}, non fumante. Rendement énergétique maximal.
    \item \textbf{Incomplète} : $\ce{O2}$ \textbf{limitant} (manque d'air, brûleur encrassé, moteur au ralenti). Produits toxiques/polluants :
    \begin{itemize}
        \item \textbf{Monoxyde de carbone $\ce{CO}$} (gaz incolore, inodore, \textbf{mortel}).
        \item \textbf{Carbone $\ce{C}$ (suie, noir de carbone)} : particules fines, fumée \textbf{noire}, flamme \textbf{jaune/orange}.
    \end{itemize}
\end{itemize}
\end{propriete}

\begin{methode}[Écrire l'équation-bilan d'une combustion incomplète]
\begin{enumerate}
    \item \textbf{Identifier le produit carboné} selon l'énoncé : $\ce{CO}$ (manque modéré $\ce{O2}$) ou $\ce{C}$ (manque sévère $\ce{O2}$). L'eau $\ce{H2O}$ est toujours formée.
    \item \textbf{Écrire le squelette} : $\ce{C_xH_y + O2 -> CO (ou C) + H2O}$.
    \item \textbf{Équilibrer} : C $\to$ H $\to$ O (attention : $\ce{CO}$ a 1 O, $\ce{C}$ a 0 O).
    \item \textbf{Mémoriser pour le méthane $\ce{CH4}$ (gaz naturel, biogaz) :}
    \begin{align*}
        \text{Complète} &: \ce{CH4 + 2 O2 -> CO2 + 2 H2O} \\
        \text{Incomplète (CO)} &: \ce{2 CH4 + 3 O2 -> 2 CO + 4 H2O} \\
        \text{Incomplète (C/soot)} &: \ce{CH4 + O2 -> C + 2 H2O}
    \end{align*}
\end{enumerate}
\end{methode}

\begin{exoresolu}[Groupe électrogène : combustion incomplète (Octane $\ce{C8H18}$)]
\textbf{Énoncé :} Un groupe électrogène à essence (octane $\ce{C8H18}$) tourne au ralenti, admission d'air bouchée. Fumée noire à l'échappement.
\begin{enumerate}
    \item Type de combustion ? Justifier (2 observations).
    \item Équation-bilan produisant du $\ce{CO}$.
    \item Équation-bilan produisant de la suie $\ce{C}$.
    \item Danger mortel en chambre fermée ?
\end{enumerate}
\tcblower
\textbf{Solution :}
\begin{enumerate}
    \item \textbf{Combustion incomplète.} 1) « Admission d'air bouchée » $\to$ $\ce{O2}$ limitant. 2) « Fumée noire » $\to$ particules de carbone (suie) non brûlées.
    \item \textbf{Produit $\ce{CO}$ :} $\ce{C8H18 + O2 -> CO + H2O}$.
    C:8 $\to$ $\ce{8 CO}$. H:18 $\to$ $\ce{9 H2O}$.
    O à droite : $8\times1 + 9 = 17$ $\to$ $8,5~\ce{O2}$. $\times 2$ :
    $\boxed{\ce{2 C8H18 + 17 O2 -> 16 CO + 18 H2O}}$.
    \item \textbf{Produit $\ce{C}$ (suie) :} $\ce{C8H18 + O2 -> C + H2O}$.
    C:8 $\to$ $\ce{8 C}$. H:18 $\to$ $\ce{9 H2O}$.
    O à droite : $9$ $\to$ $4,5~\ce{O2}$. $\times 2$ :
    $\boxed{\ce{2 C8H18 + 9 O2 -> 16 C + 18 H2O}}$.
    \item \textbf{Mortel :} $\ce{CO}$ incolore/inodore s'accumule. Se fixe sur l'hémoglobine (Hb) du sang (affinité $\times 200-250$ vs $\ce{O2}$) $\to$ \textbf{carboxyhémoglobine (HbCO)}. Le sang ne transporte plus l'$\ce{O2}$ $\to$ \textbf{hypoxie cellulaire}, coma, mort en minutes. Suie = pollution particules fines (PM2,5).
\end{enumerate}
\end{exoresolu}

\courssub{Le monoxyde de carbone : un danger invisible}

\begin{savaistu}[Le CO au Cameroun : contexte réel]
\begin{itemize}
    \item \textbf{Sources principales :} Groupes électrogènes (coupures ENEO) dans cours/garages fermés ; réchauds à pétrole/kérosène (saison sèche) ; braseros (charbon de bois) dans chambres ; moteurs moto/voiture au ralenti ; brûleurs gaz mal réglés (flamme jaune).
    \item \textbf{Statistiques :} Chaque année, des familles entières décèdent (Yaoundé, Douala, Bafoussam...) par asphyxie au CO. Le CO ne réveille pas (pas d'irritation), il endort définitivement.
    \item \textbf{Seuil de danger :} \SI{0,1}{\%} CO dans l'air (\SI{1000}{ppm}) $\to$ mortel en 1 h. \SI{0,3}{\%} $\to$ mortel en quelques minutes.
\end{itemize}
\end{savaistu}

\begin{definition}[Mécanisme toxique du monoxyde de carbone]
Le $\ce{CO}$ inhalé traverse la barrière alvéolo-capillaire et se lie au \textbf{fer (Fe$^{2+}$) de l'hémoglobine (Hb)} des globules rouges :
\[ \ce{Hb + CO <=> HbCO} \quad \text{(Carboxyhémoglobine)} \]
L'affinité $\ce{Hb-CO}$ est \textbf{200 à 250 fois plus forte} que l'affinité $\ce{Hb-O2}$. L'hémoglobine est « bloquée » par le $\ce{CO}$ et ne peut plus charger l'$\ce{O2}$ au niveau des poumons $\to$ \textbf{asphyxie tissulaire} (cerveau, cœur privés d'$\ce{O2}$).
\end{definition}

\begin{exemplebox}[Symptômes d'intoxication au CO (ordre d'apparition)]
\begin{enumerate}
    \item Maux de tête (céphalées), vertiges, nausées, vomissements.
    \item Fatigue inhabituelle, confusion, troubles de la vision, bourdonnements d'oreilles.
    \item Perte de connaissance (syncope), convulsions, arrêt cardio-respiratoire.
\end{enumerate}
\textbf{Particularité :} La peau peut devenir \textbf{rose cerise} (couleur de l'HbCO), mais ce signe est tardif et peu fiable sur peau foncée. \textbf{Ne jamais attendre les symptômes pour aérer !}
\end{exemplebox}

\courssub{Sécurité et gestes responsables autour des combustions}

\begin{methode}[Prévention du risque CO : La règle A.E.E.E.]
\begin{enumerate}
    \item \textbf{A} \textbf{ÉRER} \textbf{systématiquement} : Ouvrir fenêtres/portes \textbf{10-15 min chaque jour} (même saison sèche/froid) et \textbf{pendant} toute utilisation d'appareil à combustion (gazinière, chauffage pétrole, groupe). \textbf{Ne jamais boucher} les grilles de ventilation (haut et bas) obligatoires dans les cuisines/salles de bain.
    \item \textbf{E} \textbf{XTÉRIEUR} pour les groupes : Groupe électrogène, moto, voiture $\to$ \textbf{JAMAIS} dans pièce close, garage, sous-sol, véranda, même porte ouverte. Éloigner l'échappement (\textgreater 2 m) des portes/fenêtres/ventilation.
    \item \textbf{E} \textbf{NTRETENIR} les appareils : Vérifier bonbonne (date), détendeur, tuyau (norme NF, date péremption), brûleurs. \textbf{Flamme bleue = bonne combustion.} Flamme jaune/orange = danger $\to$ nettoyer injecteurs, faire régler par pro.
    \item \textbf{E} \textbf{QUIPER} en détecteur : Installer détecteur $\ce{CO}$ normé \textbf{EN 50291} dans \textbf{chaque chambre} et couloir (pas cuisine : vapeur cuisson = fausses alertes). Tester bouton \textbf{Test} chaque mois. Changer piles / appareil selon notice (généralement 5-7 ans).
\end{enumerate}
\textbf{Mnémotechnique :} \textbf{A.E.E.E.} = \textbf{A}érer, \textbf{E}xtérieur, \textbf{E}ntretenir, \textbf{E}quiper.
\end{methode}

\begin{exoresolu}[Analyse situation à risque : Réchaud pétrole au salon]
\textbf{Énoncé :} Saison sèche, matin frais. Famille utilise réchaud pétrole (kérosène) dans salon fenêtres fermées « pour garder la chaleur ». Pas de détecteur $\ce{CO}$. Père : mal de tête. Fille : nausées.
\begin{enumerate}
    \item Identifier 3 erreurs majeures de sécurité.
    \item Expliquer le mécanisme physiologique (lien hémoglobine).
    \item Lister les 4 actions immédiates \textbf{dans l'ordre vital}.
\end{enumerate}
\tcblower
\textbf{Solution :}
\begin{enumerate}
    \item \textbf{Erreurs :} 1) \textbf{Zéro aération} (fenêtres fermées) $\to$ accumulation $\ce{CO}$ inévitable (combustion incomplète par manque $\ce{O2}$). 2) \textbf{Appareil non évacué} (pas de conduit fumée) utilisé en chauffage continu volume clos. 3) \textbf{Absence détecteur $\ce{CO}$} $\to$ pas d'alerte précoce avant symptômes.
    \item \textbf{Mécanisme :} $\ce{CO}$ inhalé $\to$ alvéoles $\to$ sang $\to$ fixation sur \textbf{Fe$^{2+}$ de l'hémoglobine (Hb)} $\to$ \textbf{Carboxyhémoglobine (HbCO)}. Affinité $\ce{Hb-CO} \approx \mathbf{200-250\times}$ affinité $\ce{Hb-O2}$. Hb saturée en $\ce{CO}$ $\to$ \textbf{incapacité à transporter $\ce{O2}$} $\to$ \textbf{hypoxie tissulaire sévère} (cerveau, cœur).
    \item \textbf{Actions immédiates (Ordre vital) :}
    1. \textbf{ÉVACUER} \textbf{immédiatement} toutes les personnes à l'air libre (sortir de la maison).
    2. \textbf{APPELER} les secours \textbf{112} ou \textbf{119} (Cameroun) : dire « suspicion intoxication monoxyde de carbone ».
    3. \textbf{AÉRER} grand ouvert portes/fenêtres \textbf{APRÈS} avoir sorti les victimes (ne pas risquer sa vie pour ouvrir).
    4. \textbf{NE PAS RALLUMER} l'appareil. Faire contrôler par technicien agréé avant réutilisation.
\end{enumerate}
\end{exoresolu}

\courssub{Rendement d'une transformation chimique (Approfondissement)}

\begin{definition}[Rendement d'une transformation]
Le \textbf{rendement} $\eta$ (lettre grecqueêta) compare la quantité de matière \textbf{réellement obtenue} d'un produit à la quantité \textbf{théorique maximale} calculée par stœchiométrie (réactif limitatif consommé totalement).
\[ \eta = \frac{n_{\text{réel (produit)}}}{n_{\text{théo (produit)}}} \times 100 \quad (\text{en \%}) \qquad \text{avec } 0 \le \eta \le 100\% \]
$\eta < 100\%$ car : réactions parasites (combustion incomplète), pertes mécaniques, équilibre chimique non total, impuretés réactifs.
\end{definition}

\begin{methode}[Calculer le rendement (Étapes rigoureuses)]
\begin{enumerate}
    \item \textbf{Équation-bilan équilibrée.}
    \item \textbf{Quantités de matière initiales} des réactifs : $n = m/M$ (solide/gaz pur) ou $n = C \times V$ (solution).
    \item \textbf{Réactif limitatif :} Calculer avances maximales $x_{\max,i} = \frac{n_{\text{init},i}}{|\nu_i|}$. Le plus petit $x_{\max}$ correspond au réactif limitatif.
    \item \textbf{Quantité théorique du produit visé :} $n_{\text{théo}} = |\nu_{\text{produit}}| \times x_{\max}$.
    \item \textbf{Quantité réelle obtenue :} Donnée par l'énoncé (masse pesée, volume gaz mesuré, titration...). Convertir en mol ($n_{\text{réel}}$).
    \item \textbf{Rendement :} $\eta = \frac{n_{\text{réel}}}{n_{\text{théo}}} \times 100$. Vérifier $\eta \le 100\%$.
\end{enumerate}
\end{methode}

\begin{exoresolu}[Rendement combustion essence moto-taxi (Octane $\ce{C8H18}$)]
\textbf{Énoncé :} Un moto-taxi consomme \SI{1,0}{L} d'essence (octane $\ce{C8H18}$, $\rho = \SI{0,74}{g/mL}$). Le moteur rejette \SI{1,8}{kg} de $\ce{CO2}$. Combustion complète admise. ($M(\ce{C8H18})=\SI{114}{g/mol}$, $M(\ce{CO2})=\SI{44}{g/mol}$).
\begin{enumerate}
    \item Masse d'octane brûlée.
    \item Quantité théorique de $\ce{CO2}$ producible.
    \item Quantité réelle de $\ce{CO2}$ rejetée.
    \item Rendement de la transformation en $\ce{CO2}$.
\end{enumerate}
\tcblower
\textbf{Solution :}
\begin{enumerate}
    \item $V = \SI{1,0}{L} = \SI{1000}{mL}$. $m(\ce{C8H18}) = \rho V = 0,74 \times 1000 = \boxed{\SI{740}{g}}$.
    \item $n(\ce{C8H18}) = \frac{740}{114} \approx \SI{6,491}{mol}$.
    Eq. complète : $\ce{C8H18 + 12,5 O2 -> 8 CO2 + 9 H2O}$ (coeff 1 pour octane).
    Rapport stœchiométrique : $1~\ce{C8H18} \to 8~\ce{CO2}$.
    $n_{\text{théo}}(\ce{CO2}) = 8 \times 6,491 \approx \boxed{\SI{51,93}{mol}}$.
    \item $m_{\text{réel}}(\ce{CO2}) = \SI{1,8}{kg} = \SI{1800}{g}$.
    $n_{\text{réel}}(\ce{CO2}) = \frac{1800}{44} \approx \boxed{\SI{40,91}{mol}}$.
    \item $\eta = \frac{40,91}{51,93} \times 100 \approx \boxed{78,8\%}$.
    \textbf{Interprétation :} Environ 21\% du carbone de l'essence ne finit pas en $\ce{CO2}$ : combustion incomplète ($\ce{CO}$, suie $\ce{C}$), hydrocarbures imbrûlés (HC), pertes thermiques. C'est typique d'un moteur 2 temps ou 4 temps mal entretenu.
\end{enumerate}
\end{exoresolu}

\begin{attention}[Les 5 erreurs qui coûtent des points au BEPC]
\begin{enumerate}
    \item \textbf{Équation mal équilibrée :} Oublier de vérifier l'O en dernier, laisser coeff fractionnaires ($6,5~\ce{O2}$) sans multiplier par 2, \textbf{modifier les indices des formules} (ex: écrire $\ce{CO}$ au lieu de $\ce{CO2}$ pour « équilibrer »). \textit{Règle d'or : On ne touche JAMAIS aux indices (formules), seulement aux coefficients stœchiométriques.}
    \item \textbf{Confondre système fermé / ouvert (Lavoisier) :} Appliquer $m_i = m_f$ alors qu'un gaz s'échappe (bécher ouvert). \textit{Réflexe : Dessiner le système, flèches entrant/sortant. Si gaz sort $\to$ système ouvert $\to$ $m_{\text{final}} < m_{\text{initial}}$. Perte = masse gaz échappé.}
    \item \textbf{Mélanger combustion complète / incomplète :} Écrire $\ce{CO2}$ alors que l'énoncé dit « manque d'air », « flamme jaune », « fumée noire », « ralenti ». \textit{Mots-clés : Excès $\ce{O2}$ / air $\to$ $\ce{CO2}$ ; Manque $\ce{O2}$ / air $\to$ $\ce{CO}$ et/ou $\ce{C}$ (suie).}
    \item \textbf{Erreur rendement :} Calculer $n_{\text{théo}}$ sans identifier le \textbf{réactif limitatif} (utiliser celui en excès). Ou inverser la fraction $\frac{n_{\text{théo}}}{n_{\text{réel}}}$ (donne $>100\%$). \textit{Vérif systématique : Rendement $\le 100\%$.}
    \item \textbf{Sécurité CO : rédaction bâclée.} Dire « ouvrir la fenêtre » sans sortir les victimes \textbf{en premier}. Oublier numéros urgence (112/119 Cameroun). Confondre $\ce{CO}$ (monoxyde, toxique fixateur Hb) et $\ce{CO2}$ (dioxyde, asphyxiant par déplacement $\ce{O2}$ mais non fixateur). \textit{Ordre vital : 1. Sortir (Évacuer), 2. Appeler (112/119), 3. Aérer.}
\end{enumerate}
\end{attention}

\begin{exercice}[Entraînement BEPC : Réaction du zinc et acide sulfurique]
On verse \SI{50,0}{mL} d'une solution d'acide sulfurique $\ce{H2SO4}$ (\SI{1,0}{mol/L}) sur \SI{6,5}{g} de zinc métallique $\ce{Zn}$ dans un erlenmeyer bouché (système fermé) posé sur une balance. La réaction produit du sulfate de zinc $\ce{ZnSO4}$ et du dihydrogène $\ce{H2}$.
\begin{enumerate}
    \item Écrire l'équation-bilan de la réaction.
    \item Calculer les quantités de matière initiales de $\ce{Zn}$ et $\ce{H2SO4}$.
    \item Identifier le réactif limitatif et calculer l'avancement maximal $x_{\max}$.
    \item Calculer la masse théorique de $\ce{H2}$ formée.
    \item La balance indique-t-elle une variation de masse ? Justifier.
\end{enumerate}
\end{exercice}

\begin{corrige}[Exercice : Zinc + Acide sulfurique]
\begin{enumerate}
    \item $\ce{Zn(s) + H2SO4(aq) -> ZnSO4(aq) + H2(g)}$ (Équilibrée : Zn:1, H:2, S:1, O:4).
    \item $n(\ce{Zn}) = \frac{6,5}{65} = \SI{0,10}{mol}$ ($M_{\ce{Zn}}=65$). $n(\ce{H2SO4}) = C \times V = 1,0 \times 0,050 = \SI{0,050}{mol}$.
    \item Avances max : $x_{\max}(\ce{Zn}) = 0,10/1 = 0,10$ mol. $x_{\max}(\ce{H2SO4}) = 0,050/1 = \SI{0,050}{mol}$.
    $\ce{H2SO4}$ est \textbf{limitatif}. $x_{\max} = \SI{0,050}{mol}$.
    \item $n_{\text{théo}}(\ce{H2}) = 1 \times x_{\max} = \SI{0,050}{mol}$. $M(\ce{H2}) = \SI{2}{g/mol}$.
    $m_{\text{théo}}(\ce{H2}) = 0,050 \times 2 = \boxed{\SI{0,10}{g}}$.
    \item \textbf{Non, la masse ne varie pas.} Système \textbf{fermé} (erlenmeyer bouché). Le gaz $\ce{H2}$ reste piégé dans l'erlenmeyer. Loi de Lavoisier : $m_{\text{initial}} = m_{\text{final}}$.
\end{enumerate}
\end{corrige}

\begin{exercice}[Entraînement BEPC : Détecteur de CO et combustion]
\begin{enumerate}
    \item Un chauffage d'appoint au kérosène fonctionne dans une chambre de \SI{20}{m^3} (hauteur 2,5 m). La combustion est incomplète et produit \SI{5,0}{g} de $\ce{CO}$ par heure. Calculer la concentration en $\ce{CO}$ (en ppm = parties par million, volume) au bout de 2 h sans aération. (Conditions normales : $V_m = \SI{24,0}{L/mol}$ ; $1~\text{ppm} = 1~\text{cm}^3/\text{m}^3$).
    \item Le détecteur de CO placé dans le couloir sonne au bout de 30 min. Quelle action \textbf{la première} à faire ? (Choisir : a) Ouvrir la fenêtre de la chambre ; b) Sortir toute la famille dehors ; c) Éteindre le chauffage ; d) Appeler le 112).
    \item Pourquoi ne place-t-on pas le détecteur $\ce{CO}$ dans la cuisine ?
\end{enumerate}
\end{exercice}

\begin{corrige}[Exercice : Détecteur et concentration CO]
\begin{enumerate}
    \item $m(\ce{CO}) = 2 \times 5,0 = \SI{10,0}{g}$. $M(\ce{CO}) = 12+16 = \SI{28}{g/mol}$.
    $n(\ce{CO}) = 10,0 / 28 \approx \SI{0,357}{mol}$.
    $V(\ce{CO}) = n \times V_m = 0,357 \times 24,0 \approx \SI{8,57}{L} = \SI{8570}{cm^3}$.
    Volume chambre = \SI{20}{m^3} = \SI{20\,000\,000}{cm^3} (ou \SI{20000}{L}).
    Concentration (ppm) = $\frac{V(\ce{CO})~\text{en cm}^3}{V(\text{chambre})~\text{en m}^3} = \frac{8570}{20} \approx \boxed{429~\text{ppm}}$.
    \textbf{Danger :} Seuil alarme détecteur souvent 50-100 ppm. 429 ppm = \textbf{très dangereux} (maux de tête en 1-2h, mortel en 3-4h).
    \item \textbf{b) Sortir toute la famille dehors.} (Priorité absolue : évacuation vers air pur. On n'ouvre pas la fenêtre soi-même si on est déjà intoxiqué, on sort).
    \item Cuisine = vapeurs de cuisson, graisses, humidité $\to$ \textbf{fausses alertes} fréquentes / encrassement capteur. Norme : chambres + couloirs (proches zones sommeil).
\end{enumerate}
\end{corrige}

\newpage

\coursec{Exercices}

\coursec{La réaction chimique}

\courssub{Je vérifie mes connaissances}

\begin{exercice}[QCM : Transformation chimique ou physique ?]
Parmi les situations suivantes, lesquelles correspondent à une \cle{transformation chimique} ? Justifie chaque choix en une phrase.
\begin{enumerate}
    \item La \cle{fonte} de la cire d'une bougie allumée.
    \item La \cle{combustion} de la cire d'une bougie allumée.
    \item La \cle{dissolution} du sel de cuisine dans l'eau.
    \item La formation de \cle{rouille} sur un clou laissé dehors sous la pluie.
    \item L'\cle{évaporation} de l'essence renversée sur le sol.
\end{enumerate}
\end{exercice}

\begin{exercice}[Vrai ou Faux : Loi de Lavoisier et équation-bilan]
Pour chaque affirmation, réponds « Vrai » ou « Faux » et justifie brièvement.
\begin{enumerate}
    \item Lors d'une réaction chimique dans un récipient fermé, la masse totale du système augmente.
    \item Dans l'équation-bilan $\ce{2H2 + O2 -> 2H2O}$, les atomes d'hydrogène et d'oxygène sont conservés.
    \item Les \cle{réactifs} sont écrits à droite de la flèche et les \cle{produits} à gauche.
    \item La combustion complète du carbone s'écrit $\ce{C + O2 -> CO}$.
    \item Le monoxyde de carbone $\ce{CO}$ est un gaz incolore, inodore et très toxique.
\end{enumerate}
\end{exercice}

\begin{exercice}[Définitions et notations]
Donne la définition précise des termes suivants :
\begin{enumerate}
    \item \cle{Réaction chimique} (ou transformation chimique).
    \item \cle{Équation-bilan} d'une réaction chimique.
    \item \cle{Équilibrer} une équation-bilan.
    \item \cle{Combustion complète} et \cle{combustion incomplète} d'un combustible carboné.
    \item \cle{Monoxyde de carbone} : formule, origine, danger principal.
\end{enumerate}
\end{exercice}

\begin{exercice}[Calcul direct : Conservation de la masse]
On réalise la réaction suivante dans un ballon fermé posé sur une balance : $\ce{Mg + 2HCl -> MgCl2 + H2}$.
Au départ, la masse du ballon et de son contenu est de \SI{120,0}{g}. À la fin de la réaction, on observe une masse de \SI{120,0}{g}.
\begin{enumerate}
    \item Quelle loi fondamentale ce résultat illustre-t-il ? Énonce-la.
    \item Si la réaction avait été réalisée dans un tube à essai ouvert, que serait-il advenu de la masse mesurée sur la balance ? Explique.
\end{enumerate}
\end{exercice}

\courssub{Je m'entraîne}

\begin{exercice}[Équilibrage d'équations-bilans]
Équilibre les équations-bilans suivantes (coefficients entiers les plus petits) :
\begin{enumerate}
    \item $\ce{Fe + O2 -> Fe3O4}$ \quad (formation de la magnétite)
    \item $\ce{C3H8 + O2 -> CO2 + H2O}$ \quad (combustion complète du propane, gaz de bouteille)
    \item $\ce{Al + HCl -> AlCl3 + H2}$ \quad (réaction de l'aluminium avec l'acide chlorhydrique)
    \item $\ce{C2H5OH + O2 -> CO2 + H2O}$ \quad (combustion complète de l'éthanol, carburant « super »)
    \item $\ce{KClO3 -> KCl + O2}$ \quad (décomposition du chlorate de potassium, source d'$\ce{O2}$ au labo)
\end{enumerate}
\end{exercice}

\begin{exercice}[Combustion du carbone : application de la loi de Lavoisier]
On fait brûler \SI{12,0}{g} de carbone (charbon de bois) dans un excès de dioxygène, dans un récipient \textbf{fermé} de masse \SI{48,0}{g}. La masse totale initiale (récipient + carbone + dioxygène) est de \SI{100,0}{g}.
\begin{enumerate}
    \item Écrire l'équation-bilan de la combustion complète du carbone.
    \item Quelle est la masse totale du système \textbf{après} la réaction ? Justifie en citant la loi de Lavoisier.
    \item Quelle masse de dioxyde de carbone $\ce{CO2}$ a été formée ? (On admet que tout le carbone a réagi).
\end{enumerate}
\end{exercice}

\begin{exercice}[Interprétation d'expérience : Pyrolyse du sucre]
On chauffe fortement du sucre en poudre (saccharose $\ce{C12H22O11}$) dans un tube à essai à l'aide d'un bec Bunsen. On observe : le sucre fond, noircit, dégage une odeur de caramel brûlé, un gaz s'échappe et un liquide incolore se condense sur les parois froides du tube. Un résidu noir solide reste au fond.
\begin{enumerate}
    \item Nomme les deux principaux produits observés (liquide condensé et résidu solide).
    \item S'agit-il d'une combustion ? Justifie.
    \item Propose une équation-bilan \textbf{simplifiée} (non équilibrée) pour cette transformation, en notant le résidu noir $\ce{C}$.
\end{enumerate}
\end{exercice}

\begin{exercice}[Problème concret : Gazinière mal réglée]
Dans une cuisine à Douala, la flamme d'une gazinière alimentée au butane ($\ce{C4H10}$) est \textbf{jaune-orangée} et \textbf{noircit} le fond des marmites (dépôt de suie).
\begin{enumerate}
    \item Quel type de combustion se produit ? Explique l'origine de la couleur jaune et du dépôt noir.
    \item Écrire l'équation-bilan de cette combustion (non équilibrée).
    \item Quel \cle{gaz toxique}, incolore et inodore, risque de se former en quantité dangereuse dans la cuisine si la ventilation est mauvaise ?
    \item Quelle action simple sur le robinet d'arrivée d'air (ou le tiroir de la gazinière) permet de retrouver une flamme bleue et non salissante ?
\end{enumerate}
\end{exercice}

\courssub{Je me prépare au Bac}

\begin{exercice}[Sujet type BEPC 2023 : Intoxication au monoxyde de carbone]
\textbf{Document 1 : Extrait d'un article de presse (Cameroon Tribune, mars 2023).}
\emph{« Une famille de cinq personnes a été retrouvée inconsciente dans leur domicile du quartier Mvog-Ada à Yaoundé. La cause : une intoxication au monoxyde de carbone provenant d'un groupe électrogène placé dans le couloir, porte fermée, suite à une coupure d'électricité ENEO. Les secours ont mesuré un taux de CO de 800 ppm dans la pièce. »}

\textbf{Document 2 : Données chimiques.}
Le groupe électrogène fonctionne à l'essence (principalement de l'octane $\ce{C8H18}$). L'air contient environ 21\,\% de dioxygène $\ce{O2}$.
\begin{enumerate}
    \item À partir du Document 1, identifier la situation à risque (emplacement du groupe, confinement).
    \item Écrire l'équation-bilan de la \textbf{combustion incomplète} de l'octane produisant du monoxyde de carbone $\ce{CO}$ et de l'eau. Équilibre-la.
    \item Écrire l'équation-bilan de la \textbf{combustion complète} de l'octane. Équilibre-la.
    \item Expliquer pourquoi la combustion est incomplète dans le couloir (raison chimique).
    \item Citer \textbf{trois} mesures de prévention concrètes pour éviter ce drame (emplacement, ventilation, détecteur, entretien).
    \item Le taux de 800 ppm signifie 800 L de CO pour 1\,000\,000 L d'air. Sachant qu'un taux supérieur à 100 ppm est dangereux pour une exposition prolongée, commenter la gravité de la situation.
\end{enumerate}
\end{exercice}

\begin{exercice}[Sujet type BEPC 2024 : Dosage et équation-bilan — Réaction acide-base]
On souhaite mettre en évidence la réaction entre l'acide chlorhydrique (solution $\ce{HCl}$) et l'hydroxyde de sodium (solution $\ce{NaOH}$).
\textbf{Protocole :}
\begin{itemize}
    \item Bécher A : \SI{20,0}{mL} de solution d'acide chlorhydrique de concentration molaire $C_A = \SI{0,10}{mol/L}$.
    \item Bécher B : \SI{20,0}{mL} de solution d'hydroxyde de sodium de concentration molaire $C_B = \SI{0,10}{mol/L}$.
    \item On verse B dans A. La température monte légèrement. On obtient une solution incolore.
    \item On ajoute quelques gouttes de solution de nitrate d'argent $\ce{AgNO3}$ : \textbf{un précipité blanc se forme}.
    \item On évapore l'eau à sec : il reste un solide blanc cristallisé.
\end{itemize}
\begin{enumerate}
    \item Identifier les \cle{réactifs} et les \cle{produits} de la réaction. Écrire l'équation-bilan équilibrée (ions ou molécules).
    \item Calculer la quantité de matière (en mol) de $\ce{HCl}$ et de $\ce{NaOH}$ introduits. Qui est le \cle{réactif limitant} ? Justifie.
    \item Interpréter la formation du précipité avec $\ce{AgNO3}$ (rappeler le test des ions chlorure $\ce{Cl-}$).
    \item Nommer le solide blanc obtenu après évaporation. Donner son utilisation courante au Cameroun.
    \item Cette réaction est-elle une réaction d'oxydo-réduction ? Justifie par l'analyse des nombres d'oxydation.
\end{enumerate}
\end{exercice}

\begin{exercice}[Sujet type BEPC 2025 : Synthèse — Du pétrole au plastique et ses déchets]
\textbf{Document :} Le Cameroun importe une grande partie de ses produits pétroliers raffinés (essence, gasoil, butane) via la SONARA (Limbe). Les sachets plastiques (polyéthylène) envahissent les caniveaux de Yaoundé et Douala, bouchant les égouts et favorisant les inondations.
\begin{enumerate}
    \item \textbf{Origine :} Le polyéthylène $(-\ce{CH2-CH2}-)_n$ est un polymère obtenu à partir d'éthylène $\ce{C2H4}$. L'éthylène provient du \cle{craquage} du naphta (pétrole). Écrire l'équation-bilan simplifiée du craquage d'un alcane $\ce{C10H22}$ (décane) en éthylène et un autre alcane plus court.
    \item \textbf{Combustion :} Un sachet plastique brûlé à l'air libre (feu de brousse) produit une flamme jaune et de la fumée noire. Écrire l'équation-bilan de la combustion \textbf{incomplète} de l'éthylène $\ce{C2H4}$ produisant du carbone $\ce{C}$ (suie) et de l'eau. Équilibre-la.
    \item \textbf{Danger :} Expliquer pourquoi brûler les déchets plastiques à l'air libre est une source majeure de pollution de l'air (citer au moins deux polluants émis).
    \item \textbf{Gestion durable :} Proposer deux solutions techniques ou citoyennes pour réduire la pollution par les sachets plastiques au Cameroun (loi 2012, tri, recyclage, alternatives).
    \item \textbf{Calcul :} Un sachet a une masse de \SI{5,0}{g}. On le brûle complètement dans un excès d'oxygène (combustion complète en $\ce{CO2}$ et $\ce{H2O}$). Calculer la masse de dioxyde de carbone $\ce{CO2}$ rejetée dans l'atmosphère. (M(C)=12, M(H)=1, M(O)=16 g/mol).
\end{enumerate}
\end{exercice}



\newpage

\coursec{Corrigés des exercices}

\begin{corrige}[Exercice 1 — QCM : Transformation chimique ou physique ?]
\begin{enumerate}
    \item \textbf{Non, transformation physique.} La fonte est un simple changement d'état (solide $\to$ liquide) : la cire reste de la cire, aucune espèce chimique nouvelle n'apparaît. En refroidissant, on retrouve la cire solide.
    \item \textbf{Oui, transformation chimique.} La combustion de la cire consomme du dioxygène et produit des corps nouveaux : dioxyde de carbone $\ce{CO2}$, vapeur d'eau $\ce{H2O}$, et parfois de la suie $\ce{C}$. Il y a disparition de réactifs et apparition de produits.
    \item \textbf{Non, transformation physique.} Le sel se disperse en ions $\ce{Na+}$ et $\ce{Cl-}$ dans l'eau, mais il ne change pas de nature : par évaporation de l'eau, on récupère le sel identique.
    \item \textbf{Oui, transformation chimique.} La rouille (oxyde de fer hydraté) est un corps nouveau, formé par réaction du fer avec le dioxygène et l'eau. Le fer a disparu en tant que tel.
    \item \textbf{Non, transformation physique.} L'évaporation est un changement d'état (liquide $\to$ gaz) : les molécules d'essence restent les mêmes.
\end{enumerate}
\end{corrige}

\begin{attention}[Point de vigilance]
Ne pas confondre \emph{fusion/fonte} (physique) et \emph{combustion} (chimique) pour une même bougie : les deux se produisent simultanément, mais seule la combustion crée des corps nouveaux. Critère décisif : \textbf{apparition d'espèces chimiques nouvelles}.
\end{attention}

\begin{corrige}[Exercice 2 — Vrai ou Faux : Loi de Lavoisier et équation-bilan]
\begin{enumerate}
    \item \textbf{Faux.} D'après la \cle{loi de Lavoisier}, dans un récipient fermé, la masse totale du système \textbf{reste constante} au cours d'une réaction chimique : rien ne se perd, rien ne se crée.
    \item \textbf{Vrai.} Comptons les atomes. Réactifs : $2 \times 2 = 4$ atomes H et $2$ atomes O. Produits : $2 \times 2 = 4$ atomes H et $2 \times 1 = 2$ atomes O. Chaque élément est présent en même nombre des deux côtés : les atomes sont conservés, ils sont seulement réarrangés.
    \item \textbf{Faux.} C'est l'inverse : les \cle{réactifs} s'écrivent à \textbf{gauche} de la flèche et les \cle{produits} à \textbf{droite} : $\text{réactifs} \to \text{produits}$.
    \item \textbf{Faux.} $\ce{C + O2 -> CO}$ n'est d'ailleurs pas équilibrée en oxygène ; la formation de $\ce{CO}$ traduit une combustion \textbf{incomplète} : $\ce{2C + O2 -> 2CO}$. La combustion complète du carbone s'écrit $\ce{C + O2 -> CO2}$.
    \item \textbf{Vrai.} Le monoxyde de carbone $\ce{CO}$ est incolore, inodore et non irritant, donc indétectable par nos sens, et très toxique : il se fixe sur l'hémoglobine du sang à la place du dioxygène, provoquant l'asphyxie des cellules.
\end{enumerate}
\end{corrige}

\begin{attention}[Point de vigilance]
Au BEPC, préciser toujours si le système est \textbf{fermé} ou \textbf{ouvert} avant d'appliquer la loi de Lavoisier. Une variation de masse en système ouvert (gaz qui s'échappe ou $\ce{O2}$ de l'air qui entre) ne contredit jamais la loi.
\end{attention}

\begin{corrige}[Exercice 3 — Définitions et notations]
\begin{enumerate}
    \item \textbf{Réaction chimique :} transformation au cours de laquelle une ou plusieurs espèces chimiques (les \cle{réactifs}) disparaissent pour former une ou plusieurs espèces chimiques nouvelles (les \cle{produits}).
    \item \textbf{Équation-bilan :} écriture symbolique d'une réaction chimique : les formules des réactifs à gauche, celles des produits à droite, séparées par une flèche, avec éventuellement l'état physique $(s)$, $(l)$, $(g)$, $(aq)$ de chaque espèce.
    \item \textbf{Équilibrer :} placer devant chaque formule le plus petit nombre entier (le \cle{coefficient stœchiométrique}) tel que le nombre d'atomes de chaque élément soit le même dans les réactifs et dans les produits, traduisant ainsi la conservation des atomes.
    \item \textbf{Combustion complète :} combustion réalisée avec un \emph{excès} de dioxygène ; un combustible carboné ne produit alors que du dioxyde de carbone $\ce{CO2}$ et de l'eau $\ce{H2O}$. \textbf{Combustion incomplète :} combustion réalisée avec un \emph{défaut} de dioxygène ; il se forme en plus du monoxyde de carbone $\ce{CO}$ (toxique) et/ou du carbone $\ce{C}$ (suie).
    \item \textbf{Monoxyde de carbone :} formule $\ce{CO}$. Il provient des combustions incomplètes de combustibles carbonés (bois, charbon, essence, gasoil, butane) lorsque l'apport d'air est insuffisant. Danger principal : intoxication pouvant être mortelle, car il bloque le transport de l'oxygène par le sang.
\end{enumerate}
\end{corrige}

\begin{corrige}[Exercice 4 — Calcul direct : Conservation de la masse]
\begin{enumerate}
    \item Ce résultat illustre la \cle{loi de Lavoisier} (loi de conservation de la masse) : \emph{« Dans une réaction chimique réalisée en milieu fermé, la masse totale ne varie pas : la masse des réactifs consommés est égale à la masse des produits formés. »} Ici, la balance indique \SI{120,0}{g} avant et après : la masse est conservée.
    \item Dans un tube à essai \textbf{ouvert}, la masse mesurée aurait \textbf{diminué}. En effet, le dihydrogène $\ce{H2}$ produit est un gaz qui s'échapperait dans l'air : sa masse sortirait du système pesé. La perte de masse observée serait exactement égale à la masse de $\ce{H2}$ dégagé. La loi de Lavoisier n'est pas contredite : elle ne s'applique qu'aux systèmes fermés.
\end{enumerate}
\end{corrige}

\begin{attention}[Point de vigilance]
Au BEPC, préciser toujours si le système est \textbf{fermé} ou \textbf{ouvert} avant d'appliquer la loi de Lavoisier. Une variation de masse en système ouvert (gaz qui s'échappe ou $\ce{O2}$ de l'air qui entre) ne contredit jamais la loi.
\end{attention}

\begin{corrige}[Exercice 5 — Équilibrage d'équations-bilans]
Méthode : on équilibre élément par élément, en terminant par l'oxygène et l'hydrogène.
\begin{enumerate}
    \item $\ce{Fe + O2 -> Fe3O4}$ : il faut 3 Fe à gauche et 4 O à gauche, soit $2\,\ce{O2}$ :
    \[ \ce{3Fe + 2O2 -> Fe3O4} \]
    Vérification : Fe : $3 = 3$ ; O : $4 = 4$.
    \item $\ce{C3H8 + O2 -> CO2 + H2O}$ : 3 C donnent $3\,\ce{CO2}$ ; 8 H donnent $4\,\ce{H2O}$ ; O à droite : $6 + 4 = 10$, soit $5\,\ce{O2}$ :
    \[ \ce{C3H8 + 5O2 -> 3CO2 + 4H2O} \]
    Vérification : C : $3 = 3$ ; H : $8 = 8$ ; O : $10 = 10$.
    \item $\ce{Al + HCl -> AlCl3 + H2}$ : 3 Cl dans $\ce{AlCl3}$ demandent $3\,\ce{HCl}$, puis on double tout pour obtenir des coefficients entiers en H :
    \[ \ce{2Al + 6HCl -> 2AlCl3 + 3H2} \]
    Vérification : Al : $2 = 2$ ; H : $6 = 6$ ; Cl : $6 = 6$.
    \item $\ce{C2H5OH + O2 -> CO2 + H2O}$ : 2 C donnent $2\,\ce{CO2}$ ; 6 H donnent $3\,\ce{H2O}$ ; O à droite : $4 + 3 = 7$, dont 1 déjà dans l'éthanol, il en reste 6, soit $3\,\ce{O2}$ :
    \[ \ce{C2H5OH + 3O2 -> 2CO2 + 3H2O} \]
    Vérification : C : $2 = 2$ ; H : $6 = 6$ ; O : $1 + 6 = 7 = 4 + 3$.
    \item $\ce{KClO3 -> KCl + O2}$ : 3 O à gauche ; le p.p.c.m. avec 2 est 6, d'où $2\,\ce{KClO3}$ et $3\,\ce{O2}$ :
    \[ \ce{2KClO3 -> 2KCl + 3O2} \]
    Vérification : K : $2 = 2$ ; Cl : $2 = 2$ ; O : $6 = 6$.
\end{enumerate}
\end{corrige}

\begin{attention}[Point de vigilance]
On ne modifie \textbf{jamais} les indices dans les formules ($\ce{H2O}$ ne devient pas $\ce{H2O2}$) : on place uniquement des coefficients \emph{devant} les formules. Les coefficients doivent être les plus petits entiers possibles.
\end{attention}

\begin{corrige}[Exercice 6 — Combustion du carbone : application de la loi de Lavoisier]
\begin{enumerate}
    \item Combustion complète du carbone :
    \[ \ce{C(s) + O2(g) -> CO2(g)} \]
    \item La masse totale après réaction est \SI{100,0}{g}. Justification : le récipient est \textbf{fermé}, donc aucune matière n'entre ni ne sort ; d'après la \cle{loi de Lavoisier}, la masse totale du système est conservée : $m_{\text{après}} = m_{\text{avant}} = \SI{100,0}{g}$.
    \item Calcul de la masse de $\ce{CO2}$ formée :
    \begin{align*}
    n(\ce{C}) &= \frac{m(\ce{C})}{M(\ce{C})} = \frac{12,0}{12,0} = \SI{1,00}{mol} \\
    \text{D'après l'équation : } & 1\ \text{mol de } \ce{C} \to 1\ \text{mol de } \ce{CO2} \\
    n(\ce{CO2}) &= \SI{1,00}{mol} \\
    m(\ce{CO2}) &= n \times M(\ce{CO2}) = 1,00 \times (12,0 + 2 \times 16,0) = 1,00 \times 44,0 = \SI{44,0}{g}
    \end{align*}
    \textbf{Vérification par Lavoisier :} masse de dioxygène consommée $= 1,00 \times 32,0 = \SI{32,0}{g}$. Or $m(\ce{C}) + m(\ce{O2}) = 12,0 + 32,0 = \SI{44,0}{g} = m(\ce{CO2})$. La masse des réactifs consommés égale la masse du produit formé : la loi est vérifiée.
\end{enumerate}
\end{corrige}

\begin{corrige}[Exercice 7 — Interprétation d'expérience : Pyrolyse du sucre]
\begin{enumerate}
    \item Le liquide incolore qui se condense sur les parois froides est de l'\textbf{eau} $\ce{H2O}$ (on peut l'identifier au sulfate de cuivre anhydre qui bleuit). Le résidu noir solide est du \textbf{carbone} $\ce{C}$ (charbon).
    \item \textbf{Non, ce n'est pas une combustion.} Une combustion est une réaction avec le dioxygène de l'air, vive, avec flamme. Ici, le sucre est chauffé dans un tube sans apport de dioxygène : c'est une \cle{pyrolyse}, c'est-à-dire une décomposition d'un corps par la chaleur. Le sucre se décompose en carbone et en eau.
    \item Équation-bilan simplifiée :
    \[ \ce{C12H22O11 ->[chaleur] 12C + 11H2O} \]
    Vérification : C : $12 = 12$ ; H : $22 = 22$ ; O : $11 = 11$. L'équation est équilibrée : les 12 atomes de carbone du saccharose se retrouvent dans le résidu noir, et l'hydrogène et l'oxygène forment l'eau.
\end{enumerate}
\end{corrige}

\begin{attention}[Point de vigilance]
Bien distinguer \textbf{pyrolyse} (décomposition par la chaleur, sans dioxygène) et \textbf{combustion} (réaction avec le dioxygène). Le mot « on chauffe » ne signifie pas automatiquement « combustion ».
\end{attention}

\begin{corrige}[Exercice 8 — Problème concret : Gazinière mal réglée]
\begin{enumerate}
    \item Il s'agit d'une \textbf{combustion incomplète} : l'apport en dioxygène est insuffisant. La couleur jaune-orangée de la flamme est due à des particules de \textbf{carbone} portées à incandescence dans la flamme. Le dépôt noir au fond des marmites est de la \textbf{suie} $\ce{C}$, du carbone qui n'a pas brûlé.
    \item Équation-bilan (non équilibrée, produits multiples possibles) :
    \[ \ce{C4H10 + O2 -> CO2 + CO + C + H2O} \]
    \item Le gaz toxique qui risque de se former est le \textbf{monoxyde de carbone} $\ce{CO}$, incolore, inodore et mortel. Dans une cuisine mal ventilée, il s'accumule et peut intoxiquer toute la famille.
    \item Il faut \textbf{ouvrir davantage l'arrivée d'air} (régler la virole ou le tiroir d'aération du brûleur) pour augmenter l'apport en dioxygène. La combustion redevient complète : la flamme est \textbf{bleue}, plus chaude et non salissante. Il faut aussi déboucher les injecteurs encrassés par la suie.
\end{enumerate}
\end{corrige}

\begin{attention}[Point de vigilance]
Retenir le code couleur des flammes : flamme \textbf{bleue} = combustion complète (assez d'air) ; flamme \textbf{jaune} = combustion incomplète (manque d'air, production de suie et de $\ce{CO}$). C'est une question classique de sécurité au BEPC.
\end{attention}

\begin{corrige}[Exercice 9 — Sujet type BEPC 2023 : Intoxication au monoxyde de carbone]
\begin{enumerate}
    \item La situation à risque : le groupe électrogène est placé \textbf{à l'intérieur de l'habitation} (dans un couloir), dans un espace \textbf{confiné} (porte fermée). Les gaz d'échappement ne sont pas évacués vers l'extérieur et l'air ne se renouvelle pas.
    \item Combustion incomplète de l'octane produisant du $\ce{CO}$ :
    \[ \ce{2C8H18 + 17O2 -> 16CO + 18H2O} \]
    Vérification : C : $16 = 16$ ; H : $36 = 36$ ; O : $34 = 16 + 18 = 34$.
    \item Combustion complète de l'octane :
    \[ \ce{2C8H18 + 25O2 -> 16CO2 + 18H2O} \]
    Vérification : C : $16 = 16$ ; H : $36 = 36$ ; O : $50 = 32 + 18 = 50$.
    \item Dans le couloir fermé, le moteur \textbf{consomme le dioxygène} de l'air sans que celui-ci soit renouvelé. La proportion d'$\ce{O2}$ chute en dessous de ce qui est nécessaire à l'oxydation complète du carbone en $\ce{CO2}$ : la combustion devient incomplète et produit massivement du monoxyde de carbone $\ce{CO}$, qui s'accumule dans la pièce.
    \item Trois mesures de prévention :
    \begin{itemize}
        \item Installer le groupe électrogène \textbf{à l'extérieur} de la maison, loin des portes et fenêtres, jamais dans un couloir, un garage ou une pièce de vie ;
        \item Assurer une \textbf{ventilation permanente} des pièces (grilles d'aération non obstruées) et ne jamais boucher les aérations ;
        \item Installer un \textbf{détecteur de monoxyde de carbone} dans l'habitation et faire \textbf{entretenir régulièrement} le moteur (filtre à air, carburateur) pour garantir une combustion la plus complète possible.
    \end{itemize}
    \item Le taux mesuré (800 ppm) est \textbf{8 fois supérieur} au seuil de dangerosité en exposition prolongée (100 ppm). À une telle concentration, l'intoxication est rapide : maux de tête, nausées, vertiges, puis perte de conscience et mort par asphyxie en quelques heures, d'autant plus que les victimes dorment et ne perçoivent rien (gaz inodore). La situation décrite est donc \textbf{extrêmement grave} ; seuls l'évacuation immédiate et l'aération des lieux pouvaient sauver la famille.
\end{enumerate}
\end{corrige}

\begin{attention}[Barème indicatif (sur 10) et points de vigilance]
Barème : Q1 : 1 pt ; Q2 : 2 pts ; Q3 : 2 pts ; Q4 : 1,5 pt ; Q5 : 1,5 pt ; Q6 : 2 pts.\\
Points de vigilance : citer explicitement le \textbf{manque de dioxygène} comme cause chimique de la combustion incomplète ; exiger des équations \textbf{équilibrées} (coefficients entiers) ; pour la prévention, donner des mesures \textbf{concrètes} et non des généralités ; rappeler que le $\ce{CO}$ est indétectable par les sens humains.
\end{attention}

\begin{corrige}[Exercice 10 — Sujet type BEPC 2024 : Dosage et équation-bilan — Réaction acide-base]
\begin{enumerate}
    \item Réactifs : l'acide chlorhydrique $\ce{HCl}$ et l'hydroxyde de sodium $\ce{NaOH}$. Produits : le chlorure de sodium $\ce{NaCl}$ et l'eau $\ce{H2O}$.
    Équation-bilan moléculaire :
    \[ \ce{HCl(aq) + NaOH(aq) -> NaCl(aq) + H2O(l)} \]
    Équation ionique (les ions $\ce{Na+}$ et $\ce{Cl-}$ sont spectateurs) :
    \[ \ce{H+(aq) + OH-(aq) -> H2O(l)} \]
    L'élévation de température confirme qu'une réaction chimique (exothermique) a eu lieu.
    \item Quantités de matière introduites :
    \begin{align*}
    n(\ce{HCl}) &= C_A \times V_A = 0,10 \times 0,0200 = \SI{2,0 \times 10^{-3}}{mol} \\
    n(\ce{NaOH}) &= C_B \times V_B = 0,10 \times 0,0200 = \SI{2,0 \times 10^{-3}}{mol}
    \end{align*}
    L'équation montre que la réaction se fait mole à mole (rapport 1:1). Comme $n(\ce{HCl}) = n(\ce{NaOH}) = \SI{2,0 \times 10^{-3}}{mol}$, le mélange est réalisé dans les \textbf{proportions stœchiométriques} : il n'y a \textbf{pas de réactif limitant}, les deux réactifs sont entièrement consommés.
    \item Le précipité blanc est du \textbf{chlorure d'argent} $\ce{AgCl}$. C'est le test caractéristique des ions chlorure $\ce{Cl-}$ :
    \[ \ce{Ag+(aq) + Cl-(aq) -> AgCl(s)} \]
    Ce précipité blanc, qui noircit à la lumière, prouve la présence des ions $\ce{Cl-}$ en solution, apportés par le $\ce{NaCl}$ formé (et par l'acide initial).
    \item Le solide blanc obtenu après évaporation est le \textbf{chlorure de sodium} $\ce{NaCl}$, c'est-à-dire le sel de cuisine. Au Cameroun, il sert à assaisonner les plats, à conserver les aliments (poisson fumé, viandes) et à préparer les solutés de réhydratation orale.
    \item \textbf{Non, ce n'est pas une réaction d'oxydo-réduction.} Analysons les nombres d'oxydation : H est à $+1$, Cl à $-1$, Na à $+1$ et O à $-2$, aussi bien dans les réactifs que dans les produits. \textbf{Aucun nombre d'oxydation ne varie} : il n'y a pas de transfert d'électrons. C'est une réaction \textbf{acide-base}, caractérisée par un transfert de proton $\ce{H+}$ de l'acide vers la base.
\end{enumerate}
\end{corrige}

\begin{attention}[Barème indicatif (sur 10) et points de vigilance]
Barème : Q1 : 2 pts ; Q2 : 3 pts ; Q3 : 2 pts ; Q4 : 1 pt ; Q5 : 2 pts.\\
Points de vigilance : convertir les volumes en litres ($\SI{20,0}{mL} = \SI{0,0200}{L}$) avant le calcul ; justifier l'absence de réactif limitant par la \textbf{comparaison des quantités au rapport stœchiométrique}, pas seulement par l'égalité des nombres ; pour Q5, la justification par les nombres d'oxydation (ou par le transfert de proton) est exigée, une simple réponse « non » ne rapporte aucun point.
\end{attention}

\begin{corrige}[Exercice 11 — Sujet type BEPC 2025 : Synthèse — Du pétrole au plastique et ses déchets]
\begin{enumerate}
    \item Craquage du décane en éthylène et en un alcane plus court (octane) :
    \[ \ce{C10H22 -> C2H4 + C8H18} \]
    Vérification : C : $10 = 2 + 8$ ; H : $22 = 4 + 18$. L'équation est équilibrée.
    \item Combustion incomplète de l'éthylène produisant du carbone et de l'eau :
    \[ \ce{C2H4 + O2 -> 2C + 2H2O} \]
    Vérification : C : $2 = 2$ ; H : $4 = 4$ ; O : $2 = 2$.
    \item Brûler les plastiques à l'air libre se fait à basse température et avec peu d'oxygène : c'est une combustion incomplète qui émet notamment :
    \begin{itemize}
        \item du \textbf{monoxyde de carbone} $\ce{CO}$, gaz toxique qui empêche le transport de l'oxygène par le sang ;
        \item des \textbf{particules fines de carbone} (suies, fumées noires) qui pénètrent profondément dans les poumons et provoquent des maladies respiratoires ;
        \item des \textbf{composés organiques volatils} et des hydrocarbures aromatiques, dont certains sont cancérigènes.
    \end{itemize}
    \item Deux solutions pour réduire la pollution par les sachets plastiques :
    \begin{itemize}
        \item \textbf{Appliquer la réglementation camerounaise} interdisant les sachets plastiques non biodégradables de faible épaisseur, et utiliser des emballages alternatifs : sachets biodégradables, sacs en tissu ou en papier réutilisables ;
        \item \textbf{Organiser le tri, la collecte et le recyclage} des plastiques (transformation en granulés, pavés, objets manufacturés) et sensibiliser les populations à ne pas jeter les sachets dans les caniveaux, ce qui réduit aussi les inondations.
    \end{itemize}
    \item Calcul de la masse de $\ce{CO2}$ rejetée. Le polyéthylène est assimilé à son motif $\ce{C2H4}$ de masse molaire :
    \[ M(\ce{C2H4}) = 2 \times 12 + 4 \times 1 = \SI{28}{g/mol} \]
    Combustion complète :
    \[ \ce{C2H4 + 3O2 -> 2CO2 + 2H2O} \]
    \begin{align*}
    n(\ce{C2H4}) &= \frac{m}{M} = \frac{5,0}{28} \approx \SI{0,179}{mol} \\
    n(\ce{CO2}) &= 2 \times n(\ce{C2H4}) = 2 \times 0,179 \approx \SI{0,358}{mol} \\
    m(\ce{CO2}) &= n \times M(\ce{CO2}) = 0,358 \times 44 \approx \SI{15,7}{g}
    \end{align*}
    \textbf{Conclusion :} la combustion complète d'un seul sachet de \SI{5,0}{g} rejette environ \SI{15,7}{g} de dioxyde de carbone, soit plus de trois fois sa propre masse, car le dioxygène de l'air s'ajoute au carbone du plastique. À l'échelle des millions de sachets brûlés chaque jour, la contribution au réchauffement climatique est considérable.
\end{enumerate}
\end{corrige}

\begin{attention}[Barème indicatif (sur 10) et points de vigilance]
Barème : Q1 : 1,5 pt ; Q2 : 1,5 pt ; Q3 : 2 pts ; Q4 : 2 pts ; Q5 : 3 pts.\\
Points de vigilance : pour Q1, vérifier la conservation des atomes de C et de H ; pour Q2, bien écrire $\ce{C}$ (suie) et non $\ce{CO}$ puisque l'énoncé impose le carbone ; pour Q5, ne pas oublier le facteur 2 du coefficient stœchiométrique devant $\ce{CO2}$ et conclure par une phrase avec l'unité ; le résultat $m(\ce{CO2}) > m(\text{sachet})$ doit être commenté (apport d'oxygène de l'air), c'est un piège classique.
\end{attention}

\newpage

\coursec{Fiche bilan}

\begin{popfigure}
\centering
\begin{tikzpicture}[
    mindmap,
    concept color=popblue,
    level 1 concept/.append style={level distance=4.5cm, sibling angle=60},
    level 2 concept/.append style={level distance=3cm},
    every node/.style={font=\small, text=white},
    grow cyclic,
    branch1/.style={concept color=popblue},
    branch2/.style={concept color=poporange},
    branch3/.style={concept color=popgreen},
    branch4/.style={concept color=poppurple},
    branch5/.style={concept color=poppink},
    branch6/.style={concept color=popgold},
]
\node[concept, text width=2.8cm, align=center, font=\bfseries\large] {Réaction\\ chimique}
    child[branch1] { node[concept, text width=2.5cm, align=center] {Transformation\\ chimique\\ (réactifs $\to$ produits)} }
    child[branch2] { node[concept, text width=2.5cm, align=center] {Loi de\\ Lavoisier\\ (conservation\\ de la masse)} }
    child[branch3] { node[concept, text width=2.5cm, align=center] {Équation-bilan\\ et équilibrage\\ (coefficients\\ stœchiométriques)} }
    child[branch4] { node[concept, text width=2.5cm, align=center] {Combustion\\ complète\\ $\ce{CO2 + H2O}$} }
    child[branch5] { node[concept, text width=2.5cm, align=center] {Combustion\\ incomplète\\ $\ce{CO}$ et $\ce{C}$\\ (danger)} }
    child[branch6] { node[concept, text width=2.5cm, align=center] {Sécurité\\ prévention CO\\ (aération,\\ détecteur)} };
\end{tikzpicture}
\legende{Carte mentale de la leçon : la réaction chimique, de l'observation expérimentale à la prévention des risques liés aux combustions.}
\end{popfigure}

\begin{bilan}

\courssub{Définitions clés}
\begin{itemize}
    \item \cle{Transformation chimique} : transformation au cours de laquelle des espèces chimiques (les réactifs) disparaissent et de nouvelles espèces chimiques (les produits) apparaissent. À l'échelle microscopique, les atomes ne disparaissent pas : ils se \textbf{réarrangent} pour former de nouvelles molécules.
    \item \cle{Réactifs} : espèces chimiques présentes avant la réaction (écrites à gauche de la flèche).
    \item \cle{Produits} : espèces chimiques formées par la réaction (écrites à droite de la flèche).
    \item \cle{Équation-bilan} : représentation symbolique de la réaction : $\ce{Reactifs -> Produits}$, avec les états physiques $(s)$ solide, $(l)$ liquide, $(g)$ gazeux, $(aq)$ en solution aqueuse, et les coefficients stœchiométriques.
    \item \cle{Combustion} : réaction chimique rapide entre un \cle{combustible} (qui brûle) et un \cle{comburant} (le dioxygène $\ce{O2}$ de l'air), qui dégage de la chaleur et de la lumière.
    \item \cle{Combustion complète} : le dioxygène est en quantité suffisante ; si le combustible contient du carbone et de l'hydrogène, les produits sont uniquement $\ce{CO2}$ et $\ce{H2O}$. Exemple : $\ce{CH4 + 2O2 -> CO2 + 2H2O}$.
    \item \cle{Combustion incomplète} : le dioxygène est en quantité insuffisante ; il se forme du \cle{monoxyde de carbone $\ce{CO}$} (gaz toxique, incolore et inodore) et/ou du carbone $\ce{C}$ (suie noire). Exemple : $\ce{2CH4 + 3O2 -> 2CO + 4H2O}$.
\end{itemize}

\courssub{Lois et formules à connaître par cœur}
\begin{center}
\begin{tabularx}{\linewidth}{|>{\bfseries}l|X|}\hline
\rowcolor{popblueL}\textbf{Notion} & \textbf{Énoncé / Formule} \\\hline
Loi de Lavoisier & « Rien ne se perd, rien ne se crée, tout se transforme. » Dans une transformation chimique, la masse totale des réactifs est égale à la masse totale des produits (système fermé). \\\hline
Équilibrage & Il doit y avoir le même nombre d'atomes de chaque élément des deux côtés de la flèche. On ne modifie \textbf{que les coefficients} placés devant les formules, jamais les formules elles-mêmes. \\\hline
Combustion complète du méthane & $\ce{CH4 + 2O2 -> CO2 + 2H2O}$ \\\hline
Combustion complète du butane (gaz des bonbonnes) & $\ce{2C4H10 + 13O2 -> 8CO2 + 10H2O}$ \\\hline
Combustion complète de l'octane (essence) & $\ce{2C8H18 + 25O2 -> 16CO2 + 18H2O}$ \\\hline
Test du $\ce{CO2}$ & L'eau de chaux $\ce{Ca(OH)2}$ se trouble en présence de $\ce{CO2}$ (formation d'un précipité blanc de calcaire $\ce{CaCO3}$). \\\hline
Test de l'eau $\ce{H2O}$ & Le sulfate de cuivre anhydre (blanc) bleuit en présence d'eau. \\\hline
\end{tabularx}
\end{center}

\begin{attention}[Pièges fréquents au BEPC]
\begin{itemize}
    \item Ne jamais modifier les indices dans une formule pour équilibrer : écrire $\ce{H2O2}$ à la place de $\ce{H2O}$ change l'espèce chimique ! On place un coefficient \textbf{devant} : $\ce{2H2O}$.
    \item La loi de Lavoisier n'est valable que si l'on compte \textbf{tous} les réactifs et \textbf{tous} les produits, y compris les gaz. Une bougie qui brûle « perd de la masse » seulement en apparence : les gaz produits s'échappent.
    \item Le monoxyde de carbone $\ce{CO}$ ne sent rien et ne se voit pas : on ne peut pas le détecter soi-même. Seul un détecteur de $\ce{CO}$ ou une bonne aération protège.
    \item Ne pas confondre $\ce{CO}$ (monoxyde de carbone, toxique) et $\ce{CO2}$ (dioxyde de carbone, non toxique mais étouffant en grande quantité).
\end{itemize}
\end{attention}

\courssub{Méthodes express}
\begin{methode}[Équilibrer une équation-bilan]
    \begin{enumerate}
        \item Écrire les formules brutes correctes des réactifs et des produits (ne plus jamais les modifier ensuite).
        \item Compter les atomes de chaque élément de chaque côté de la flèche.
        \item Commencer par l'élément présent dans le moins de formules (pour une combustion : C d'abord, puis H, puis O en dernier).
        \item Placer des coefficients devant les formules pour égaliser les nombres d'atomes.
        \item Vérifier élément par élément que les deux côtés sont égaux.
        \item Si tous les coefficients sont pairs (ou multiples d'un même nombre), les simplifier par leur plus grand diviseur commun.
    \end{enumerate}
\end{methode}

\begin{methode}[Identifier le type de combustion]
    \begin{enumerate}
        \item Observer la flamme : bleue et propre $\rightarrow$ combustion complète ; jaune-orange avec dépôt noir de suie $\rightarrow$ combustion incomplète.
        \item Tester les produits : eau de chaux qui se trouble ($\ce{CO2}$) et sulfate de cuivre anhydre qui bleuit ($\ce{H2O}$) confirment une combustion complète.
        \item Si un détecteur de $\ce{CO}$ sonne, ou si des personnes ressentent maux de tête et fatigue soudaine près d'un appareil de chauffage ou d'un groupe électrogène $\rightarrow$ combustion incomplète : DANGER, évacuer et aérer immédiatement.
    \end{enumerate}
\end{methode}

\end{bilan}

\begin{exoresolu}[Application de la loi de Lavoisier]
On fait brûler complètement $\SI{4,4}{g}$ de butane $\ce{C4H10}$ dans le dioxygène. Il se forme $\SI{13,2}{g}$ de dioxyde de carbone et $\SI{7,2}{g}$ d'eau. Calculer la masse de dioxygène consommée.
\tcblower
D'après la loi de Lavoisier, la masse totale des réactifs est égale à la masse totale des produits :
\begin{align*}
m(\ce{C4H10}) + m(\ce{O2}) &= m(\ce{CO2}) + m(\ce{H2O})\\
\SI{4,4}{g} + m(\ce{O2}) &= \SI{13,2}{g} + \SI{7,2}{g}\\
m(\ce{O2}) &= \SI{20,4}{g} - \SI{4,4}{g} = \SI{16}{g}
\end{align*}
La masse de dioxygène consommée est $\SI{16}{g}$.
\end{exoresolu}

\begin{aretenir}[Checklist avant le BEPC]
\begin{itemize}
    \item $\square$ Savoir définir transformation chimique, réactifs, produits, combustion, combustible et comburant.
    \item $\square$ Citer la loi de Lavoisier et l'appliquer à un calcul de masse (masse des réactifs = masse des produits).
    \item $\square$ Écrire une équation-bilan avec les états physiques et la flèche $\rightarrow$.
    \item $\square$ Équilibrer une équation simple (combustion, synthèse) par la méthode pas à pas.
    \item $\square$ Écrire l'équation de la combustion complète du méthane, du butane et de l'octane (essence).
    \item $\square$ Expliquer pourquoi la combustion incomplète est dangereuse (formation de $\ce{CO}$).
    \item $\square$ Décrire les symptômes d'une intoxication au $\ce{CO}$ : maux de tête, nausées, vertiges, fatigue, coma, mort.
    \item $\square$ Citer 3 gestes de prévention : aérer les pièces, entretenir les appareils (cuisinière, groupe électrogène, chauffe-eau), installer un détecteur de $\ce{CO}$ ; ne jamais faire fonctionner un groupe électrogène dans une pièce fermée.
    \item $\square$ Savoir identifier le $\ce{CO2}$ (eau de chaux qui se trouble) et l'eau (sulfate de cuivre anhydre qui bleuit).
    \item $\square$ Différencier combustible (ce qui brûle : bois, gaz, essence) et comburant (le dioxygène de l'air).
    \item $\square$ Justifier la couleur de la flamme : bleue = combustion complète ; jaune = combustion incomplète (particules de carbone incandescentes).
\end{itemize}
\end{aretenir}

\findecours

\end{document}
